\newcommand{\DoNotLoadEpstopdf}{}
\pdfinfoomitdate=1
\pdftrailerid{}
\pdfsuppressptexinfo=-1
\documentclass[11pt,letterpaper]{article}
\usepackage[T1]{fontenc}
\usepackage{lmodern}
\usepackage{amsmath,amssymb,amsthm,mathtools}
\usepackage[margin=1in]{geometry}
\usepackage{booktabs,longtable,array,pict2e}
\usepackage{microtype}
\usepackage[colorlinks=true,linkcolor=black,citecolor=black,urlcolor=blue]{hyperref}
\usepackage{xurl}
\setlength{\parskip}{4pt}
\setlength{\parindent}{1.2em}
\setlength{\emergencystretch}{2em}
\setcounter{tocdepth}{1}
\numberwithin{equation}{section}
\newtheorem{theorem}{Theorem}[section]
\newtheorem{lemma}[theorem]{Lemma}
\newtheorem{proposition}[theorem]{Proposition}
\newtheorem{corollary}[theorem]{Corollary}
\theoremstyle{remark}\newtheorem{remark}[theorem]{Remark}
% Macros of the five sources. \PP is P in Parts III-V; Parts I-II defined it as P_n
% and their uses were rewritten as \PP_n. \Li is the operator Li; Part V defined it
% as Li_2 and its uses were rewritten as \Li_2. Every other shared name has the same
% definition in each source.
\newcommand{\PP}{\mathbb P}
\newcommand{\EE}{\mathbb E}
\newcommand{\QQ}{\mathbb Q}
\newcommand{\ZZ}{\mathbb Z}
\newcommand{\NN}{\mathbb N}
\newcommand{\floor}[1]{\left\lfloor #1\right\rfloor}
\newcommand{\ceil}[1]{\left\lceil #1\right\rceil}
\newcommand{\dd}{\,\mathrm d}
\newcommand{\one}{\mathbf 1}
\newcommand{\area}{\mathcal A}
\newcommand{\Ai}{\operatorname{Ai}}
\DeclareMathOperator{\Li}{Li}
\DeclareMathOperator{\sech}{sech}
% Added in the merge (batch 107).
\newcommand{\Prb}{\mathbb P}
\newcommand{\file}[1]{\nolinkurl{#1}}
\newenvironment{writenote}[1][]{\par\smallskip\begingroup\small
 \noindent\textbf{[write]}\if\relax\detokenize{#1}\relax\else~\textbf{#1.}\fi\ }%
 {\par\endgroup\smallskip}
\newcommand{\wtag}{\textup{\textbf{[write]}}}
\newenvironment{partabstract}{\begin{quote}\small\noindent\textbf{Abstract of this Part (as delivered).}\ }{\end{quote}}
\newcommand{\partsource}[1]{\par\noindent\begin{minipage}{\textwidth}\small #1\end{minipage}\par\medskip}
\hypersetup{pdftitle={Subdiagonal and superdiagonal partitions (OEIS A238875 and A238873)},pdfauthor={},pdfsubject={Merged report a238873-diagonal-partitions (batch 107): proofs of the Archibald-Blecher-Elizalde-Knopfmacher Conjectures 5.2 and 5.3 (Report 156); the subdiagonal expansion to all orders (Report 159); the superdiagonal rate (Report 158), Airy term and limit shape (Report 160), and full equivalent conditional on arXiv:2607.27504v1 (Report 162)}}
\title{\textbf{Subdiagonal and superdiagonal partitions\\ (OEIS A238875 and A238873)}\\[7pt]
\large Two conjectures of Archibald, Blecher, Elizalde and Knopfmacher proved;\\ the subdiagonal ratio to all orders; the superdiagonal rate, Airy term and limit shape;\\ the full superdiagonal equivalent, conditional on a preprint}
\author{Five research reports of one external research session\\(bundle Reports 156, 159, 158, 160 and 162), merged into one ProveIt report}
\date{All Parts: 3 October 2026\quad merged: 6 October 2026}
\begin{document}
\hypersetup{pageanchor=false}
\maketitle
\begin{abstract}
Write partitions in weakly increasing order $\lambda_1\le\cdots\le\lambda_k$. Let $s_r(n)$ count those of $n$ with $\lambda_i\le i+r$ for every $i$, so that $s_0(n)$ is OEIS A238875 (subdiagonal partitions), and let $A(n)$ count those with $\lambda_i\ge i$, OEIS A238873 (superdiagonal partitions). Archibald, Blecher, Elizalde and Knopfmacher (Afrika Matematika 36:77, 2025) conjectured $q(n)/A(n)\to0$, with $q(n)$ the number of partitions into distinct parts (their Conjecture~5.2), and $s_0(n)/p(n)\to\frac12$ (Conjecture~5.3). This report prints five manuscripts as Parts~I--V. Part~I proves both conjectures: $0\le\mathbb P_n(R\le r)-s_r(n)/p(n)=O(1/\log n)$ uniformly in $r\ge0$, with $R$ Dyson's rank, a logistic crossover in $r/\sqrt n$, a Lambert-$W$ inverse, and $\liminf\log A(n)/\sqrt n\ge B_*=\sqrt{\pi^2/3+2\log^22}>\pi/\sqrt3$. Part~II proves $s_0(n)/p(n)=\sum_{j\le R}c_j\beta^j+O(\beta^{R+1})$ for every $R$, $\beta=\pi/\sqrt{6n}$, $c_0=\frac12$, $c_1=-\frac18$, with integer inverse windows: a second proof of Conjecture~5.3, at $r=0$ only. Part~III proves $\log A(n)=B\sqrt n+O(n^{1/3})$ with $B=\sqrt{\pi^2/3+4\log^22}=2.2829\ldots$, which replaces Part~I's lower rate. Part~IV proves $\log A(n)=B\sqrt n-Dn^{1/6}+o(n^{1/6})$, $D=\zeta\log2\,(\pi^2/12+\log^22)^{-1/6}$ with $-\zeta$ the first zero of $\operatorname{Ai}$, using Mallein's published killed-walk estimates, and the limit shape of a uniform superdiagonal partition. Part~V proves $A(n)\sim Kn^{-11/12}\exp(B\sqrt n-Dn^{1/6})$ with an explicit $K$, \emph{conditionally} on the unrefereed preprint arXiv:2607.27504v1 of Yu-Tian Li; the unconditional superdiagonal statements are those of Parts~III and~IV. The write adds a proof that partitions with choosable initial intervals are exactly the superdiagonal ones, which settles Gus Wiseman's conjecture in A238873 (the same argument was posted in A387112 by Clément Garrot on 31~August 2026), and the defect $\mathbb P_n(R\le0)-s_0(n)/p(n)=\beta/4+(\frac7{16}-\frac3{2\pi^2})\beta^2+O(\beta^3)$, from Part~II. Effective constants, fluctuations, all orders for $A(n)$ and an unconditional amplitude remain open. The report is unrefereed, and nothing in it is formalized.
\end{abstract}
\hypersetup{pageanchor=true}
{\small\setlength{\parskip}{0pt}\tableofcontents}
\newpage

\section*{Guide to this report}\phantomsection\label{dgp:sec:guide}
\addcontentsline{toc}{section}{Guide to this report}
The five Parts are five manuscripts of the session bundle of Reports 1--243 (arrival commit \texttt{60f54ea06}), placed in this directory by commit \texttt{3988bf5c4} (batch~107, cluster 107-DIAGP). All five write partitions in weakly increasing order $\lambda_1\le\cdots\le\lambda_k$ and study the classes of Archibald, Blecher, Elizalde and Knopfmacher \cite{abek}:
\[
 s_r(n)=\#\{\lambda\vdash n:\lambda_i\le i+r\ \forall i\},\qquad A(n)=\#\{\lambda\vdash n:\lambda_i\ge i\ \forall i\},
\]
with $s(n)=s_0(n)=\mathrm{A238875}(n)$ and $A(n)=\mathrm{A238873}(n)$, the empty partition counted at $n=0$.
\begin{center}
\small
\begin{tabular}{@{}l l >{\raggedright\arraybackslash}p{0.47\textwidth} l@{}}
\toprule
Part & Bundle report & Delivered title & Labels\\
\midrule
I & 156 (base) & Proofs of two conjectures on subdiagonal and superdiagonal partitions: Quantitative rank coupling and exponential separation & \texttt{dgp:conj:}\\
II & 159 & All orders asymptotics for subdiagonal partitions: Endpoint reductions and integer inverse localization & \texttt{dgp:sub:}\\
III & 158 & Exact exponential growth of increasing superdiagonal partitions: An elementary proof and quantitative threshold inversion & \texttt{dgp:rate:}\\
IV & 160 & Airy correction and limit shape for superdiagonal partitions: Pointwise enumeration and macroscopic structure & \texttt{dgp:airy:}\\
V & 162 & Full prefactor and refined inverse for superdiagonal partitions: Exact coefficient asymptotics for A238873 \textbf{(conditional; see below)} & \texttt{dgp:pre:}\\
\bottomrule
\end{tabular}
\end{center}
\emph{Arrangement.} The base of the merge is Report~156: it is the only manuscript that treats both classes, every shift $r\ge0$, and both conjectures of \cite[Section~5]{abek}; its \file{Report156.tex} was staged as this \file{article.tex}, and it is printed first, as Part~I. The other four follow in dependency order, grouped by class. Part~II (subdiagonal, $r=0$ only) cites Part~I and sharpens it at $r=0$. Parts~III--V are superdiagonal: Part~III replaces Part~I's lower rate by the exact rate; Part~IV adds the Airy term and uses Part~III's coarse lower bound for its limit-shape theorem; Part~V, whose proof is independent of Parts~I--IV but refines their statements under an extra, unrefereed input, comes last. The bundle numbers do not follow this order (159 is Part~II).

\emph{Part~V is conditional.} Part~V's Theorem~\ref{dgp:pre:thm:main} and Corollary~\ref{dgp:pre:cor:inverse} use Theorem~1 and Corollary~2 of Yu-Tian Li's preprint arXiv:2607.27504v1 (submitted 29~July 2026) \cite{li}, which is unrefereed and whose proof nobody at intake or write verified. They are \emph{conditional} on that preprint, and are labelled so in their statements. The unconditional superdiagonal statements of this report are those of Parts~III and~IV (whose external inputs are elementary or published).

\emph{Numbering.} Sections are numbered continuously, and statements and equations within sections. Part~I keeps Report~156's Sections~1--9 and all its numbers. Report~159's Section~$k$ is Section~$k+9$ here, Report~158's is $k+19$, Report~160's is $k+29$ and Report~162's is $k+38$; a statement or equation $k.j$ of a manuscript becomes $(k+\text{offset}).j$. For example, Report~159's Theorem~1.1 is Theorem~\ref{dgp:sub:thm:allorders}, Report~158's Theorem~1.1 is Theorem~\ref{dgp:rate:thm:main}, Report~160's Theorems~1.1 and~1.3 are Theorems~\ref{dgp:airy:thm:airy} and~\ref{dgp:airy:thm:shape}, and Report~162's Theorem~1.1 is Theorem~\ref{dgp:pre:thm:main}. Subsection~\ref{dgp:conj:sub:choosable} (with Proposition~\ref{dgp:conj:prop:choosable}) and Section~\ref{dgp:sec:further} were added in the write; no delivered number moved.

Notes added in the write are set in small type and tagged \textbf{[write]}, with their date (6~October 2026). They give cross-references between the Parts, mark questions that a later Part answers and statements that a later Part replaces, mark restated material, and record what the write checked. No delivered statement was changed. The write's interventions in delivered text are listed under ``Provenance and merge decisions'': six restated proofs replaced by pointers (each disclosed in place), a ``conditional'' tag in two statement headings of Part~V, unified citation keys, and two \LaTeX\ macro names. Each Part ends with its manuscript's own limits or questions; Section~\ref{dgp:sec:further} gathers every open question and non-claim for the whole report under Vladimir's standing rule of 4~October 2026.

\section*{What is proved, and where}\phantomsection\label{dgp:sec:status}
\addcontentsline{toc}{section}{What is proved, and where}
Here $p(n)$ and $q(n)$ count all partitions and partitions into distinct parts; $R$ is Dyson's rank (largest part minus number of parts) of a uniform partition of $n$; $c=\pi/\sqrt6$ and $\beta=c/\sqrt n=\pi/\sqrt{6n}$; $b=\log2$, $C=\pi^2/12+b^2$, $B=2\sqrt C=\sqrt{\pi^2/3+4\log^22}=2.2829\ldots$; $a_1=-2.3381\ldots$ is the largest zero of $\operatorname{Ai}$, $\zeta=-a_1$, and $D=\zeta bC^{-1/6}=1.5507\ldots$.
\begingroup\small
\renewcommand{\theHtable}{status.\arabic{table}}
\begin{longtable}{@{}>{\raggedright\arraybackslash}p{0.60\textwidth} >{\raggedright\arraybackslash}p{0.35\textwidth}@{}}
\toprule
Statement & Where\\
\midrule
\endhead
\multicolumn{2}{@{}l}{\emph{Subdiagonal partitions (unconditional)}}\\
$0\le\Prb_n(R\le r)-s_r(n)/p(n)\le C'/\log n$ uniformly in $r\ge0$; hence $s(n)/p(n)\to\frac12$: \textbf{ABEK Conjecture~5.3 proved} & Part~I, Theorem~\ref{dgp:conj:thm:coupling}\\
Uniform logistic crossover $s_r(n)/p(n)=(1+e^{-\pi r/\sqrt{6n}})^{-1}+O(1/\log n)$ (uses Dousse--Mertens) & Part~I, Corollary~\ref{dgp:conj:cor:logistic}\\
Lambert-$W$ inverse of $s$ with window $O(\log y/\log\log y)$ & Part~I, Theorem~\ref{dgp:conj:thm:inverse}\\
$s(n)/p(n)=\sum_{j\le R}c_j\beta^j+O_R(\beta^{R+1})$ for every $R$, $c_j\in\mathbb Q[\pi^{-2}]$, $c_0..c_3$ explicit ($c_1=-\frac18$): a second, sharper proof of Conjecture~5.3, \emph{at $r=0$ only} (uses Liu--Zhou and Zhou) & Part~II, Theorem~\ref{dgp:sub:thm:allorders}\\
Integer inverse windows $\lceil X_J\mp C_Jx_0^{-J/2}\rceil$ & Part~II, Theorem~\ref{dgp:sub:thm:inverse}\\
$\Prb_n(R\le0)-s(n)/p(n)=\frac\beta4+\bigl(\frac7{16}-\frac3{2\pi^2}\bigr)\beta^2+O(\beta^3)$ & \wtag\ note at Part~I's Section~\ref{dgp:conj:sec:questions}, \eqref{dgp:conj:eq:defect}\\
\midrule
\multicolumn{2}{@{}l}{\emph{Superdiagonal partitions (unconditional)}}\\
$\liminf\log A(n)/\sqrt n\ge B_*=\sqrt{\pi^2/3+2\log^22}=2.0617\ldots>\pi/\sqrt3$, so $q(n)/A(n)\to0$: \textbf{ABEK Conjecture~5.2 proved} & Part~I, Theorem~\ref{dgp:conj:thm:super}\\
$C/t-O(t^{-1/3})\le\log\sum_nA(n)e^{-tn}\le C/t+O(1)$; $B\sqrt n-O(n^{1/3})\le\log A(n)\le B\sqrt n+O(1)$: the exact rate $B$, \textbf{replacing Part~I's lower rate $B_*$}; inverse to $O((\log y)^{5/3})$ & Part~III, Theorems~\ref{dgp:rate:thm:main} and~\ref{dgp:rate:thm:inverse}\\
$\log A(n)=B\sqrt n-Dn^{1/6}+o(n^{1/6})$ and its radial analogue (uses Mallein's published estimates); inverse $T(y)=x+\zeta bC^{-2/3}x^{2/3}+o(x^{2/3})$, $x=(\log y/B)^2$ & Part~IV, Theorem~\ref{dgp:airy:thm:airy}, Corollary~\ref{dgp:airy:cor:inverse}\\
Limit shape: $tK\to2b$, $K/\sqrt n\to2\log2/\sqrt C=1.2145\ldots$; cumulative and row profiles $g$, $f$; half the limiting rows on the diagonal (macroscopically) & Part~IV, Theorem~\ref{dgp:airy:thm:shape}, Figure~\ref{dgp:airy:fig:shape}\\
\midrule
\multicolumn{2}{@{}l}{\emph{Superdiagonal partitions, \textbf{conditional on arXiv:2607.27504v1}}}\\
$A(n)\sim Kn^{-11/12}\exp(B\sqrt n-Dn^{1/6})$, $K=\sqrt2C^{5/12}/(\sqrt\pi\operatorname{Ai}'(a_1))=1.2705\ldots$; the radial amplitude $K_F$ and first radial correction $\alpha$; inverse to the $\sqrt x$ term & Part~V, Theorem~\ref{dgp:pre:thm:main}, Corollary~\ref{dgp:pre:cor:inverse}\\
\midrule
\multicolumn{2}{@{}l}{\emph{OEIS}}\\
A partition has choosable initial intervals iff it is superdiagonal, so their number is $A(n)$: Wiseman's A238873 conjecture (same argument published by Clément Garrot in A387112) & \wtag\ Proposition~\ref{dgp:conj:prop:choosable}\\
\midrule
\multicolumn{2}{@{}l}{\emph{Open}}\\
\multicolumn{2}{@{}p{0.96\textwidth}@{}}{The defect for fixed $r\ge1$ and growing $r$; $r/\sqrt n$-sensitive errors; effective onsets everywhere; convergence and higher coefficients of Part~II's series; superdiagonal fluctuations, discrete diagonal contacts and extreme rows; a pointwise relative correction and all orders for $A(n)$; an $x^{1/3}$ inverse term; an \emph{unconditional} proof of the amplitude $K$; literature. See Section~\ref{dgp:sec:further}.}\\
\bottomrule
\end{longtable}
\addtocounter{table}{-1}% the uncaptioned longtable must not consume a table number
\endgroup
The five companions check finite mechanisms in exact arithmetic (counts, injections, identities, finite inequalities, coefficient algebra); no asymptotic statement rests on them.

\section*{The two journal conjectures, and the OEIS entries}\phantomsection\label{dgp:sec:conjectures}
\addcontentsline{toc}{section}{The two journal conjectures, and the OEIS entries}
\emph{The conjectures.} M.~Archibald, A.~Blecher, S.~Elizalde and A.~Knopfmacher, \emph{Subdiagonal and superdiagonal partitions}, Afrika Matematika \textbf{36}, article~77 (2025), received 12~November 2024, accepted 27~February 2025, published online 7~April 2025, open access under CC~BY~4.0 \cite{abek}. In their Section~5 (``Asymptotics'', p.~15), $P_n$ is the set of weakly increasing partitions of $n$, $A_n\subseteq P_n$ and $S_n\subseteq P_n$ are the superdiagonal ($\lambda_i\ge i$) and subdiagonal ($\lambda_i\le i$) ones, and $Q_n\subseteq P_n$ is the subset of partitions with distinct parts, with the remark $Q_n\subseteq A_n$. The two conjectures read, verbatim (the write read the open-access text on 6~October 2026):
\begin{quote}
\textbf{Conjecture 5.2}\quad $\displaystyle\lim_{n\to\infty}\frac{|Q_n|}{|A_n|}=0.$

``We conjecture that, asymptotically, half of the partitions are subdiagonal.''\\
\textbf{Conjecture 5.3}\quad $\displaystyle\lim_{n\to\infty}\frac{|S_n|}{|P_n|}=\frac12.$
\end{quote}
After each, the paper says that similar problems were studied in its references [4] and [5]. It proves $|A_n|/|P_n|\to0$ (Proposition~5.1), that balanced partitions (largest part equal to the number of parts) have density zero (Lemma~5.4, by the no-ones injection that Part~I's Lemma~\ref{dgp:conj:lem:atom} extends to every rank), and $\limsup|S_n|/|P_n|\le\frac12$ (Proposition~5.5). Its closing paragraph says that it ``would be interesting'' to describe the asymptotic behaviour of $|A_n|$ similarly. These are conjectures of a refereed journal article, not OEIS conjectures; neither OEIS entry states them.

\emph{What this report proves about them.} Part~I proves both: Conjecture~5.3 by Theorem~\ref{dgp:conj:thm:coupling} (with $|S_n|=s_0(n)$, $|P_n|=p(n)$, and the rate $O(1/\log n)$), and Conjecture~5.2 by Theorem~\ref{dgp:conj:thm:super} ($|Q_n|=q(n)$, $|A_n|=A(n)$; Theorem~\ref{dgp:conj:thm:super} gives $q(n)/A(n)\le\exp(-(B_*-\pi/\sqrt3-\eta)\sqrt n)$ eventually, for every $\eta>0$) \textup{(wording corrected after the independent check of 6~October 2026; the first wording read ``the ratio decays like $\exp(-(B_*-\pi/\sqrt3-\eta)\sqrt n)$'', but the theorem gives an upper bound only)}. The write read both proofs in full and found no gap; the intake re-derived their key steps (the interior lemma, the fixed-rank injection, the prefix--tail injection and the tail comparison). Part~II proves Conjecture~5.3 a second time, at $r=0$ only, with the full expansion. Part~III answers the paper's closing question at the logarithmic scale (the limit of $\log A(n)/\sqrt n$ exists and equals $B$); Part~IV adds the next term, unconditionally; Part~V gives the full equivalent, conditionally. With Part~III's rate, Conjecture~5.2 holds with the exponent $B-\pi/\sqrt3=0.4691\ldots$ in place of Part~I's $B_*-\pi/\sqrt3=0.2479\ldots$.

\emph{The OEIS entries}, as read by the write on 6~October 2026 (A238873 revision \#36 of Oct~07 2025; A238875 revision \#33 of Apr~14 2025, the revisions Part~I's bibliography records):
\begin{itemize}
\item \textbf{A238875} (Joerg Arndt, Mar 24 2014), ``Subdiagonal partitions: number of partitions (p1, p2, p3, ...) of n with pi <= i.'', with the comment that the partitions are weakly increasing lists, a g.f.\ via A129176, and the link to \cite{abek} (``See p.~9''). No conjecture.
\item \textbf{A238873} (Joerg Arndt, Mar 23 2014), ``Number of superdiagonal partitions: partitions (p1, p2, p3, ...) of n such that pi >= i.'', with a comment that the weakly decreasing convention gives A003114 (Gus Wiseman, Oct 06 2025), the link to \cite{abek} (``See p.~5''), and one conjecture, unrelated to those of \cite{abek}, which reads verbatim:\\
\verb!Conjecture: Also the number of integer partitions y of n with choosable!\\
\verb!initial intervals, meaning it is possible to choose distinct positive!\\
\verb!integers (z_1, z_2, ...) such that z_i <= y_i for all i.!\\
\verb!- _Gus Wiseman_, Sep 26 2025!
\end{itemize}
No manuscript addresses Wiseman's conjecture. It is true: Proposition~\ref{dgp:conj:prop:choosable} (added by the write, with a complete proof) shows that a partition has choosable initial intervals if and only if its increasing arrangement is superdiagonal. The argument is not new: on 31~August 2026 Clément Garrot added to OEIS A387112 (revision \#19 of Sep~07 2026) a comment with the same proof \cite{oeisA387112}, stating that it also proves the corresponding conjectures in A238873, A387118 and A388711. The A238873 entry still displays the line as a conjecture. The intake record proposed the argument without knowing of Garrot's comment; the write found it. Nothing was submitted to the OEIS.

\emph{The inverses} (relation to the transseries volume \cite{TSvol}). Statement by statement:
\begin{itemize}
\item Part~I, Theorem~\ref{dgp:conj:thm:inverse}: the centre $x_0(y)$ of \eqref{dgp:conj:eq:inversecenter} is an \emph{instance} of the volume's Lambert core \file{p0:thm:lambert-core} after the volume's monomial reduction \file{p0:lem:alpha-reduction} with $\alpha=\frac12$: the model $M(x)=e^{2c\sqrt x}/(8\sqrt3\,x)$ is of exponential--power type (\file{p0:def:model}) with data $(1/(8\sqrt3),2c,\frac12,-1,\cdot,0)$; in $\xi=\sqrt x$ the equation $M=y$ is $2c\xi-2\log\xi=\log y+\log(8\sqrt3)$, the core with $a=2c$, $b=-2<0$, whose large solution \file{p0:eq:core-lambert-minus} is $\xi=c^{-1}(-W_{-1}(-c/\sqrt{8\sqrt3\,y}))$, i.e.\ $x_0=\xi^2$. The window $O(\sqrt{x_0}/\log x_0)$, the bracketing and the nested-log initializer \eqref{dgp:conj:eq:loginverse} are Part~I's own.
\item Part~II, Theorem~\ref{dgp:sub:thm:inverse}: the same centre $x_0$ (an instance as above); the centres $X_J$ are, after the same reduction, the reversion of \file{p0:thm:lambert-centered} for the model $M_J$ of \eqref{dgp:sub:eq:MJ}, whose data in $\xi$ are $a=2c$, $b=-2$ and $Q(\tau)=\log(1+\sum_ja_jc^j\tau^j)$, squared back by \file{p0:eq:alpha-inverse-transport}: an \emph{instance} of the coefficient formula. The write checked the first two coefficients against the volume's recurrence \file{p0:eq:c-bell} ($n=1,2$): it gives $b_0=h$ and $b_1=ht-a_2+h^2/2$, Part~II's values. The ceiling windows and Lemma~\ref{dgp:sub:lem:window} are Part~II's own, and they are \emph{not} an instance of \file{p0:thm:staircase}: that theorem presupposes an admissible interpolation of the sequence, while Part~II brackets the root of a smooth model of $s(n)$ directly at the integers.
\item Parts~III, IV and~V, Theorem~\ref{dgp:rate:thm:inverse} and Corollaries~\ref{dgp:airy:cor:inverse} and~\ref{dgp:pre:cor:inverse}: not instances of anything in the volume. Part~III has only an $O(n^{1/3})$ error, and the logarithms of Parts~IV and~V contain the growing fractional power $n^{1/6}$, outside \file{p0:def:model} (whose corrections are $b\log x$ and a series in $x^{-\delta}$). Each inverse is proved directly by monotonicity and integer bracketing. The power-and-logarithm balance of Part~V's inverse is at most analogous to the volume's monomial--logarithmic template \file{plt:thm:lw-template}.
\end{itemize}
None of the manuscripts claims novelty for the inversion mechanics.

\section*{Notation across the five Parts}\phantomsection\label{dgp:sec:notation}
\addcontentsline{toc}{section}{Notation across the five Parts}
\emph{Common to all Parts.} Increasing order; $A(n)$ the superdiagonal count with $A(0)=1$; $Z_j$ (Part~I: $Z_1$) the multiplicity of the part $j$; $T(y)=\min\{n\ge0:\cdot\ge y\}$ a first threshold; $p(n)$ the partition number. \emph{Fixed for this guide and Section~\ref{dgp:sec:further}:} $s_r$, $s=s_0$, $c=\pi/\sqrt6$, $\beta=c/\sqrt n$, $b=\log2$, $C=\pi^2/12+b^2$, $B=2\sqrt C$, $a_1<0$ the largest zero of $\operatorname{Ai}$, $\zeta=-a_1>0$, $D=\zeta bC^{-1/6}$. No symbol of a manuscript was renamed by the write (only two macro names, which do not change the printed text). Part~V itself renames Li's $A_q(z)$ to $H_q(y)$ because $A(n)$ is taken; Li's scaled argument $\zeta$ is Part~V's $s$. The symbols below change meaning between the Parts or inside one; each Part fixes its own meaning, and opens with a short table of its readings.
\begingroup\small
\renewcommand{\theHtable}{notation.\arabic{table}}
\begin{longtable}{@{}>{\raggedright\arraybackslash}p{0.11\textwidth} >{\raggedright\arraybackslash}p{0.85\textwidth}@{}}
\toprule
Symbol & Meanings\\
\midrule
\endhead
$b$ & Part~I: $b=\pi/\sqrt3$, the rate of $q(n)$ \eqref{dgp:conj:eq:constants}. Part~II: Zhou's variable $b=\pi/\sqrt{6(n-1/24)}$ (Section~\ref{dgp:sub:sec:analytic}). Parts~III--V: $b=\log2$. \emph{Tempting false reading:} Part~I's $b$ is not $\log2$.\\
$B$, $\mathcal B$ & Part~I: lower rates $B_*$, $B_0$, $B_h$ (\eqref{dgp:conj:eq:constants}, \eqref{dgp:conj:eq:baseline}, \eqref{dgp:conj:eq:Bh}). Part~II: $B(r)$ a cubic rank coefficient \eqref{dgp:sub:eq:ab}; $B$ a number of top rows and $B_{\rm src}$, $B_{\rm tgt}$ complements (Lemma~\ref{dgp:sub:lem:mixed}); $\mathcal B_n$, $\mathcal B(x;n)$ bottom events. Parts~III and~V: $B=2\sqrt C$, the exact rate (Part~IV writes $2\sqrt C$, and $B$ in its inverse); Part~III also $B_k$ tail heights; Part~V $\mathcal B_t$ an index block. \emph{False reading:} Part~I's $B_*$ is a lower bound, not the rate.\\
$C$, $\mathcal C$ & Part~I: $C$ a generic constant (Theorem~\ref{dgp:conj:thm:coupling}); $C_M$ Catalan numbers, $\mathcal C_M$ prefixes. Part~II: $C_i$ cumulative multiplicities; $C_J$, $C$ inverse constants. Parts~III--V: $C=\pi^2/12+\log^22$ fixed; $C_j$ cumulative multiplicities. Part~IV reserves $C$ and writes generic constants $c_*$, $K_*$.\\
$t$, $\beta$ & Parts~I, III, IV, V: $t$ a Boltzmann (radial) variable, $e^{-t}$ the activity (Part~I in Section~\ref{dgp:conj:sec:interior}). Part~II: $t=c^{-2}=6/\pi^2$, and its radial variable is $\beta=c/\sqrt n$ (also $\beta_x$, $\beta_0$). Part~I: $\beta=c/\sqrt n$ (proof of Corollary~\ref{dgp:conj:cor:logistic}). Parts~IV, V: $\beta=d-2b-\log(1-e^{-d})>0$ a buffer constant. \emph{False reading:} Part~II's $t$ is a number, not a variable.\\
$\zeta$, $a$ & Part~IV: $\zeta=-a_1\approx2.3381$. Part~V: $a=a_1<0$, $a_2$; Li's preprint: $\zeta$ the scaled argument (Part~V's $s$). Part~I: $a_h=\log(\rho_h/2)$. Part~II: $a=\frac3{16}-\frac t4$ \eqref{dgp:sub:eq:ab}, $a_j$ model coefficients. Part~IV: $a$, $d_0$ interval ends, $a(t)=t^{-19/12}$ \eqref{dgp:airy:eq:ascale}. \emph{False reading:} Part~IV's $-\zeta b$ and Part~V's $ba$ are the same number.\\
$D$ & Part~I: $D(n)=p(n)-p(n-1)$; $D_h(M)$ Dyck counts; $D$ a constant in the proof of Theorem~\ref{dgp:conj:thm:inverse}. Part~II: $D_k(n,a)$ a counting table; $D_J$ constants. Part~III: $D_h(M)$; $D$ a constant in the inverse. Part~IV: $D_M(t)=e^{2bM-tM(M+1)/2}$; $D=\zeta bC^{-1/6}$ (Section~\ref{dgp:airy:sec:pointwise}). Part~V: $D=-abC^{-1/6}$, the same number; $D_M=4^Me^{-tM(M+1)/2}$, equal to Part~IV's $D_M(t)$; $d$ a cutoff $>b$ and $d$ an inverse coefficient \eqref{dgp:pre:eq:inversecoeffs}.\\
$K$ & Parts~I, II: $K$ the length cutoff (\eqref{dgp:conj:eq:scales}, \eqref{dgp:sub:eq:cutoffs}). Part~III: constants $K_0,\dots,K_3$. Part~IV: $K$ the number of parts; constants $K_0,\dots,K_8$. Part~V: the amplitudes $K$, $K_F$; $K_N$ a length cap; $\mathcal K$ a compact set. \emph{False reading:} Part~IV's $K$ is a length, Part~V's an amplitude.\\
$L$ & Parts~I, II: the largest part; Part~II also $L_{x,y,r}(n)$. Parts~III, IV: $L(x)=x\log4-x^2/2+\operatorname{Li}_2(e^{-x})$ (the same function in both; Part~IV also $L$ a grid bound). Part~V: $L=2b=\log4$, as in Li's preprint.\\
$A$, $H$ & $A(n)$ throughout; \cite{abek}'s $A(q,t)$ its generating function; Part~II: Zhou's normalizer $A$; Part~I: $H=\log y$; Part~II: $H(x)$ the Rademacher term \eqref{dgp:sub:eq:H}; Parts~IV, V: heights $H_j=j-C_j$; Part~V: $H_q(y)$ Ramanujan's entire function (Li's $A_q$); $\mathcal A_M$ an area (Parts~IV, V).\\
$N$ & Parts~I, II: $N(r,n)$ rank counts. Part~IV: $N=J-M$; $\mathcal N$ the size. Part~V: $N=\sum jZ_j$ the size; $N_t(s)$ Li's normalization \eqref{dgp:pre:eq:N}; $\mathcal N(0,2C)$ a normal law.\\
$s$, $q$ & Part~I: $s_r(n)$; $q(n)$ distinct partitions. Part~II: $s(n)=s_0(n)$; $q=e^{-\beta}$; $s$ an index. Part~III: $q$ rational. Part~IV: $s$ a tilt; $q$ Mallein's potential. Part~V: $s$ the scaled Airy argument, $s_t$ the first-zero coordinate; $q=e^{-t}$.\\
$R$, $\rho$ & Parts~I, II: $R$ Dyson's rank; Part~II also $R$ a truncation order. Parts~I, III: $\rho_h=4\cos^2(\pi/(h+2))$. Part~IV: $\rho$ a ramp parameter. Part~V: $\rho=\rho_1<\rho_2<\cdots$ the zeros of $H_q$; $R=y_t(s)$ a radius.\\
$M$, $m$ & Part~I: $M$ a prefix length, $M(x)$ a model; $m$ the interior scale. Part~II: $M$ the endpoint cutoff, $M_J(x)$ models; $m$ as in Part~I. Parts~III--V: $M\approx b/t$ the cutoff. Part~IV: $m(x)$ the limiting mean profile; $m$ a walk length. Part~V: $m_t$ a centre; $m=\lfloor1/t\rfloor$; $m$ a target in Proposition~\ref{dgp:pre:prop:llt}.\\
$f$, $g$, $h$ & Part~I: $f_M(N)$ distinct tails; $f(x)=-\log(1-e^{-tx})$; $g=\frac14\operatorname{sech}^2(x/2)$; $h(x)=\log M(x)$; $h$ a height. Part~II: $f(z)=1/(1+e^z)$; $h=\frac7{24}+\frac t2$. Part~III: $f(x)=-\log(1-e^{-x})$; $h$ a height. Part~IV: $g$ and $f=g^{-1}$ the limit profiles \eqref{dgp:airy:eq:g}--\eqref{dgp:airy:eq:f}; $h(x)=-\log(1-e^{-x})$. Part~V: $f(x)=-\log(1-e^{-x})$; $g_j$, $G_q$ product coefficients. \emph{False reading:} Part~IV's $f$ is the row profile, not Part~III's $f$ (Part~IV's $h$).\\
$W$, $G$, $Q$ & Parts~I, III: $W_h(M)$ a weight bound; Part~I also $W_{-1}$ (Lambert). Part~IV: $W_P$, $W_T$, $\mathcal W_t$. Part~II: $G_b$ gaps; $Q_\beta$ the Boltzmann law. Part~III: $G_M(t)$ a tail sum. Part~IV: $G$ a buffer; $Q_{t,M}(h)$ survival; $\mathbb Q_t$ a product law. Part~V: $G$ a buffer, $G_q$; $Q_1$, $Q_2$ Airy densities.\\
$p$, $P$, $\mathbb P$ & $p(n)$ partitions (Parts~I, II, and in the guide); $p$ a geometric parameter (Parts~III--V), $p_*$; Part~V also $p=11/12$. Part~II: $P_r(n)$ cumulative rank counts; Parts~I, II: $P(q)$ the partition product; Parts~IV, V: $P_M$ a tail product. $\mathbb P_n$: uniform law on partitions of $n$ (Parts~I, II) or on superdiagonal partitions of $n$ (Part~IV); $\mathbb P_t$ Boltzmann laws (Parts~III--V).\\
$E$, $\alpha$ & Part~II: $E_n$, $E_n^{(M)}$ failure counts. Part~III: $E(t)$ an error. Part~IV: $E_t$, $E_H$ events. Part~V: $E_0=\pi^2/12-b^2$; $\alpha=a^2(\frac14-\frac b{30})$. Part~I: $\alpha=\lim r_n/\sqrt n$.\\
\bottomrule
\end{longtable}
\addtocounter{table}{-1}% the uncaptioned longtable must not consume a table number
\endgroup

\section*{Provenance and merge decisions}\phantomsection\label{dgp:sec:provenance}
\addcontentsline{toc}{section}{Provenance and merge decisions}
Five manuscripts of the session bundle of Reports 1--243. Their author lines read ``Report 156'', ``Report 159'', ``Report 158'', ``Report 160'' and ``Report 162''; none names an author, a tool or an addressee, each sets an empty PDF author field, and none carries ``prepared for private review'' wording. All are dated 3~October 2026, pin no ProveIt commit and contain no repository path. They arrived in commit \texttt{60f54ea06} and were placed by \texttt{3988bf5c4} (batch~107, cluster 107-DIAGP).
\begin{itemize}
\item \textbf{Part~I}, Report~156, archive \file{Two_Partition_Conjectures_Proofs_and_Quantitative_Extensions_Source.zip} (\file{Report156.tex}, 861 lines, 14 pages); the base. Contributes both conjecture proofs, the uniform shifted coupling and fixed-rank atom bound, the logistic crossover, the Lambert inverse, the Catalan and bounded-height prefix lower bounds.
\item \textbf{Part~II}, Report~159, archive \file{Subdiagonal_Partitions_All_Orders_Expansion_and_Integer_Inverses_Source.zip} (\file{Report159.tex}, 1019 lines, 17 pages). Contributes the all-orders expansion of $s(n)/p(n)$ with $c_0,\dots,c_3$, the mixed row--column bijection, and the integer inverse windows.
\item \textbf{Part~III}, Report~158, archive \file{Superdiagonal_Partitions_Exact_Exponential_Growth_and_Inverse_Source.zip} (\file{Report158.tex}, 749 lines, 13 pages). Contributes the exact rate $B$ with the tilted upper bound, the survival lemma, the uniform spectral bound, the elementary coefficient transfer and the inverse.
\item \textbf{Part~IV}, Report~160, archive \file{Superdiagonal_Partitions_Airy_Correction_and_Limit_Shape_Source.zip} (\file{Report160.tex}, 1308 lines, 17 pages). Contributes the Airy term (radial and pointwise), the exact-weight tail lemma, the second-order inverse, and the length and shape theorem.
\item \textbf{Part~V}, Report~162, archive \file{Superdiagonal_Partitions_Full_Prefactor_and_Refined_Inverse_Source.zip} (\file{Report162.tex}, 933 lines, 15 pages). Contributes, conditionally, the radial amplitude and first correction, the pointwise equivalent, the refined inverse; unconditionally, the first-pole transfer lemma, the cutoff lemma and the local limit theorem, and the two source-display repairs.
\end{itemize}
No package embeds another (no byte-identical file across the five). Where the merge had to choose:
\begin{enumerate}
\item \emph{Order and base}: as in the Guide.
\item \emph{Restatements printed once.} The barrier lemma, the survival lemma, the monotonicity injection and the dilogarithm evaluation $L(b)=C$ are first printed in Part~III, and the Dyck representation in Part~I. Six later \emph{proofs} of them are replaced by pointers, each with a dated note saying what is not reprinted: Part~III's proof of Lemma~\ref{dgp:rate:lem:dyck} (Part~I, Sections~\ref{dgp:conj:sec:prefix}--\ref{dgp:conj:sec:height}); Part~IV's proof of the barrier equivalence (Section~\ref{dgp:airy:sec:results}), of the survival bound (Section~\ref{dgp:airy:sec:radial}) and of the dilogarithm evaluation (Section~\ref{dgp:airy:sec:representation}), and Part~IV's and Part~V's monotonicity injections (Sections~\ref{dgp:airy:sec:pointwise} and~\ref{dgp:pre:sec:inverse}). Every statement and every labelled display stays. Restatements woven into a longer proof stay, with a note: Part~II's interior step (Section~\ref{dgp:sub:sec:localize}; Part~I's Lemma~\ref{dgp:conj:lem:interior}), Part~V's barrier sentence, its summation by parts \eqref{dgp:pre:eq:sumbyparts} (Part~IV's \eqref{dgp:airy:eq:prefixweight}) and its ruin bound (Part~III's Lemma~\ref{dgp:rate:lem:survival}, Part~IV's \eqref{dgp:airy:eq:ruin}). The Fourier and local-limit arguments of Part~IV (Section~\ref{dgp:airy:sec:fourier}) and Part~V (Section~\ref{dgp:pre:sec:local}) use the same buffer, ruin and three-arc cover but are assembled differently (a tilted local limit theorem with exact weights, against a moment-generating-function limit with total-variation constraint removal); both are printed, as two routes.
\item \emph{Conditional status.} The headings of Part~V's Theorem~\ref{dgp:pre:thm:main} and Corollary~\ref{dgp:pre:cor:inverse} carry the added tag ``\wtag\ conditional on arXiv:2607.27504v1''; nothing else in them changed.
\item \emph{Front matter and bibliography.} The five title blocks, abstracts and scope paragraphs are printed at the heads of their Parts; one table of contents replaces five. The five bibliographies are merged into one: an entry cited by several Parts is printed once, with every Part's annotation kept verbatim; citation keys were unified (for example Part~I's \texttt{oeissuper} and Part~III's \texttt{oeis} are both \texttt{oeisA238873}; Part~I's \texttt{dlmf} is DLMF~\S27.14, Part~IV's is DLMF Table~9.9.1, now two keys). The entries ``Report~156'' and ``Report~158'' point to Parts~I and~III. The write added the entries for the transseries volume and for OEIS A387112, and dated annotations to the entries of the papers it read.
\item \emph{Macros.} Parts~I--II define \verb!\PP! as $\mathbb P_n$ and Parts~III--V as $\mathbb P$; Part~III defines \verb!\Li! as $\operatorname{Li}$ and Part~V as $\operatorname{Li}_2$. One definition each is kept and the uses were rewritten (\verb!\PP! as \verb!\PP_n! in Parts~I--II, \verb!\Li! as \verb!\Li_2! in Part~V); the printed text is unchanged.
\item \emph{Write additions.} This front matter, the reading-conventions tables of the Parts, Subsection~\ref{dgp:conj:sub:choosable} with Proposition~\ref{dgp:conj:prop:choosable}, the defect formula \eqref{dgp:conj:eq:defect}, Section~\ref{dgp:sec:further}, the dated notes, and the label prefixes. Everything else is delivered text.
\end{enumerate}

\section*{What was checked, and what was not}\phantomsection\label{dgp:sec:trust}
\addcontentsline{toc}{section}{What was checked, and what was not}
The report is unrefereed and not formalized. At intake (dossier of batch~107, 5~October 2026) the five manuscripts were read in full and no mathematical claim was found false; only two source-display repairs (Part~V's) and no wrong statement. The intake re-derived the key steps of both conjecture proofs by hand; recomputed $A(n)$ to $n=16000$ and $s(n)$ to $n=1500$ by its own dynamic programs (equal to the companions' exact tables for $n\le400$ and $n\le300$); recomputed every constant ($B_*$, $B_0$, $B$, $C$, $a_1$, $\operatorname{Ai}'(a_1)$, $D$, $K_F$, $K$, $\alpha$, $2D/B=\zeta bC^{-2/3}=1.3585577\ldots$, $2\log2/\sqrt C$) to 30 digits; checked that $\bigl(s/p-\sum_{j\le2}c_j\beta^j\bigr)/\beta^3$ drifts towards $c_3=-0.3162$ ($-0.544$ at $n=100$, $-0.388$ at $n=1500$); and checked the normalization of Li's Theorem~1 against Part~V. It fitted Part~V's constants by least squares (float dynamic program to $n=16000$): the pointwise residual $\log A(n)-B\sqrt n+Dn^{1/6}+\frac{11}{12}\log n-\log K$ on $n=2000,2250,\dots,16000$ in the basis $1,n^{-1/6},n^{-1/3}[,n^{-1/2}]$ gives constant terms $-0.038$ and $-0.026$; the radial residual on 23 equally spaced $t\in[0.016,0.06]$ in the basis $1,t^{1/3},t^{2/3}[,t]$ gives constant terms $-0.025$ and $-0.088$, and with the constant fixed at $0$ a $t^{1/3}$ coefficient $1.34$ against $\alpha=1.24$ (maximal residual $3.6\cdot10^{-4}$). So $K$, $K_F$ and $\alpha$ agree to the few-percent level these fits can see (an error of order $\log2$ in a normalization would be visible); the fits are ill-conditioned and prove nothing. The approach is slow: $\log A(n)/\sqrt n=2.1177$ at $n=8000$ against $B=2.2829$, which the term $-Dn^{1/6}$ explains.

The write (6~October 2026) re-read the five proofs and checked against the sources: both conjectures in the open-access text of \cite{abek}, and its Theorem~2.1, whose display sums from $k\ge1$ (Part~V's first repair); the statements of Dousse--Mertens' Theorem~1.2 \cite{dm}, Liu--Zhou's Theorem~1.3 \cite{lz}, Zhou's Theorem~4.1 \cite{zhou} (hypotheses at $k=2$, $m=0$, fixed $j=|r|$), Mallein's Lemmas~2.6 and~2.7 \cite{mallein} and the $q$-Airy example of \v{S}tampach--\v{S}\v{t}ov\'{\i}\v{c}ek \cite{ss} (Part~V's second repair), each in the arXiv version named in the notes, with the journal versions not compared; the arXiv record of Li's preprint (one version, v1 of 29~July 2026, no journal reference, on 6~October 2026). It recomputed the defect \eqref{dgp:conj:eq:defect} numerically (with exact $p(n)$ and exact rank counts) and checked Proposition~\ref{dgp:conj:prop:choosable} by exhaustive search for $n\le24$. On copies of the arrival archives (Windows, Python~3.14.4) the five \texttt{counts} commands regenerated the shipped count tables (after removing the carriage returns of Windows text output), and Part~V's 14 tests passed.

What was \emph{not} done: the analytic estimates (the Rademacher and rank inputs and their specialization, the endpoint reduction and mixed bijection to every order, the spectral, Chernoff, local-limit and Fourier estimates, the killed-walk specialization, the first-pole transfer) were read and their displayed algebra checked, but they are the manuscripts' proofs, not an independent verification. Nobody verified the proofs of the external theorems: Hardy--Ramanujan--Rademacher (DLMF), Dousse--Mertens, Liu--Zhou, Zhou, Mallein, \v{S}tampach--\v{S}\v{t}ov\'{\i}\v{c}ek, and Li. Li's Theorem~1 and Corollary~2 were read at intake, where their hypotheses were found to be met; the write did not read them again. Zhou's coefficient recipe \eqref{dgp:sub:eq:rankrecipe} was not re-derived from his Proposition~2.2.

\section*{Relation to neighbouring reports and formal projects}\phantomsection\label{dgp:sec:neighbours}
\addcontentsline{toc}{section}{Relation to neighbouring reports and formal projects}
No other report of the repository treats A238873, A238875, the paper \cite{abek}, superdiagonal or subdiagonal partitions in this sense, or the turning-point asymptotics of Ramanujan's entire function (repository search, 6~October 2026; the word ``superdiagonal'' elsewhere refers to matrices). The partition reports of this collection (for example \file{a022629-distinct-partition-norms}, \file{a033552-catalan-partitions}, \file{a097356-sqrt-restricted-partitions}, \file{a291698-moving-fugacity-partitions}) share only the generic Hardy--Ramanujan and Boltzmann methods. \file{a082161-airy-amplitudes} also meets the first Airy zero through a killed walk with an area or linear potential (its Dirichlet operator $-\partial_x^2+8x$); the mechanism is analogous, the theorems are not instances of each other. \file{a380274-mahonian-growing-powers/literature.md} cites R.~Zhang's paper on Ramanujan's entire function $A_q(z)$ for a different problem. The transseries volume \file{Analysis/Transseries/docs/series-and-transseries/Transseries_And_Inversion/transseries_and_inversion.tex} is related to the inverses as stated above. Placement in the collection confers no formal status: no Lean or Rocq file in the repository treats these sequences, and no statement of this report is formalized.
\clearpage

\part{Proofs of two conjectures on subdiagonal and superdiagonal partitions}\label{dgp:conj:part}
\partsource{Bundle Report~156, archive \file{Two_Partition_Conjectures_Proofs_and_Quantitative_Extensions_Source.zip}, delivered as \file{Report156.tex} (14 pages), the base of the merge. Delivered title block ``Proofs of two conjectures on subdiagonal and superdiagonal partitions / Quantitative rank coupling and exponential separation / Report 156 / 3 October 2026''. Printed as delivered, with \textbf{[write]} notes; section, statement and equation numbers are those of the delivery. Subsection~\ref{dgp:conj:sub:choosable} is added in the write.}
\begin{center}
\small
\begin{tabular}{@{}>{\raggedright\arraybackslash}p{0.46\textwidth} >{\raggedright\arraybackslash}p{0.46\textwidth}@{}}
\toprule
Reading conventions of this Part & Elsewhere in the report\\
\midrule
$b=\pi/\sqrt3$ (the rate of $q(n)$); $B_*$, $B_0$, $B_h$ lower rates & Parts~III--V: $b=\log2$, $B=2\sqrt C$ the exact rate\\
$c=\pi/\sqrt6$; $\beta=c/\sqrt n$ (Section~\ref{dgp:conj:sec:logistic}); $t=c/\sqrt n$ (Section~\ref{dgp:conj:sec:interior}) & Part~II: $t=6/\pi^2$, radial variable $\beta$\\
$C$ a generic constant; $C_M$ Catalan numbers; $D(n)=p(n)-p(n-1)$; $K$, $m$ cutoffs; $L$ the largest part & Parts~III--V: $C=\pi^2/12+\log^22$; Part~IV: $K$ the length\\
\bottomrule
\end{tabular}
\end{center}
\begin{partabstract}
For weakly increasing partitions of $n$, let $s_r(n)$ count those satisfying
$\lambda_i\leq i+r$ at every index, and let $A(n)$ count those satisfying
$\lambda_i\geq i$. We prove both conjectures in Section~5 of Archibald,
Blecher, Elizalde and Knopfmacher, \emph{Subdiagonal and superdiagonal
partitions}, Afrika Matematika \textbf{36}, 77 (2025).

An extension of their rank-zero injection, together with an elementary
interior estimate and the classical Hardy--Ramanujan formula, gives
\[
 0\leq\PP_n(R\leq r)-\frac{s_r(n)}{p(n)}=O((\log n)^{-1})
 \qquad(r\in\ZZ_{\geq0}),
\]
uniformly over all nonnegative integer shifts, where $R$ is Dyson's rank.
In particular $s_0(n)/p(n)=1/2+O((\log n)^{-1})$, proving their
Conjecture~5.3. We derive its Lambert $W$ inverse with the growing
localization error $O(\log y/\log\log y)$. A separate application of
the Dousse--Mertens rank theorem gives a uniform logistic crossover.

For superdiagonal partitions, Catalan prefixes followed by distinct
large parts give explicit finite lower bounds. Restricting the prefix
to fixed-height Dyck paths yields
\[
 \liminf_{n\to\infty}\frac{\log A(n)}{\sqrt n}
 \geq \sqrt{\pi^2/3+2(\log2)^2}>\frac\pi{\sqrt3}.
\]
Thus distinct partitions are exponentially negligible relative to
$A(n)$ on the $\sqrt n$ scale, proving Conjecture~5.2. This is a lower
bound, not an exact superdiagonal growth rate. No all-orders expansion
is established. An exact, independently implemented companion checks
sequence prefixes and the finite combinatorial mechanisms.
\end{partabstract}
\noindent\textbf{Attribution and scope.}
The definitions, generating functions, conjectures, and rank-zero
injection are due to the cited sources. The argument below extends
that injection to fixed ranks and supplies the quantitative coupling
and prefix constructions needed for the two limits. A bounded search
of the current article, OEIS entries, relevant literature, and the
available report catalogue found no identical later solution as of
3 October 2026. This is not a claim of exhaustive scholarly priority.
The constants and starting indices in asymptotic error terms are not
numerically certified.
\medskip

\section{Definitions and the two conjectures}\label{dgp:conj:sec:targets}
A partition will always be written in \emph{weakly increasing} order:
\[
 \lambda=(\lambda_1,\ldots,\lambda_k),\qquad
 1\leq\lambda_1\leq\cdots\leq\lambda_k,\qquad |\lambda|=n.
\]
Let $p(n)$ be the ordinary partition number and $q(n)$ the number of
partitions into distinct parts. For an integer $r\geq0$, put
\begin{align}
 s_r(n)&=\#\{\lambda\vdash n:\lambda_i\leq i+r\text{ for all }i\},\label{dgp:conj:eq:subdef}\\
 A(n)&=\#\{\lambda\vdash n:\lambda_i\geq i\text{ for all }i\}.\label{dgp:conj:eq:superdef}
\end{align}
The empty partition gives $s_r(0)=A(0)=p(0)=q(0)=1$.
The unshifted sequences are $s_0(n)=\mathrm{A238875}(n)$ and
$A(n)=\mathrm{A238873}(n)$ in OEIS~\cite{oeisA238875,oeisA238873}.
Changing the order convention changes the problem: weakly decreasing
superdiagonal partitions belong to the Rogers--Ramanujan class
A003114, not the class studied here.

Archibald et al.~\cite{abek}, Section~5, use exactly this increasing
convention. Their Conjecture~5.2 asserts $q(n)/A(n)\to0$; their
Conjecture~5.3 asserts $s_0(n)/p(n)\to1/2$. They prove
$A(n)/p(n)\to0$ in Proposition~5.1 and
$\limsup s_0(n)/p(n)\leq1/2$ in Proposition~5.5. The following
quantitative results settle the two conjectures.

\begin{theorem}[Uniform subdiagonal coupling]\label{dgp:conj:thm:coupling}
For a uniformly random partition of $n$, let $R$ be its largest part
minus its number of parts. As $n\to\infty$,
\begin{equation}\label{dgp:conj:eq:coupling}
 0\leq\PP_n(R\leq r)-\frac{s_r(n)}{p(n)}
 \leq\frac{C}{\log n}\qquad\text{for every integer }r\geq0,
\end{equation}
with one constant $C$ and one sufficiently large starting index,
independent of $r$. In particular,
\begin{align}
 \frac{s_0(n)}{p(n)}&=\frac12+O((\log n)^{-1}),\label{dgp:conj:eq:half}\\
 s_0(n)&=\frac{\exp(\pi\sqrt{2n/3})}{8\sqrt3\,n}
       \left(1+O((\log n)^{-1})\right).\label{dgp:conj:eq:subeq}
\end{align}
The one-sided upper error in \eqref{dgp:conj:eq:half} is $O(n^{-1/2})$.
\end{theorem}

\begin{theorem}[Exponential separation]\label{dgp:conj:thm:super}
Let
\begin{equation}\label{dgp:conj:eq:constants}
 b=\frac\pi{\sqrt3},\qquad
 B_*=\sqrt{\frac{\pi^2}{3}+2(\log2)^2}.
\end{equation}
Then
\begin{equation}\label{dgp:conj:eq:strong}
 \liminf_{n\to\infty}\frac{\log A(n)}{\sqrt n}\geq B_*.
\end{equation}
For every $\eta>0$, there is an index after which
\begin{equation}\label{dgp:conj:eq:ratio}
 \frac{q(n)}{A(n)}
 \leq \exp\bigl(-(B_*-b-\eta)\sqrt n\bigr).
\end{equation}
Since $B_*>b$, this proves $q(n)/A(n)\to0$.
\end{theorem}

Here and throughout, logarithms are natural. Both theorems are
asymptotic statements with unspecified constants and onset. The
finite inequalities used in the superdiagonal proof, however, hold
at every integer satisfying their explicitly stated feasibility
conditions. The remaining sections also distinguish the elementary
proof of Theorem~\ref{dgp:conj:thm:coupling} from the external rank theorem
used for the logistic corollary.

\begin{writenote}[6 October 2026]
The two conjectures are quoted verbatim from \cite{abek} in the front matter (``The two journal conjectures, and the OEIS entries''); Theorem~\ref{dgp:conj:thm:coupling} at $r=0$ is Conjecture~5.3 and Theorem~\ref{dgp:conj:thm:super} implies Conjecture~5.2. Later Parts sharpen both theorems, as follows. \emph{Part~II} proves Theorem~\ref{dgp:conj:thm:coupling}'s statement \eqref{dgp:conj:eq:half} again \emph{at $r=0$ only}, with the expansion $s_0(n)/p(n)=\frac12-\frac\beta8+O(\beta^2)$, $\beta=\pi/\sqrt{6n}$, to every order (Theorem~\ref{dgp:sub:thm:allorders}); so the error in \eqref{dgp:conj:eq:half} is $O(n^{-1/2})$ in both directions, \eqref{dgp:conj:eq:subeq} holds with relative error $O(n^{-1/2})$, and the ratio is below $\frac12$ for all large $n$ (by exact computation at the write, for every $10\le n\le400$; it equals $\frac12$ at $n=2$ and $n=9$ and exceeds it at $n=0,1,3,5$\textup{; added after the independent check of 6~October 2026: it is below $\frac12$ at $n=4,6,7,8$, so the list is complete for $n<10$, and a floating-point dynamic program extends ``below'' to every $10\le n\le2000$}). Part~II says nothing about $r\ge1$ or the logistic regime: for those, Theorem~\ref{dgp:conj:thm:coupling} and Corollary~\ref{dgp:conj:cor:logistic} are the only statements in this report. \emph{Part~III} replaces the lower rate $B_*$ of Theorem~\ref{dgp:conj:thm:super} by the exact rate: $\log A(n)/\sqrt n\to B=\sqrt{\pi^2/3+4\log^22}=2.2829\ldots>B_*=2.0617\ldots$ (Theorem~\ref{dgp:rate:thm:main}). Hence \eqref{dgp:conj:eq:ratio} holds with $B-b-\eta=0.4691\ldots-\eta$ in place of $B_*-b-\eta=0.2479\ldots-\eta$ (here $b=\pi/\sqrt3$). Part~IV adds the term $-Dn^{1/6}$, and Part~V, conditionally on an unrefereed preprint, the full equivalent. Theorem~\ref{dgp:conj:thm:super} and its proof stand as printed: they are the first proof of Conjecture~5.2, and the shortest.
\end{writenote}
\subsection{\wtag\ Choosable initial intervals}\label{dgp:conj:sub:choosable}
\begin{writenote}[6 October 2026]
This subsection is added in the write. It concerns the conjecture that OEIS A238873 displays (Gus Wiseman, 26~September 2025; quoted verbatim in the front matter), which no manuscript addresses: that $A(n)$ also counts the partitions $y$ of $n$ for which one can choose distinct positive integers $z_1,z_2,\dots$ with $z_i\le y_i$ for every $i$. The proof below is the write's; the same argument was posted earlier, on 31~August 2026, by Clément Garrot in OEIS A387112 \cite{oeisA387112}, with the remark that it proves the corresponding conjectures in A238873, A387118 and A388711.
\end{writenote}
\begin{proposition}\label{dgp:conj:prop:choosable}
Let $y$ be a partition of $n\ge0$, with parts $y_1,\dots,y_k$ listed in any order, and let $\lambda_1\le\cdots\le\lambda_k$ be its weakly increasing arrangement. There are distinct positive integers $z_1,\dots,z_k$ with $z_i\le y_i$ for every $i$ if and only if $\lambda_i\ge i$ for every $i$. Consequently, for every $n\ge0$, the number of partitions of $n$ with choosable initial intervals is $A(n)$ of \eqref{dgp:conj:eq:superdef}, and the number of those without is $p(n)-A(n)$.
\end{proposition}
\begin{proof}
Whether such $z_i$ exist depends only on the multiset of parts: if $\sigma$ permutes the indices, a choice $(z_i)$ for $(y_i)$ gives the choice $(z_{\sigma(i)})$ for $(y_{\sigma(i)})$. So we may assume $y_i=\lambda_i$.

If $\lambda_i\ge i$ for every $i$, take $z_i=i$: these are distinct, and $z_i=i\le\lambda_i$.

Conversely, let $z_1,\dots,z_k$ be distinct positive integers with $z_i\le\lambda_i$, and fix $j\in\{1,\dots,k\}$. For $i\le j$ we have $z_i\le\lambda_i\le\lambda_j$, so $z_1,\dots,z_j$ are $j$ distinct elements of $\{1,\dots,\lambda_j\}$, a set with $\lambda_j$ elements. Hence $j\le\lambda_j$.

For $n=0$ the empty partition satisfies both conditions vacuously (with the empty choice), in agreement with $A(0)=1$. Since the two classes coincide partition by partition, their counts agree for every $n$, and the complement has $p(n)-A(n)$ members.
\end{proof}
\begin{writenote}[6 October 2026]
Proposition~\ref{dgp:conj:prop:choosable} proves Wiseman's conjecture in A238873, and also the one in A387118 (Wiseman, 4~October 2025: the partitions without choosable initial intervals are counted by the non-superdiagonal weakly increasing partitions). Since the equivalence holds partition by partition and preserves the number of parts, it also proves the refinement by number of parts conjectured in A388711 (Wiseman's entry of 23~September 2025) and the prime-index form conjectured in A387112 \cite{oeisA387112} (the numbers whose weakly increasing prime indices $y$ satisfy $y_i\ge i$), as Garrot's comment says \textup{(sentence added after the independent check of 6~October 2026)}. Exhaustive search at the write, with an independent backtracking search for the distinct choices, confirms the equivalence partition by partition and the counts $1,1,1,2,3,3,5,7,\dots,289$ of A238873 for $0\le n\le24$. Nothing was submitted to the OEIS.
\end{writenote}

\section{Classical input and a fixed rank injection}\label{dgp:conj:sec:rank}
Put $c=\pi/\sqrt6$, so $2c=\pi\sqrt{2/3}$. We use the classical
refined Hardy--Ramanujan estimate~\cite{dlmfpart}, obtainable from
the first term of the Rademacher expansion there,
\begin{equation}\label{dgp:conj:eq:hr}
 p(n)=\frac{e^{2c\sqrt n}}{4\sqrt3\,n}
       \left(1+O(n^{-1/2})\right).
\end{equation}
For distinct partitions we need only the classical leading
formula, also recalled in \cite[Section~5.2]{abek},
\begin{equation}\label{dgp:conj:eq:distinct}
 q(n)\sim\frac{e^{b\sqrt n}}{4\,3^{1/4}n^{3/4}},
 \qquad \log q(n)=b\sqrt n+O(\log n).
\end{equation}
These are inputs; no asymptotic enumeration of the constrained
classes is being assumed.

For $n\geq1$, write $L=\lambda_k$, $k=\ell(\lambda)$,
$Z_1=\#\{i:\lambda_i=1\}$, and $R=L-k$. Let
\[
 N(r,n)=\#\{\lambda\vdash n:R(\lambda)=r\},\qquad
 D(n)=p(n)-p(n-1).
\]
Deleting one part equal to $1$ shows that $D(n)$ counts partitions
of $n$ with no ones.

\begin{lemma}[Uniform rank atom bound]\label{dgp:conj:lem:atom}
For every $n\geq2$ and every integer $r$,
\begin{equation}\label{dgp:conj:eq:atom}
 N(r,n)\leq D(n),\qquad \frac{D(n)}{p(n)}=O(n^{-1/2}).
\end{equation}
Consequently the probability of any interval of $m$ consecutive
integer ranks is at most $mD(n)/p(n)$.
\end{lemma}
\begin{proof}
Fix the source rank $r$. For a partition other than the all-ones
partition, delete its $j=Z_1$ ones and add $j$ to one largest remaining
part. For $j=0$ this is the identity. The resulting partition $\nu$
has no ones, has the same weight, and satisfies
\begin{equation}\label{dgp:conj:eq:rankshift}
 R(\nu)=(L+j)-(k-j)=r+2j.
\end{equation}
Thus $\nu$ and the specified rank $r$ determine
$j=(R(\nu)-r)/2$. If $j>0$, the enlarged part is the unique largest
part, even if several parts were largest before the map. Subtracting
$j$ from this part and restoring $j$ ones recovers the source. If
$j=0$, the inverse is the identity. Outputs coming from these two
cases cannot coincide at fixed $r$, by \eqref{dgp:conj:eq:rankshift}.
This proves the injection into the $D(n)$ no-ones partitions.

The excluded partition has rank $1-n$. This rank is attained only
by the all-ones partition: $L\geq1$ and $k\leq n$, with equality in
$L-k\geq1-n$ only in that case. Since the singleton $(n)$ has no ones
for $n\geq2$, $D(n)\geq1$ and \eqref{dgp:conj:eq:atom} holds at this rank too.
Empty rank classes cause no difficulty.

Finally, directly taking the ratio in \eqref{dgp:conj:eq:hr} gives
\[
 \frac{p(n-1)}{p(n)}
 =\frac{n}{n-1}e^{2c(\sqrt{n-1}-\sqrt n)}
    \left(1+O(n^{-1/2})\right)
 =1+O(n^{-1/2}).
\]
Together with $D(n)\geq0$, this proves the stated bound. No sharper
one-step error is needed. Summing the atom bound gives the interval
assertion.
\end{proof}

\noindent\textbf{Attribution.}
The rank-zero instance of this map is the injection in
\cite[Lemma~5.4]{abek}. Its extension to each specified source rank,
using \eqref{dgp:conj:eq:rankshift} to recover the number of deleted ones,
is the step needed here.
\begin{writenote}[6 October 2026]
At rank zero the bound reads $N(0,n)\le D(n)=O(n^{-1/2}p(n))$; the true size is $N(0,n)/p(n)=\frac\beta4+O(\beta^2)$ (Liu--Zhou, used in Part~II at \eqref{dgp:sub:eq:rankcubic}), so the order $n^{-1/2}$ of \eqref{dgp:conj:eq:atom} is attained at $r=0$.
\end{writenote}

Ferrers conjugation exchanges the largest part and the number of
parts, so $N(r,n)=N(-r,n)$. Therefore
\begin{equation}\label{dgp:conj:eq:rankhalf}
 \PP_n(R\leq0)=\frac12+\frac{N(0,n)}{2p(n)}
            =\frac12+O(n^{-1/2}).
\end{equation}
A subdiagonal partition necessarily has $L\leq k$, explaining the
source paper's upper bound. The task is to show that the interior
indices exclude only a negligible additional proportion.

\section{Interior control and exceptional events}\label{dgp:conj:sec:interior}
\begin{lemma}[Three regions]\label{dgp:conj:lem:interior}
Let $r\geq0$ and $m\geq1$ be integers. Suppose that a partition
$\lambda\vdash n$ has $k$ parts and satisfies
\begin{equation}\label{dgp:conj:eq:interiorhyp}
 k\geq2m+1,\quad Z_1\geq m,\quad L\leq k+r-m,
 \quad (m+1)(k-m+1)>n.
\end{equation}
Then $\lambda_i\leq i+r$ for every index $i$.
\end{lemma}
\begin{proof}
For $i\leq m$, the part is $1$. For $i>k-m$, the largest-part
bound gives $\lambda_i\leq k+r-m\leq i+r$. If a remaining index
$m<i\leq k-m$ violates the conclusion, integrality and $r\geq0$
imply $\lambda_i\geq i+r+1\geq i+1$. The final $k-i+1$ parts
then give
\[
 n\geq(i+1)(k-i+1).
\]
This is a concave quadratic in $i$, so its minimum on
$[m+1,k-m]$ occurs at an endpoint. Those endpoint values are
$(m+2)(k-m)$ and $(m+1)(k-m+1)$; their difference is
$k-2m-1\geq0$. Both are at least the last expression in
\eqref{dgp:conj:eq:interiorhyp}, a contradiction.
\end{proof}
\begin{writenote}[6 October 2026]
Part~II restates this interior step, with the same cutoffs \eqref{dgp:conj:eq:scales} and the same short-length bound \eqref{dgp:conj:eq:length}, in its Section~\ref{dgp:sub:sec:localize} (\eqref{dgp:sub:eq:cutoffs}, \eqref{dgp:sub:eq:short}). The restatement is kept there, because Part~II's endpoint analysis continues from it.
\end{writenote}

For all sufficiently large $n$, choose
\begin{equation}\label{dgp:conj:eq:scales}
 K=\floor{\frac{\sqrt n\log n}{4c}},\qquad
 m=\ceil{\frac{5c\sqrt n}{\log n}}.
\end{equation}
Then $K/m\to\infty$, $m=o(\sqrt n)$, and
\begin{equation}\label{dgp:conj:eq:areascale}
 \frac{(m+1)(K-m+1)}{n}\longrightarrow\frac54.
\end{equation}
Thus $K\geq2m+1$ and $(m+1)(K-m+1)>n$ eventually. Both
conditions remain true with $k\geq K$ replacing $K$.
The constants in \eqref{dgp:conj:eq:scales} are convenient for the proof;
no all-$n$ feasibility assertion is made.

\subsection{Too few ones}
Deleting exactly $m$ ones is a bijection between partitions of $n$
with at least $m$ ones and all partitions of $n-m$. Hence
\begin{equation}\label{dgp:conj:eq:ones}
 \PP_n(Z_1<m)=1-\frac{p(n-m)}{p(n)}=O((\log n)^{-1}).
\end{equation}
To justify the variable shift directly, $n-m\sim n$ and
\eqref{dgp:conj:eq:hr} give
\begin{align*}
 \log\frac{p(n-m)}{p(n)}
 &=2c(\sqrt{n-m}-\sqrt n)+\log\frac n{n-m}+O(n^{-1/2})\\
 &=O(m/\sqrt n)=O((\log n)^{-1}).
\end{align*}
Exponentiation and $p(n-m)\leq p(n)$ prove \eqref{dgp:conj:eq:ones}.
This does not rely on subtracting two uncontrolled asymptotic
remainders or on a sharpened one-step ratio.

\subsection{Too few parts}
By conjugation, $\PP_n(k<K)$ is the probability that the largest part
is less than $K$. Let $p_{\leq K}(n)$ count partitions whose parts
are at most $K$. With $t=c/\sqrt n$ and $z=e^{-t}$, coefficient
positivity in the ordinary partition product gives
\begin{equation}\label{dgp:conj:eq:productbound}
 p_{\leq K}(n)\leq z^{-n}P(z)\prod_{j>K}(1-z^j),
 \qquad P(z)=\prod_{j\geq1}(1-z^j)^{-1}.
\end{equation}
The function $f(x)=-\log(1-e^{-tx})$ is positive and decreasing;
its logarithmic singularity at zero is integrable. Thus
\begin{align*}
 \log P(e^{-t})&=\sum_{j\geq1}f(j)
 \leq\int_0^\infty f(x)\dd x\\
 &=\sum_{\ell\geq1}\frac1\ell\int_0^\infty e^{-t\ell x}\dd x
 =\frac{\pi^2}{6t}.
\end{align*}
The exchange uses nonnegative terms. In particular,
\[
 z^{-n}P(z)\leq e^{tn+\pi^2/(6t)}=e^{2c\sqrt n},
 \qquad \frac{z^{-n}P(z)}{p(n)}=O(n).
\]
Furthermore, $\log(1-u)\leq-u$ implies
\begin{equation}\label{dgp:conj:eq:tailproduct}
 \prod_{j>K}(1-z^j)
 \leq\exp\left(-\frac{z^{K+1}}{1-z}\right),\qquad
 \frac{z^{K+1}}{1-z}\sim c^{-1}n^{1/4}.
\end{equation}
Indeed, the floor in \eqref{dgp:conj:eq:scales} gives
$z^{K+1}=n^{-1/4}(1+O(n^{-1/2}))$, while
$1-z\sim c/\sqrt n$. For a fixed $c_0>0$ small enough,
\begin{equation}\label{dgp:conj:eq:length}
 \PP_n(k<K)\leq\frac{p_{\leq K}(n)}{p(n)}
 =O\bigl(n e^{-c_0n^{1/4}}\bigr).
\end{equation}

\subsection{The rank window}
For every integer $r$, the interval $r-m<R\leq r$ contains exactly
$m$ integer ranks. Lemma~\ref{dgp:conj:lem:atom} therefore gives
\begin{equation}\label{dgp:conj:eq:window}
 \PP_n(r-m<R\leq r)\leq\frac{mD(n)}{p(n)}
 =O((\log n)^{-1}),
\end{equation}
uniformly in $r$. The treatment of the all-ones partition in that
lemma avoids any exceptional-rank convention in this bound.

\subsection{Completion of the first conjecture}
Every $r$-subdiagonal partition satisfies $L\leq k+r$, hence
$R\leq r$. Conversely, among partitions satisfying $R\leq r$,
remove the events $Z_1<m$, $k<K$, and $r-m<R\leq r$.
All remaining partitions satisfy Lemma~\ref{dgp:conj:lem:interior}.
For every sufficiently large $n$ and every integer $r\geq0$, we
obtain the exact union-bound inequality
\begin{align}
 0\leq\PP_n(R\leq r)-\frac{s_r(n)}{p(n)}
 &\leq \PP_n(Z_1<m)+\PP_n(k<K)\nonumber\\
 &\hspace{1em}+\PP_n(r-m<R\leq r).\label{dgp:conj:eq:union}
\end{align}
No independence among these events is assumed. Estimates
\eqref{dgp:conj:eq:ones}, \eqref{dgp:conj:eq:length}, and \eqref{dgp:conj:eq:window} prove the
uniform coupling. Setting $r=0$ and using \eqref{dgp:conj:eq:rankhalf} gives
\eqref{dgp:conj:eq:half}, including its sharper one-sided upper bound.
Finally \eqref{dgp:conj:eq:hr} gives \eqref{dgp:conj:eq:subeq}. This completes the
proof of Theorem~\ref{dgp:conj:thm:coupling}, and thus of Conjecture~5.3.

\section{A separate logistic crossover}\label{dgp:conj:sec:logistic}
The preceding proof uses no rank limit theorem. The following
stronger description of growing shifts invokes the positive-rank
part of Dousse--Mertens~\cite[Theorem~1.2]{dm}.

\begin{corollary}\label{dgp:conj:cor:logistic}
Uniformly for all integers $r\geq0$,
\begin{equation}\label{dgp:conj:eq:logistic}
 \frac{s_r(n)}{p(n)}
 =\frac1{1+\exp(-\pi r/\sqrt{6n})}+O((\log n)^{-1}).
\end{equation}
Consequently, if $r_n/\sqrt n\to\alpha\in[0,\infty)$, the ratio
tends to $(1+e^{-c\alpha})^{-1}$. If $r_n/\sqrt n\to\infty$,
the ratio tends to $1$.
\end{corollary}
\begin{proof}
Set $\beta=c/\sqrt n$, $J=\floor{\sqrt n\log n/(\pi\sqrt6)}$,
$g(x)=\tfrac14\sech^2(x/2)$, and $F(x)=(1+e^{-x})^{-1}$.
The cited theorem yields, uniformly for $1\leq j\leq J$,
\begin{equation}\label{dgp:conj:eq:dm}
 \frac{N(j,n)}{p(n)}
 =\beta g(\beta j)\left(1+O(\beta^{1/2}j^{1/3})\right).
\end{equation}
Only positive ranks are used here; the zero rank is controlled
by Lemma~\ref{dgp:conj:lem:atom}. Since $g$ decreases on $[0,\infty)$,
\begin{equation}\label{dgp:conj:eq:riemann}
 \left|\beta\sum_{j=1}^r g(\beta j)
      -\int_0^{\beta r}g(x)\dd x\right|\leq\beta g(0)
 \qquad(0\leq r\leq J).
\end{equation}
In particular the main-term masses have bounded sum. Summing the
relative error in \eqref{dgp:conj:eq:dm} thus costs at most
$O(\beta^{1/2}J^{1/3})$, not $J$ times this quantity. By
\eqref{dgp:conj:eq:rankhalf}, \eqref{dgp:conj:eq:riemann}, and $F'=g$,
\[
 \PP_n(R\leq r)=F(\beta r)
 +O\bigl(n^{-1/12}(\log n)^{1/3}\bigr)
 \qquad(0\leq r\leq J).
\]
At the endpoint $\beta J=(\log n)/6+O(n^{-1/2})$, so
$1-F(\beta J)=O(n^{-1/6})$. For $r>J$, both distribution
functions lie between their values at $J$ and $1$; monotonicity
therefore extends the same error bound to all such $r$.
This error is $o(1/\log n)$. Combining it with
Theorem~\ref{dgp:conj:thm:coupling} proves \eqref{dgp:conj:eq:logistic}.
\end{proof}
\begin{writenote}[6 October 2026]
The write read the statement of the cited theorem \cite{dm} in arXiv:1406.6848: for $|m|\le\sqrt n\log n/(\pi\sqrt6)$, $N(m,n)=\frac\beta4\operatorname{sech}^2(\frac{\beta m}2)\,p(n)\bigl(1+O(\beta^{1/2}|m|^{1/3})\bigr)$ with $\beta=\pi/\sqrt{6n}$. With $g(x)=\frac14\operatorname{sech}^2(x/2)$ this is \eqref{dgp:conj:eq:dm} for $1\le j\le J$. The theorem's range includes $m=0$, where its displayed relative error vanishes; the proof above does not use it there, but takes the rank-zero term from Lemma~\ref{dgp:conj:lem:atom}. \textup{(Added after the independent check of 6~October 2026.)} Read literally at $m=0$, the statement would give $N(0,n)=\frac\beta4p(n)$ exactly for all large $n$, which is false: by Part~II's \eqref{dgp:sub:eq:rankcubic}, $N(0,n)/p(n)=\frac\beta4+a\beta^2+O(\beta^3)$ with $a=\frac3{16}-\frac3{2\pi^2}=0.0355\ldots\ne0$, and in exact arithmetic $N(0,2000)/(\beta p(2000))=0.25103$. (Their Remark~1.3 obtains the factor $|m|^{1/3}$ from the choice $\alpha(m)=|m|^{-1/3}$, which is undefined at $m=0$.) The rank-zero case of \cite{dm} is used nowhere in this report. Corollary~\ref{dgp:conj:cor:logistic} is the only statement of the report that depends on \cite{dm}.
\end{writenote}

\section{The inverse at the precision actually proved}\label{dgp:conj:sec:inverse}
For real $y>1$, define the first threshold
\begin{equation}\label{dgp:conj:eq:threshold}
 T(y)=\min\{n\geq0:s_0(n)\geq y\}.
\end{equation}
Prepending a part $1$ is an injection from subdiagonal partitions of
$n$ to those of $n+1$: each old index increases by one. Thus $s_0(n)$
is nondecreasing; \eqref{dgp:conj:eq:subeq} makes it unbounded. It is in fact
strictly increasing for $n\geq4$. At every odd target size at least
$5$, a partition $(1,2,\ldots,2)$ has exactly one one; at every even
target size at least $6$, a partition $(1,2,\ldots,2,3)$ does too.
Both are subdiagonal and outside the prepending map's image, whose
members have at least two ones when the source is nonempty.

\begin{theorem}[Limited-precision inverse]\label{dgp:conj:thm:inverse}
Let $W_{-1}$ be the real branch of Lambert's function with values at
most $-1$. For sufficiently large $y$, put
\begin{equation}\label{dgp:conj:eq:inversecenter}
 x_0(y)=c^{-2}
 W_{-1}\!\left(-\frac{c}{\sqrt{8\sqrt3\,y}}\right)^2.
\end{equation}
Then
\begin{align}
 T(y)&=x_0(y)+O\left(\frac{\sqrt{x_0(y)}}{\log x_0(y)}+1\right)\label{dgp:conj:eq:inverseerror}\\
 &=x_0(y)+O\left(\frac{\log y}{\log\log y}\right).\nonumber
\end{align}
Writing $H=\log y$, an equivalent explicit initializer is
\begin{equation}\label{dgp:conj:eq:loginverse}
 T(y)=\frac{\left[H+2\log H+\log(12\sqrt3/\pi^2)\right]^2}{4c^2}
       +O\left(\frac{H}{\log H}\right).
\end{equation}
\end{theorem}
\begin{proof}
For $x>0$ define
\[
 M(x)=\frac{e^{2c\sqrt x}}{8\sqrt3\,x},\qquad
 h(x)=\log M(x)=2c\sqrt x-\log x-\log(8\sqrt3).
\]
The large solution of $M(x)=y$ is \eqref{dgp:conj:eq:inversecenter}:
put $u=\sqrt x$, take the positive square root, and rearrange as
$(-cu)e^{-cu}=-c/\sqrt{8\sqrt3\,y}$. The branch $W_{-1}$
selects $u\to\infty$. Also
\begin{equation}\label{dgp:conj:eq:logmodel}
 \log s_0(n)=h(n)+O((\log n)^{-1}),\qquad
 h'(x)=\frac c{\sqrt x}-\frac1x\sim\frac c{\sqrt x}.
\end{equation}
Let $x_0=x_0(y)$ and, for a fixed sufficiently large $D$, put
$\Delta=D\sqrt{x_0}/\log x_0$. For
$x\in[x_0-\Delta-1,x_0+\Delta+1]$, one has $x\sim x_0$,
$\log x\sim\log x_0$, and $h'(x)\sim c/\sqrt{x_0}$ uniformly.
Thus
\[
 h(x_0\pm\Delta)-h(x_0)
 =\pm\frac{cD+o(1)}{\log x_0}.
\]
Choose $D$ large enough to dominate the fixed implicit constant
in \eqref{dgp:conj:eq:logmodel}. Rounding the lower point down and the upper
point up changes $h$ by only $O(x_0^{-1/2})=o(1/\log x_0)$.
For all sufficiently large $y$, therefore,
\[
 s_0(\floor{x_0-\Delta})<y,
 \qquad s_0(\ceil{x_0+\Delta})\geq y.
\]
Monotonicity brackets the first threshold and proves
\eqref{dgp:conj:eq:inverseerror}. Since $\sqrt{x_0}\sim H/(2c)$ and
$\log x_0\sim2\log H$, the second error form follows.

To obtain \eqref{dgp:conj:eq:loginverse}, set $v=2c\sqrt{x_0}$.
The model equation becomes
\[
 v-2\log v=H+\log(2\sqrt3/c^2).
\]
It follows first that $v\sim H$, and substitution back into this
equation gives
\[
 v=H+2\log H+\log(2\sqrt3/c^2)+O(\log H/H).
\]
Squaring changes $x_0$ by $O(\log H)$, which is smaller than
$O(H/\log H)$, and $2\sqrt3/c^2=12\sqrt3/\pi^2$.
\end{proof}

The window in \eqref{dgp:conj:eq:inverseerror} grows. This theorem does not
identify the exact integer threshold by rounding $x_0$, does not
produce two adjacent candidate integers, and does not supply a
numerically certified interval for any specified $y$. Such statements
would need explicit error constants and substantially sharper
information. The elementary nested-log initializer is not an
all-orders expansion for the constrained counting sequence.
\begin{writenote}[6 October 2026]
Part~II's Theorem~\ref{dgp:sub:thm:inverse} sharpens Theorem~\ref{dgp:conj:thm:inverse}, which concerns only $r=0$: with the same centre $x_0$, the threshold lies between $\lceil X_J-C_Jx_0^{-J/2}\rceil$ and $\lceil X_J+C_Jx_0^{-J/2}\rceil$ for every fixed $J$, with $X_1=x_0+\frac7{24}+\frac3{\pi^2}$. The centre $x_0$ is an instance of the transseries volume's Lambert core after the reduction $\xi=\sqrt x$ (front matter, ``The inverses''). Strict monotonicity of $s_0$ from $n=4$ on was confirmed by exact computation for $4\le n\le400$ (and $s_0(3)=s_0(4)=2$, so it does not hold from $n=3$). The injection that prepends a part~$1$ is used again in Part~II, Section~\ref{dgp:sub:sec:inverse}.
\end{writenote}

\section{Catalan prefixes and distinct tails}\label{dgp:conj:sec:prefix}
We now prove Conjecture~5.2 independently of the subdiagonal
argument. The first construction already yields a strict
exponential separation. A refinement in the next section improves
the constant.

For $M\geq1$, let $\mathcal C_M$ consist of weakly increasing
lists $\lambda=(\lambda_1,\ldots,\lambda_M)$ satisfying
\begin{equation}\label{dgp:conj:eq:prefixdef}
 i\leq\lambda_i\leq M\qquad(1\leq i\leq M).
\end{equation}
Reverse complementation,
$\mu_j=M+1-\lambda_{M+1-j}$, is a bijection from
$\mathcal C_M$ to weakly increasing lists with $1\leq\mu_j\leq j$.
Placing the $j$th north step at horizontal coordinate $\mu_j-1$
and then completing with east steps gives the usual Dyck path from
$(0,0)$ to $(M,M)$. The inverse reads those coordinates. The
reflection principle therefore gives
\begin{equation}\label{dgp:conj:eq:catalan}
 |\mathcal C_M|=C_M=\binom{2M}{M}-\binom{2M}{M+1}
                  =\frac1{M+1}\binom{2M}{M}.
\end{equation}
Every prefix has weight $w(\lambda)\leq M^2$.

Let $f_M(N)$ count distinct partitions of $N$ with all parts greater
than $M$, with $f_M(0)=1$ and $f_M(N)=0$ for negative $N$.
For a prefix $\lambda\in\mathcal C_M$, append such a distinct tail
$\tau_1<\cdots<\tau_\ell$. Its $j$th part is at least $M+j$,
so the combined list is superdiagonal at index $M+j$ as well as at
every prefix index. It is weakly increasing, since the prefix parts
are at most $M$ and the tail parts exceed $M$. The cutoff $M$
recovers the prefix and tail uniquely. Consequently
\begin{equation}\label{dgp:conj:eq:exactprefix}
 A(n)\geq\sum_{\lambda\in\mathcal C_M}f_M(n-w(\lambda)).
\end{equation}

\begin{lemma}[Tail comparison]\label{dgp:conj:lem:tail}
For $N\geq M+1$, the values $f_M(N)$ are positive and nondecreasing,
and
\begin{equation}\label{dgp:conj:eq:tailcompare}
 2^M f_M(N)\geq q(N).
\end{equation}
\end{lemma}
\begin{proof}
The singleton $(N)$ shows positivity. Increasing the unique largest
part of a distinct partition by $1$ preserves distinctness and the
lower cutoff. It is injective: on its image, subtract $1$ from the
still unique largest part. Thus $f_M(N+1)\geq f_M(N)$.

Every distinct partition is uniquely determined by its subset of
small parts and its remaining large parts. Hence
\begin{equation}\label{dgp:conj:eq:subsetdecomp}
 q(N)=\sum_{S\subseteq\{1,\ldots,M\}}
 f_M\left(N-\sum_{j\in S}j\right).
\end{equation}
Each summand is at most $f_M(N)$. For arguments at least $M+1$,
use monotonicity; arguments between $1$ and $M$ contribute zero;
argument zero contributes $1\leq f_M(N)$; negative arguments
contribute zero. There are $2^M$ subsets.
\end{proof}

\begin{proposition}[Finite Catalan lower bound]\label{dgp:conj:prop:finite}
For integers $M\geq1$ and $n-M^2\geq M+1$,
\begin{equation}\label{dgp:conj:eq:finitebase}
 A(n)\geq C_M f_M(n-M^2)\geq C_M2^{-M}q(n-M^2).
\end{equation}
\end{proposition}
\begin{proof}
Every argument in \eqref{dgp:conj:eq:exactprefix} is at least $n-M^2$,
inside the monotonicity range. Apply Lemma~\ref{dgp:conj:lem:tail}.
\end{proof}

Fix $0<\varepsilon<1$ and take $M=\floor{\varepsilon\sqrt n}$.
The feasibility condition in Proposition~\ref{dgp:conj:prop:finite} holds
for all sufficiently large $n$. Stirling's formula in
\eqref{dgp:conj:eq:catalan} and \eqref{dgp:conj:eq:distinct} give
\[
 \log(C_M2^{-M})=M\log2+O(\log M),\qquad
 \log q(n-M^2)=b\sqrt{n-M^2}+O(\log n).
\]
It follows that
\begin{equation}\label{dgp:conj:eq:baselineepsilon}
 \liminf_{n\to\infty}\frac{\log A(n)}{\sqrt n}
 \geq\varepsilon\log2+b\sqrt{1-\varepsilon^2}.
\end{equation}
The maximum of the right side is
\begin{equation}\label{dgp:conj:eq:baseline}
 B_0=\sqrt{b^2+(\log2)^2}>b,
 \qquad\varepsilon=\frac{\log2}{B_0}.
\end{equation}
This follows from Cauchy--Schwarz applied to the unit vector
$(\varepsilon,\sqrt{1-\varepsilon^2})$, with equality at the
stated value. Thus the simple construction alone proves
Conjecture~5.2. It makes no assertion that $B_0$ is the actual
growth rate.

\section{Fixed height prefixes improve the exponential rate}\label{dgp:conj:sec:height}
Fix an integer $h\geq3$. Restrict the prefixes further to
\begin{equation}\label{dgp:conj:eq:heightprefix}
 i\leq\lambda_i\leq\min(M,i+h-1).
\end{equation}
Under the reverse complementation in Section~\ref{dgp:conj:sec:prefix},
this becomes $1\leq\mu_j\leq j$ and $\mu_j\geq j-h+1$.
In the associated Dyck path, the height $y-x$ just after its $j$th
north step is $j-\mu_j+1$. Maxima occur just after north steps,
so \eqref{dgp:conj:eq:heightprefix} corresponds exactly to Dyck paths of
semilength $M$ with maximum height at most $h$.

Let $D_h(M)$ denote their number. If $T_h$ is the adjacency matrix
of the path graph on vertices $0,1,\ldots,h$, recording successive
heights gives
\begin{equation}\label{dgp:conj:eq:walk}
 D_h(M)=(T_h^{2M})_{00}.
\end{equation}
The standard sine eigenvectors can also be checked directly:
for $1\leq j\leq h+1$ and $0\leq a\leq h$,
\[
 v_j(a)=\sqrt{\frac2{h+2}}\sin\frac{(a+1)j\pi}{h+2},\qquad
 T_hv_j=2\cos\frac{j\pi}{h+2}\,v_j.
\]
The boundary values at $a=-1,h+1$ vanish, and the sine addition
formula verifies the eigenvalue equation. The vectors are an
orthonormal basis by finite sine orthogonality. Therefore
\begin{equation}\label{dgp:conj:eq:spectral}
 D_h(M)=\frac2{h+2}\sum_{j=1}^{h+1}
 \sin^2\frac{j\pi}{h+2}
 \left(2\cos\frac{j\pi}{h+2}\right)^{2M}.
\end{equation}
The eigenvalues of largest modulus are opposites, corresponding to
$j=1$ and $j=h+1$. Since $2M$ is even, their contributions add.
For each \emph{fixed} $h$ this proves
\begin{equation}\label{dgp:conj:eq:heightasymp}
 D_h(M)=c_h\rho_h^M(1+O_h(\theta_h^M)),\qquad
 \rho_h=4\cos^2\frac\pi{h+2},\quad
 c_h=\frac4{h+2}\sin^2\frac\pi{h+2}>0,
\end{equation}
where $0\leq\theta_h<1$. In particular $\rho_h>2$ for $h\geq3$.
No uniformity in $h$ is asserted in this estimate.

Every restricted prefix has weight at most
\begin{equation}\label{dgp:conj:eq:weight}
 W_h(M)=\frac{M(M+1)}2+(h-1)M.
\end{equation}
This bound is deliberately simple; the upper cutoff $M$ can make
it nonsharp at the final indices. The same injection and tail
comparison now give the finite inequality
\begin{equation}\label{dgp:conj:eq:finiteheight}
 A(n)\geq D_h(M)2^{-M}q(n-W_h(M))
 \quad\text{if }n-W_h(M)\geq M+1.
\end{equation}
Here $h,M,n$ are integers, and the statement itself is exact.

Keep $h$ fixed and put $M=\floor{\varepsilon\sqrt n}$, with
$0<\varepsilon<\sqrt2$. Then $W_h(M)/n\to\varepsilon^2/2$, so
the feasibility condition eventually holds. Using
\eqref{dgp:conj:eq:heightasymp}, \eqref{dgp:conj:eq:distinct}, and
\eqref{dgp:conj:eq:finiteheight}, we obtain
\begin{equation}\label{dgp:conj:eq:heightepsilon}
 \liminf_{n\to\infty}\frac{\log A(n)}{\sqrt n}
 \geq\varepsilon\log(\rho_h/2)
       +b\sqrt{1-\varepsilon^2/2}.
\end{equation}
Write $a_h=\log(\rho_h/2)>0$. Cauchy--Schwarz, now with the
unit vector $(\varepsilon/\sqrt2,\sqrt{1-\varepsilon^2/2})$,
shows that the maximum is
\begin{equation}\label{dgp:conj:eq:Bh}
 B_h=\sqrt{b^2+2a_h^2},\qquad
 \varepsilon_h=\frac{2a_h}{\sqrt{b^2+2a_h^2}}\in(0,\sqrt2).
\end{equation}
Thus the left side of \eqref{dgp:conj:eq:heightepsilon} is at least $B_h$
for \emph{every fixed} integer $h\geq3$. Taking the supremum of
these already established lower bounds and using $\rho_h\to4$
gives \eqref{dgp:conj:eq:strong}.

This order of limits is essential: first fix $h$, then take
$n\to\infty$, then take the supremum over fixed heights. No
estimate with a height growing with $n$ is used or needed.
By \eqref{dgp:conj:eq:distinct},
\[
 \limsup_{n\to\infty}\frac1{\sqrt n}\log\frac{q(n)}{A(n)}
 \leq b-B_*.
\]
This is equivalent to \eqref{dgp:conj:eq:ratio} and completes the proof of
Theorem~\ref{dgp:conj:thm:super}.

For orientation, the three constants are approximately
\[
 b=1.8137993642,\qquad B_0=1.9417314819,\qquad
 B_*=2.0617405660.
\]
These decimals merely illustrate exact expressions. The trivial
upper bound $A(n)\leq p(n)$ yields
\begin{equation}\label{dgp:conj:eq:unknownrate}
 B_*\leq\liminf\frac{\log A(n)}{\sqrt n}
 \leq\limsup\frac{\log A(n)}{\sqrt n}
 \leq\pi\sqrt{2/3}=2.5650996603\ldots.
\end{equation}
Neither existence of a limit nor equality with $B_*$ follows from
these inequalities.
\begin{writenote}[6 October 2026]
\emph{Settled by Part~III.} The limit exists and equals $B=\sqrt{\pi^2/3+4\log^22}=2.2829104645\ldots$, strictly between $B_*$ and $\pi\sqrt{2/3}$ (Theorem~\ref{dgp:rate:thm:main}). Part~III removes both losses that Part~I's ``further questions'' name: an exponential tilt gives a matching upper bound, and a geometric tail with a positive survival probability replaces the distinct tail and its comparison $2^Mf_M(N)\ge q(N)$. The bounded-height Dyck prefix of this section is Part~III's Lemma~\ref{dgp:rate:lem:dyck}, used there with a height growing like $t^{-1/3}$ through the exact spectral bound \eqref{dgp:rate:eq:spectralbound}, which is the sum \eqref{dgp:conj:eq:spectral} with all but two terms dropped.
\end{writenote}

\section{Exact computation and reproducibility}\label{dgp:conj:sec:computation}
The source package supplies a Python standard-library companion.
Its exact integer computations test the finite mechanisms, not the
truth of an asymptotic statement by numerical extrapolation.
The empty-partition convention and increasing order are included
in all sequence checks.

\subsection{Independent enumeration}
One method directly enumerates increasing partitions and tests the
index inequalities. A second method processes part sizes and their
multiplicities, recording the number of parts and total weight.
For a block of $a$ new parts equal to $v$ following $k$ smaller
parts, the subdiagonal condition is $v\leq k+1+r$ when $a>0$;
the superdiagonal condition is $v\geq k+a$. These are respectively
the strongest inequalities at the first and last index of the
block. Thus the multiplicity method enforces the definitions
without reusing the direct enumerator.

The original OEIS numeric prefixes have 57 entries for A238875
and 61 for A238873, both beginning at $n=0$. The companion stores
the extracted integers and source URLs, and compares them with
its independently computed values. Exact data extend through
$n=200$. The following sample values are intended only as
illustrations of the finite sequences.
\begin{center}
\begin{tabular}{r r r r r}
\toprule
$n$ & $p(n)$ & $s_0(n)$ & $A(n)$ & $q(n)$\\
\midrule
25 & 1958 & 898 & 349 & 142\\
40 & 37338 & 17319 & 3997 & 1113\\
50 & 204226 & 95502 & 16582 & 3658\\
100 & 190569292 & 91171308 & 5323937 & 444793\\
200 & 3972999029388 & 1930845918142 & 25548856857 & 487067746\\
\bottomrule
\end{tabular}
\end{center}
For example, at $n=200$ the ratios are about $0.485992$ and
$0.0190642$. These do not establish an error constant, a valid
starting index, or the exact superdiagonal rate.

\subsection{Finite mechanisms checked}
The companion verifies the fixed-rank map and its inverse, detects
collisions within each rank class, and treats the all-ones rank
separately. It checks the deterministic interior lemma whenever its
hypotheses hold; shifted subdiagonal counts are compared with the
appropriate rank cumulative counts. The prefix tests compare
Catalan enumeration with the exact binomial formula, and bounded
height prefix enumeration with integer walk counts. Distinct-tail
monotonicity, the small-subset decomposition, prefix--tail validity
and disjointness, and both finite lower bounds are tested with
integer arithmetic. The shipped verification exhausts 215,308 partitions, including the
empty partition, through $n=40$; checks 215,267 injection and inverse
cases and 446,242 interior-lemma cases; and verifies 1,690 Catalan and
9,995 bounded-height finite lower bounds. The reproducible output
records the remaining ranges and numbers of successful checks.

Finite verification does not replace any proof above. In
particular, the asymptotic choices \eqref{dgp:conj:eq:scales} need not satisfy
the interior hypotheses at modest $n$, and finite data cannot
certify the eventual spectral or partition asymptotics. The exact
threshold command, where supported within a finite table, must
search actual integer counts; the asymptotic center is not an
integer-threshold certificate.

\subsection{Package contract}
The accompanying README gives bounded commands, exact result
files, and regression tests that run both normally and with Python
optimization enabled. Validation does not depend on Python
\texttt{assert} statements. Output paths must be new regular files
in trusted directories; existing paths, symlinks, and special-file
targets are rejected. The software has explicit finite workload
limits independent of the unrestricted mathematical theorems.

The PDF builder uses a clean exclusive directory, a fixed epoch,
fixed metadata, an isolated LaTeX format, disabled shell escape,
and stable references. A fixed allowlist and SHA-256 manifest
control the source ZIP; working files and third-party
journal PDFs are not included. Two clean builds and two archive
replays must be byte-identical within the recorded toolchain.
This is a same-toolchain reproducibility claim, not a promise of
identical output across different TeX versions or fonts. The
trusted-directory and modified-TeX qualifications are stated in
the README.
\begin{writenote}[6 October 2026]
The source package is shipped with the prefix \file{156-conjectures-}: \file{companion.py} and \file{test_companion.py}, \file{build_pdf.py}, \file{make_zip.py}, \file{release_tools.py} and \file{test_release.py} under \file{code/}, and the two result files, \file{data/oeis_prefixes.json} and \file{SOURCE_PROVENANCE.json} under \file{data/}. The PDF, the README (replaced by this report's) and \file{SHA256SUMS} are not shipped. The report's README gives the rerun route from the arrival archive. At intake, on a Windows copy, the exact data were regenerated byte for byte; two of the 37 companion tests fail there, because they exercise the POSIX-only safe-output code.
\end{writenote}

\section{Further questions and present limits}\label{dgp:conj:sec:questions}
\paragraph{The subdiagonal defect.}
The coupling reduces the half-partition conjecture to the fact that
interior violations are rare. It does not identify the first term of
\[
 \PP_n(R\leq0)-s_0(n)/p(n).
\]
The $O(1/\log n)$ estimate comes from a deliberately coarse union
bound and uniform rank atom estimate. Determining the actual defect
scale, its leading constant, or an all-orders expansion requires
more precise joint information on the rank, small-part
multiplicities, and interior shape. The known expansion of $p(n)$
alone does not provide an expansion of $s_0(n)$.
\begin{writenote}[6 October 2026]
\emph{Answered at $r=0$ by Part~II.} By \eqref{dgp:conj:eq:rankhalf}, $\PP_n(R\le0)=\frac12+N(0,n)/(2p(n))$. Part~II's \eqref{dgp:sub:eq:rankcubic} at $r=0$ gives $N(0,n)/p(n)=\frac\beta4+a\beta^2+O(\beta^3)$ with $a=\frac3{16}-\frac3{2\pi^2}$, and its \eqref{dgp:sub:eq:cubic} gives $s_0(n)/p(n)=\frac12-\frac\beta8+c_2\beta^2+O(\beta^3)$ with $c_2=-\frac{11}{32}+\frac3{4\pi^2}$; here $\beta=\pi/\sqrt{6n}$. Since $\frac a2-c_2=\frac3{32}-\frac3{4\pi^2}+\frac{11}{32}-\frac3{4\pi^2}$,
\begin{equation}\label{dgp:conj:eq:defect}
 \PP_n(R\le0)-\frac{s_0(n)}{p(n)}=\frac\beta4+\Bigl(\frac7{16}-\frac3{2\pi^2}\Bigr)\beta^2+O(\beta^3).
\end{equation}
So the defect has exact order $n^{-1/2}$, with leading term $\pi/(4\sqrt{6n})$; the $(\log n)^{-1}$ of the union bound \eqref{dgp:conj:eq:union} is not its true size. The leading term needs only Liu--Zhou's $N(0,n)/p(n)=\frac\beta4+O(\beta^2)$ and Part~II's $c_0$, $c_1$; the $\beta^2$ term uses Part~II's $a$ and $c_2$. A numerical check at the write, with exact $p(n)$, exact $N(0,n)$ (the Atkin--Swinnerton-Dyer rank formula, checked against enumeration for $n\le30$) and a floating-point dynamic program for $s(n)$: the defect divided by $\beta$ is $0.2957$, $0.2623$, $0.2599$ at $n=100$, $1000$, $1500$, and $(\text{defect}-\beta/4)/\beta^2$ is $0.356$, $0.303$, $0.299$ there, against $\frac7{16}-\frac3{2\pi^2}=0.2855$. Part~II also gives $s_0(n)$ itself to every order. The defect at fixed $r\ge1$ and at growing $r$ remains open (Section~\ref{dgp:sec:further}, item~\ref{dgp:q:defect}).
\end{writenote}

\paragraph{Sharper shifts.}
Theorem~\ref{dgp:conj:thm:coupling} is uniform over all nonnegative integer
$r$, including $n$-dependent shifts. A natural refinement is an
error sensitive to $r/\sqrt n$, with smaller error in regions where
the logistic tail is small. Negative integer shifts do not give a
symmetric extension of this definition: $r\leq-1$ already makes
$\lambda_1\leq1+r$ impossible for every nonempty partition.
\begin{writenote}[6 October 2026]
Open; Section~\ref{dgp:sec:further}, item~\ref{dgp:q:defect}.
\end{writenote}

\paragraph{The superdiagonal exponential constant.}
The prefix construction sacrifices information twice: it replaces
prefix weights by a common upper bound, and it bounds every
small-part subset contribution by the largest tail count. Keeping
the full weighted fixed-height transfer matrix, or using a sharper
tail comparison, may improve $B_*$. A matching upper bound and a
proof that $\log A(n)/\sqrt n$ has a limit remain open within this
report. A leading amplitude, correction series, and precise inverse
for $A(n)$ are therefore also beyond the present results.
\begin{writenote}[6 October 2026]
\emph{Answered by Part~III, and in part by Parts~IV and~V.} The limit exists and equals $B$, with the matching upper bound (Theorem~\ref{dgp:rate:thm:main}). The next term is $-Dn^{1/6}$, $D=\zeta bC^{-1/6}$, with an inverse to $o(x^{2/3})$ (Part~IV, unconditional); the amplitude $K$, the power $n^{-11/12}$ and an inverse to $o(\sqrt x)$ are Part~V's, conditional on an unrefereed preprint. A correction series remains open (Section~\ref{dgp:sec:further}, items~\ref{dgp:q:amplitude}--\ref{dgp:q:allorders}).
\end{writenote}

\paragraph{Explicit numerical onset.}
The superdiagonal finite inequalities are directly computable.
The asymptotic results would become numerical certificates only
after supplying explicit classical remainder bounds and tracking
all constants through the coupling and inverse. In particular, the
window in Theorem~\ref{dgp:conj:thm:inverse} cannot be turned into a
specified finite search interval merely by choosing an arbitrary
multiplier in its big-$O$ term.
\begin{writenote}[6 October 2026]
Open for every Part of this report; Section~\ref{dgp:sec:further}, item~\ref{dgp:q:onset}.
\end{writenote}

\clearpage
\part{All orders asymptotics for subdiagonal partitions}\label{dgp:sub:part}
\partsource{Bundle Report~159, archive \file{Subdiagonal_Partitions_All_Orders_Expansion_and_Integer_Inverses_Source.zip}, delivered as \file{Report159.tex} (17 pages). Delivered title block ``All orders asymptotics for subdiagonal partitions / Endpoint reductions and integer inverse localization / Report 159 / 3 October 2026''. Printed as delivered, with \textbf{[write]} notes. Delivered Section~$k$ is Section~$k+9$ here, and its statement or equation $k.j$ is $(k+9).j$. It treats only $s=s_0$ (shift $r=0$).}
\begin{center}
\small
\begin{tabular}{@{}>{\raggedright\arraybackslash}p{0.46\textwidth} >{\raggedright\arraybackslash}p{0.46\textwidth}@{}}
\toprule
Reading conventions of this Part & Elsewhere in the report\\
\midrule
$t=c^{-2}=6/\pi^2$, a \emph{number}; $\beta=c/\sqrt n$ the radial variable & Parts~I, III--V: $t$ a variable\\
$b=\pi/\sqrt{6(n-1/24)}$ (Zhou); $a=\frac3{16}-\frac t4$, $B(r)$ rank coefficients; $h=\frac7{24}+\frac t2$ & Part~I: $b=\pi/\sqrt3$; Parts~III--V: $b=\log2$, $B$ the rate\\
$R$ the rank and a truncation order; $C_i$, $G_b$ multiplicity and gap sums; $E_n$ failure counts & Part~I: $R$ the rank; Parts~III--V: $C_j$, $G$ a buffer\\
\bottomrule
\end{tabular}
\end{center}
\begin{partabstract}
Let $s(n)$ count weakly increasing partitions
$1\leq\lambda_1\leq\cdots\leq\lambda_k$ of $n$ satisfying
$\lambda_i\leq i$ for every index, and let $p(n)$ be the ordinary
partition number. For $\beta=\pi/\sqrt{6n}$, we prove a unique,
algorithmically computable formal expansion
\[
 \frac{s(n)}{p(n)}\sim\sum_{j\geq0}c_j\beta^j,
 \qquad c_j\in\QQ[\pi^{-2}],
\]
with a remainder $O_R(\beta^{R+1})$ after every fixed order $R$.
The first four coefficients are
\[
 c_0=\frac12,\quad c_1=-\frac18,\quad
 c_2=-\frac{11}{32}+\frac{3}{4\pi^2},\quad
 c_3=-\frac{317}{384}+\frac{329}{64\pi^2}-\frac9{8\pi^4}.
\]
A uniform endpoint estimate reduces each requested order to finitely
many rank counts with fixed shifts. A mixed row and column insertion
bijection handles intersections of bottom and top obstructions,
including explicitly enumerable, polynomially bounded exceptional
classes. Zhou's published complete fixed-rank expansion then supplies
the analytic coefficients. Six small endpoint patterns explain the
cubic term directly.

The ordinary threshold inverse has arbitrarily small algebraic
windows about computable smooth centers, with ceilings retained on
both endpoints. Its second explicit center is
\[
 x_0+\frac7{24}+\frac3{\pi^2}
 +\left(\frac{23}{32}+\frac{45}{2\pi^4}\right)
       \frac{\pi}{\sqrt{6x_0}},
\]
with window radius $O(x_0^{-1})$, where $x_0$ is the large Lambert
$W$ inverse of the leading model. Error constants and onset are not
numerically specified. The companion uses exact integer and rational
arithmetic and implements coefficient extraction only through order
three.
\end{partabstract}
\noindent\textbf{Attribution and scope.}
The sequence is OEIS A238875. Archibald, Blecher, Elizalde and
Knopfmacher formulated the half-proportion conjecture in their 2025
paper~\cite{abek}. Report~156~\cite{r156} proved it with an
$O(1/\log n)$ error by an elementary rank coupling. The present report
is a distinct refinement: it proves an all-orders endpoint reduction
and sharper integer inverse localization, using additional published
rank asymptotics. It makes no claim about a superdiagonal expansion.
A bounded check of the source literature, OEIS data, and the available
report catalogue found no duplicate of this refinement as of
3 October 2026; it is not exhaustive scholarly priority clearance.
\medskip

\section{Definitions and main results}\label{dgp:sub:sec:main}
A partition is written in weakly increasing order unless a decreasing
row notation is explicitly used. Define
\[
 s(n)=\#\{(\lambda_1,\ldots,\lambda_k):
 1\leq\lambda_1\leq\cdots\leq\lambda_k,\quad
 \textstyle\sum_i\lambda_i=n,\quad\lambda_i\leq i\},
\]
and put $s(0)=p(0)=1$ for the empty partition. The increasing order is
essential. The source definition and numerical sequence are
A238875~\cite{oeisA238875}; the first values for $n=0,\ldots,10$ are
\[
 1,1,1,2,2,4,5,7,10,15,18.
\]
For a nonempty partition let $k$ be its length, $L$ its largest part,
and $R=L-k$ its Dyson rank. Write
\[
 N(r,n)=\#\{\lambda\vdash n:R=r\},\qquad
 P_r(n)=\sum_{u\leq r}N(u,n),\qquad c=\frac\pi{\sqrt6},
 \qquad \beta=\frac c{\sqrt n}.
\]
The empty partition has rank zero. Counts at negative weight are zero.
Conjugation gives $N(r,n)=N(-r,n)$.
All limits below have the integer truncation order fixed in advance.

\begin{theorem}[Complete asymptotic expansion]\label{dgp:sub:thm:allorders}
There is a unique sequence of computable coefficients
$c_j\in\QQ[\pi^{-2}]$ such that, for every integer $R\geq0$,
\begin{equation}\label{dgp:sub:eq:allorders}
 \frac{s(n)}{p(n)}=\sum_{j=0}^{R}c_j\beta^j
                   +O_R(\beta^{R+1}).
\end{equation}
In particular,
\begin{align}\label{dgp:sub:eq:cubic}
 \frac{s(n)}{p(n)}={}&\frac12-\frac\beta8
 +\left(-\frac{11}{32}+\frac3{4\pi^2}\right)\beta^2\nonumber\\
 &+\left(-\frac{317}{384}+\frac{329}{64\pi^2}
                 -\frac9{8\pi^4}\right)\beta^3+O(\beta^4).
\end{align}
At any requested order, a terminating finite combinatorial reduction
followed by the published fixed-rank expansion computes the
coefficients. No convergence or computational efficiency is asserted.
\end{theorem}

The first correction is negative; the half-proportion is approached
from below eventually. In powers of $n$, the second-order statement is
\[
 \frac{s(n)}{p(n)}=\frac12-\frac\pi{8\sqrt{6n}}
 +\left(\frac18-\frac{11\pi^2}{192}\right)\frac1n
 +O(n^{-3/2}).
\]
The full cubic formula has remainder $O(n^{-2})$.
\begin{writenote}[6 October 2026]
``Report~156'' \cite{r156} is Part~I. Theorem~\ref{dgp:sub:thm:allorders} proves Conjecture~5.3 of \cite{abek} a second time, with a far sharper error, but only for the unshifted class $s=s_0$; Part~I's uniform statement for every shift $r$ and its logistic Corollary~\ref{dgp:conj:cor:logistic} have no counterpart here. Combined with Part~I's \eqref{dgp:conj:eq:rankhalf} it gives the defect formula \eqref{dgp:conj:eq:defect}. ``Eventually'' in the first sentence of this paragraph is needed: by exact computation (write) $s(n)/p(n)<\frac12$ for every $10\le n\le400$, while it equals $\frac12$ at $n=2$ and $n=9$ and exceeds it at $n=0,1,3,5$. Numerically (intake, $n\le1500$), $\bigl(s/p-\sum_{j\le2}c_j\beta^j\bigr)/\beta^3=-0.544,-0.629,-0.517,-0.441,-0.411,-0.388$ at $n=100,200,400,700,1000,1500$, drifting towards $c_3=-0.3162$, and the remainder after the cubic term divided by $\beta^4$ stays near $-2$ to $-3.5$; this supports $c_0,\dots,c_3$ and proves nothing.
\end{writenote}

To state the inverse, set
\begin{equation}\label{dgp:sub:eq:F}
 F(x)=\frac{e^{2c\sqrt x}}{8\sqrt3\,x},\qquad
 x_0(y)=c^{-2}W_{-1}\!\left(-\frac{c}{\sqrt{8\sqrt3\,y}}\right)^2,
\end{equation}
where $W_{-1}$ is the real branch with values at most $-1$. Thus
$F(x_0(y))=y$ on the large branch. Define the integer threshold
\begin{equation}\label{dgp:sub:eq:T}
 T(y)=\min\{n\in\ZZ_{\geq0}:s(n)\geq y\}.
\end{equation}

\begin{theorem}[Integer inverse localization]\label{dgp:sub:thm:inverse}
For each fixed $J\geq1$ there are computable coefficients
$b_0,\ldots,b_{J-1}\in\QQ[\pi^{-2}]$ such that, with
$\beta_0=c/\sqrt{x_0(y)}$ and
\begin{equation}\label{dgp:sub:eq:XJ}
 X_J(y)=x_0(y)+\sum_{r=0}^{J-1}b_r\beta_0^r,
\end{equation}
there exist constants $C_J,y_J>0$ for which
\begin{equation}\label{dgp:sub:eq:inverseceil}
 \ceil{X_J-C_Jx_0^{-J/2}}\leq T(y)
 \leq\ceil{X_J+C_Jx_0^{-J/2}}\qquad(y\geq y_J).
\end{equation}
The first coefficients and second center are
\begin{align}
 b_0&=\frac7{24}+\frac3{\pi^2},\qquad
 b_1=\frac{23}{32}+\frac{45}{2\pi^4},\label{dgp:sub:eq:b01}\\
 X_2&=x_0+\frac7{24}+\frac3{\pi^2}
       +\left(\frac{23}{32}+\frac{45}{2\pi^4}\right)
          \frac\pi{\sqrt{6x_0}}.\label{dgp:sub:eq:X2}
\end{align}
The constants $C_J,y_J$ are not numerically explicit here.
\end{theorem}

For $J\geq1$, the two ceilings in \eqref{dgp:sub:eq:inverseceil} eventually
agree or are adjacent. Agreement determines $T(y)$; disagreement
leaves two candidates. This is not an assertion that
$T(y)=X_J(y)+o(1)$, nor a certified finite-input rounding algorithm.

\section{The analytic inputs and their fixed rank specialization}\label{dgp:sub:sec:analytic}
The combinatorial part of the proof is given in full below. Its
analytic inputs are the classical partition expansion and Zhou's
complete expansion for each fixed Dyson rank. Their provenance and
specialization matter, particularly at rank zero.

\subsection{The smooth partition term}
Put
\begin{equation}\label{dgp:sub:eq:H}
 H(x)=\frac{e^{2c\sqrt{x-1/24}}}{4\sqrt3(x-1/24)}
           \left(1-\frac1{2c\sqrt{x-1/24}}\right).
\end{equation}
The Rademacher expansion~\cite{dlmfpart} gives, for every fixed
$\varepsilon\in(0,c)$,
\begin{equation}\label{dgp:sub:eq:Rademacher}
 p(n)=H(n)+O\!\left(e^{(c+\varepsilon)\sqrt n}\right).
\end{equation}
Expanding the explicit expression $H$ gives the full algebraic
expansion of $p(n)$ and each fixed shift. Finite differences may be
computed from derivatives of $H$ on fixed neighborhoods; the
exponentially smaller error in \eqref{dgp:sub:eq:Rademacher} is harmless.
In particular, for fixed positive $a_1,\ldots,a_d$,
\begin{equation}\label{dgp:sub:eq:boxordinary}
 [q^n]P(q)\prod_{h=1}^{d}(1-q^{a_h})=O(\beta^d p(n)),
 \qquad P(q)=\prod_{j\geq1}(1-q^j)^{-1}.
\end{equation}
Indeed the finite difference of $H$ is an integral of $H^{(d)}$ over
a fixed box, and $H^{(d)}(x)=O(\beta^dH(n))$ there. This reasoning
does not subtract uncontrolled big-$O$ terms.

It is convenient to write $t=c^{-2}=6/\pi^2$. Through degree three,
\begin{align}
 \frac{H'}H&=\beta-t\beta^2+\left(\frac{t^2}4+\frac t{48}\right)\beta^3
              +O(\beta^4),\label{dgp:sub:eq:Hratios}\\
 \frac{H''}H&=\beta^2-\frac{5t}2\beta^3+O(\beta^4),\qquad
 \frac{H'''}H=\beta^3+O(\beta^4).\nonumber
\end{align}
Consequently, for every fixed $w\geq0$,
\begin{align}\label{dgp:sub:eq:pshift}
 \frac{p(n-w)}{p(n)}={}&1-w\beta+\left(wt+\frac{w^2}2\right)\beta^2\nonumber\\
 &-\left[w\left(\frac{t^2}4+\frac t{48}\right)
          +\frac{5w^2t}4+\frac{w^3}6\right]\beta^3+O_w(\beta^4).
\end{align}

\subsection{An explicit published recipe for the rank coefficients}
Liu and Zhou~\cite[Theorem 1.3]{lz} give
$N(r,n)/p(n)=\beta/4+O_r(\beta^2)$ for every fixed integer $r$,
including zero. For all orders we use Zhou~\cite[Theorem 4.1]{zhou},
specialized to his parameters $k=2$, $m=0$, $j=|r|$.
These values satisfy his hypotheses. No rank increasing with $n$
occurs in any finite reduction below.

For clarity, the specialized coefficient recipe can be written
using only finite formal expansions. Set
\[
 b=\frac\pi{\sqrt{6(n-1/24)}},\quad
 A=\frac{\sqrt3\,b^2e^{\pi^2/(3b)}}{2\pi^2},\quad
 f(z)=\frac1{1+e^z},
\]
and define
\begin{align}
 d_h(j)&=(1/2-j)^h-(-1)^h(1/2+j)^h,\label{dgp:sub:eq:dh}\\
 \gamma_\ell(\mu)&=\frac1{8^\ell\ell!}
    \prod_{a=1}^{\ell}\bigl(4\mu^2-(2a-1)^2\bigr),
       \qquad\gamma_0(\mu)=1.\label{dgp:sub:eq:gamma}
\end{align}
The $m=0$ specialization of the source operators gives the formal
asymptotic recipe
\begin{equation}\label{dgp:sub:eq:rankrecipe}
 \frac{N(r,n)}A\sim
 \sum_{h\geq1}\frac{d_h(|r|)}{h!}(-b)^h
 \sum_{s,\ell\geq0}\frac{(3/2)^s}{s!}
 \left(\frac3{\pi^2}\right)^\ell
 \gamma_\ell(-s-h-3/2)f^{(h+2s)}(0)(-b)^{s+\ell}.
\end{equation}
To any requested degree, the double sum is finite after retaining
$h+s+\ell$ up to that degree. Formula \eqref{dgp:sub:eq:gamma} is the
$v=0$ coefficient from the source's Proposition~2.2. Terms containing
positive powers of the evaluation variable vanish at $m=0$.
The derivatives $f^{(a)}(0)$ are rational and are computed by formal
inversion of $1+e^z$. Since
$p(n)/A=1-3b/\pi^2$ up to exponentially small error, division and the
change from $b$ to $\beta$ give coefficients in $\QQ[\pi^{-2}]$.
Taking one additional finite order of the published expansion
supplies every big-$O$ remainder used below.

This is a use of an established coefficient-asymptotic theorem,
not a general inference from a local generating-function expansion.
In particular, the all-orders proof does not assume that an arbitrary
local power series can be transferred termwise to coefficients.
\begin{writenote}[6 October 2026]
The write read the statements of both inputs in the cited arXiv versions. Liu--Zhou's Theorem~1.3 \cite{lz}: for $m=o(n^{3/4})$, $N_k(m,n)/p(n)=\frac{\pi}{4\sqrt{6n}}\operatorname{sech}^2\bigl(\frac{\pi m}{2\sqrt{6n}}\bigr)\bigl(1+O(\frac{n+m^2}{n^{3/2}})\bigr)$, with $N_2$ the rank count; at fixed $r$ this is $\beta/4+O_r(\beta^2)$, as used, rank zero included. Zhou's Theorem~4.1 \cite{zhou} expands $N_k(m+j,n)$ for fixed $j$, with $m+j\ge0$ and $m^2\beta_n^3=o(1)$; $k=2$, $m=0$, $j=|r|$ meets these hypotheses. The write did not re-derive the recipe \eqref{dgp:sub:eq:rankrecipe} from Zhou's Proposition~2.2; the direct check of the cubic coefficient below is Part~II's own. Zhou's variable $b$ here is not Part~I's $b=\pi/\sqrt3$ nor the $b=\log2$ of Parts~III--V, and $t=6/\pi^2$ is a number in this Part (front matter, ``Notation'').
\end{writenote}

\subsection{The rank input through cubic order}
Let
\begin{equation}\label{dgp:sub:eq:ab}
 a=\frac3{16}-\frac t4,\qquad
 B(r)=\frac{53}{192}-\frac{r^2}{16}-\frac{89t}{192}+\frac{t^2}{16}.
\end{equation}
The specialization just given yields
\begin{equation}\label{dgp:sub:eq:rankcubic}
 \frac{N(r,n)}{p(n)}=\frac\beta4+a\beta^2+B(r)\beta^3+O_r(\beta^4).
\end{equation}
Here is a direct check of the third coefficient. Write $x=\pi^{-2}$.
The $b^2$ coefficient inside the $h=1$ source operator is
\[
 -\frac9{32}+\frac{27x}8-\frac{27x^2}4,
\]
using $f'(0)=-1/4$, $f'''(0)=1/8$, $f^{(5)}(0)=-1/4$,
$\gamma_1(-7/2)=6$, and $\gamma_2(-5/2)=3$.
The $h=2$ terms vanish at the needed orders because the positive even
derivatives of $f$ at zero vanish. Since $d_3(r)=3r^2+1/4$,
the $h=3$ contribution to $N/A$ is
$-b^3(r^2/16+1/192)$. Thus its cubic coefficient is
\[
 \frac{53}{192}-\frac{r^2}{16}-\frac{27x}8+\frac{27x^2}4.
\]
Dividing by $1-3bx$ and using $b=\beta+x\beta^3/8+O(\beta^5)$
gives exactly \eqref{dgp:sub:eq:rankcubic}. For fixed $r,w$ it follows that
\begin{align}\label{dgp:sub:eq:Nshift}
 \frac{N(r,n-w)}{p(n)}={}&\frac\beta4+\left(a-\frac w4\right)\beta^2\nonumber\\
 &+\left[B(r)-aw+\frac{w^2}8+\frac{3wt}8\right]\beta^3
 +O_{r,w}(\beta^4).
\end{align}
The rank-independent second coefficient and the polynomial degree
of the cubic coefficient will force the endpoint cancellations.

\section{Uniform localization of the obstructions}\label{dgp:sub:sec:localize}
Write $Z_a$ for the multiplicity of part $a$. Let $E_n$ count
partitions with $R\leq0$, $Z_1\geq1$, and at least one failed
subdiagonal inequality. Deleting one $1$ gives, for $n\geq2$, the
exact identity
\begin{equation}\label{dgp:sub:eq:source}
 s(n)=P_1(n-1)-E_n.
\end{equation}
For the all-ones partition of size at least two, deletion also
preserves the rank-threshold equivalence; the singleton gives the
same identity separately. Hence \eqref{dgp:sub:eq:source} also holds at
$n=1$. We only use it asymptotically.

\subsection{Boltzmann conditioning}
Under the partition Boltzmann probability $Q_\beta$, a partition
$\lambda$ has probability $q^{|\lambda|}/P(q)$ with $q=e^{-\beta}$.
The random variables $Z_a$ are independent and
\[
 Q_\beta(Z_a=z)=(1-q^a)q^{az},\qquad z\geq0.
\]
For any set of partitions $\mathcal E$, positivity gives
\begin{equation}\label{dgp:sub:eq:condition}
 \PP_n(\mathcal E)\leq\frac{q^{-n}P(q)}{p(n)}Q_\beta(\mathcal E)
                       =O(n)Q_\beta(\mathcal E).
\end{equation}
To check the last estimate without a local limit theorem, the positive
decreasing function $-\log(1-e^{-\beta x})$ has integral
$\pi^2/(6\beta)$, so
\[
 \log P(e^{-\beta})\leq\frac{\pi^2}{6\beta},\qquad
 q^{-n}P(q)\leq e^{2c\sqrt n},\qquad p(n)\asymp n^{-1}e^{2c\sqrt n}.
\]
Let $C_i=Z_1+\cdots+Z_i$. There are
$\binom{2i-1}{i}\leq4^i$ nonnegative $i$-tuples with sum at most
$i-1$. Therefore
\begin{equation}\label{dgp:sub:eq:Cbound}
 Q_\beta(C_i\leq i-1)\leq4^i\prod_{h=1}^i(1-q^h)
                    \leq4^i\beta^ii!\leq(4\beta i)^i.
\end{equation}
Conjugation preserves this probability and exchanges small-part
multiplicities with the top gaps used below.

\subsection{Excluding the middle and short partitions}
Choose
\begin{equation}\label{dgp:sub:eq:cutoffs}
 K=\floor{\frac{\sqrt n\log n}{4c}},\qquad
 m=\ceil{\frac{5c\sqrt n}{\log n}}.
\end{equation}
Eventually $K\geq2m+1$ and $(m+1)(K-m+1)>n$, since the ratio of
the latter product to $n$ tends to $5/4$.
A failure at an increasing index $i$ forces
\[
 n\geq(i+1)(k-i+1).
\]
If $k\geq K$ and $m<i\leq k-m$, concavity in $i$ shows that the
right side is at least $(m+1)(k-m+1)>n$. Thus every failure is at
$i\leq m$ or at a decreasing index $j=k-i+1\leq m$.
This uses the same interior mechanism as Report~156, here retained
with much finer endpoint information.

By conjugation and coefficient positivity, the short-length event
satisfies
\begin{align}\label{dgp:sub:eq:short}
 \PP_n(k<K)&\leq\frac{q^{-n}P(q)}{p(n)}\prod_{a>K}(1-q^a)\nonumber\\
 &\leq O(n)\exp\!\left(-\frac{q^{K+1}}{1-q}\right)
   =O(ne^{-c_0n^{1/4}})
\end{align}
for some $c_0>0$, because $q^{K+1}/(1-q)\sim c^{-1}n^{1/4}$.
This is smaller than every fixed power of $\beta$.
\begin{writenote}[6 October 2026]
The cutoffs \eqref{dgp:sub:eq:cutoffs}, the interior step of the preceding paragraph and the short-length bound \eqref{dgp:sub:eq:short} restate Part~I's \eqref{dgp:conj:eq:scales}, Lemma~\ref{dgp:conj:lem:interior} and \eqref{dgp:conj:eq:productbound}--\eqref{dgp:conj:eq:length}. They are kept, since the endpoint analysis below starts from them; the conditioning bound \eqref{dgp:sub:eq:condition} is the estimate Part~I proves in Section~\ref{dgp:conj:sec:interior} for $z^{-n}P(z)/p(n)$.
\end{writenote}

\subsection{Fixed endpoint cutoffs with arbitrary accuracy}
Use decreasing rows $\mu_1\geq\mu_2\geq\cdots$, padded by zero,
and gaps $G_b=\mu_b-\mu_{b+1}$. These gaps are the multiplicities
of parts $b$ in the conjugate partition. In the class counted by
$E_n$, a bottom failure at $i$ means $C_i\leq i-1$; necessarily
$i\geq2$. At the top put $R=-d$, $d\geq0$. A failure at decreasing
index $j$ is equivalent to
\begin{equation}\label{dgp:sub:eq:topfailure}
 0\leq d\leq j-2,\qquad
 G_1+\cdots+G_{j-1}\leq j-2-d,
\end{equation}
since $L=k-d$ and $\mu_j\geq k-j+2$.

Fix $M\geq2$. Let $E_n^{(M)}$ be the union of the bottom events
\[
 R\leq0,\quad Z_1\geq1,\quad C_i\leq i-1
                     \quad(2\leq i\leq M)
\]
and the top events
\[
 R=-d,\quad Z_1\geq1,\quad
 \sum_{b=1}^{j-1}G_b\leq j-2-d,
 \quad 0\leq d\leq j-2,\quad 2\leq j\leq M+1.
\]
These really are failure events for all $n>(M+1)^2$: if a required
index does not exist, $R\leq0$ and length at most $M$ force the
weight to be at most $M^2$.

For $a_i=(4\beta i)^i$ and $M+1\leq i<m$,
\[
 \frac{a_{i+1}}{a_i}
 =4\beta(i+1)(1+1/i)^i\leq4e\beta(m+1)\leq\frac12
\]
eventually. Hence
\[
 \sum_{i=M+1}^{m}(4\beta i)^i=O_M(\beta^{M+1}).
\]
The omitted top events have $j\geq M+2$, so their gap-prefix length
$j-1$ is also at least $M+1$. Equations \eqref{dgp:sub:eq:condition},
\eqref{dgp:sub:eq:Cbound}, and \eqref{dgp:sub:eq:short} now prove
\begin{equation}\label{dgp:sub:eq:tail}
 0\leq E_n-E_n^{(M)}=O_M(\beta^{M-1}p(n)).
\end{equation}
The loss of two powers comes from $O(n)=O(\beta^{-2})$ in
conditioning. This is a uniform bound on a growing endpoint range;
no fixed-rank expansion has been used outside a fixed range.

\section{Exact mixed row and column reductions}\label{dgp:sub:sec:mixed}
The all-orders argument must handle simultaneous conditions at both
ends. Separate bottom and top estimates do not suffice.

\begin{lemma}[Small rectangle complements]\label{dgp:sub:lem:rectangle}
For fixed integers $B\geq1$, $C\geq0$, the number of partitions of
$n$ with $\mu_B\leq C$ is
$O_{B,C}((n+1)^{B+C-1})$. These counts remain explicitly enumerable
after imposing fixed rank or additional marked-coordinate conditions.
\end{lemma}
\begin{proof}
Specify at most $B-1$ initial row lengths, each at most $n$, and
the multiplicities of the possible sizes $1,\ldots,C$ below these
rows. There are at most $(n+1)^{B+C-1}$ descriptions. Enforcing
ordering and total weight only reduces this upper bound. All entries
can be enumerated in finite boxes and tested against the extra
conditions.
\end{proof}
Such exceptional counts are polynomially bounded; they need not be
polynomial functions of $n$. Dividing any fixed polynomial bound by
$p(n)$ gives an error smaller than every algebraic power of $\beta$.

\begin{lemma}[Mixed lower bounds]\label{dgp:sub:lem:mixed}
Fix $A\geq0$, $B\geq1$, and nonnegative integer vectors
$x=(x_1,\ldots,x_A)$, $y=(y_1,\ldots,y_B)$. Let
$L_{x,y,r}(n)$ count partitions of rank $r$ satisfying
$Z_a\geq x_a$ for $a\leq A$ and $G_b\geq y_b$ for $b\leq B$.
Set
\[
 u=\sum_a x_a,\qquad v=\sum_b y_b,\qquad
 w=\sum_a ax_a+\sum_b by_b.
\]
Then the exact identity is
\begin{equation}\label{dgp:sub:eq:mixed}
 L_{x,y,r}(n)=N(r+u-v,n-w)-B_{\rm src}+B_{\rm tgt},
\end{equation}
where $B_{\rm src}$ counts cores of weight $n-w$, rank $r+u-v$,
and $\nu_B\leq A$, while $B_{\rm tgt}$ counts the target partitions
in $L_{x,y,r}(n)$ with $\mu_B\leq A+y_B$.
Both complements are polynomially bounded and explicitly enumerable.
\end{lemma}
\begin{proof}
Start with a core $\nu$ having $\nu_B>A$. Insert $x_a$ rows of
length $a$ for each $a\leq A$. They lie below the first $B$ rows,
which do not change. Now add $y_b$ columns of height $b$, for
$b\leq B$; equivalently, add $\sum_{j=b}^{B}y_j$ to row $b$.
Each gap $G_b$ increases by $y_b$ in this second step. Inserting the
small rows can change the boundary gap $G_B$, since row $B+1$ can
change, but the gap remains nonnegative. Thus the final lower bounds
on every gap hold. No invariance of $G_B$ under row insertion is
required.

The affected top rows remain larger than $A$, so column insertion
preserves all marked small-part multiplicities. Both operations act
on the same top rows and inserted small rows in either order, so
they commute on this region. Weight increases by $w$, length by
$u$, and largest part by $v$. The output therefore has rank $r$
and $\mu_B=\nu_B+y_B>A+y_B$.

Conversely, remove the indicated columns from any target with
$\mu_B>A+y_B$. The gap lower bounds make all resulting gaps
nonnegative. Its first $B$ rows remain larger than $A$, hence no
marked small-part multiplicity changes. Remove the specified small
rows; the first $B$ rows stay fixed and the original good core is
recovered. No gap lower bound is imposed on that core. These maps
are mutually inverse.

The good source and target sets are therefore equinumerous.
Subtracting the bad-source complement and adding the independently
defined bad-target complement gives \eqref{dgp:sub:eq:mixed}. The latter
are covered by Lemma~\ref{dgp:sub:lem:rectangle}. Empty cores and negative
core weights cause no missing case under the stated conventions.
In particular, targets with $n<w$ are necessarily bad targets.
\end{proof}

When $B=0$, the same argument inserts only small rows, using the
good-core condition ``largest part greater than $A$.'' Both
complements have largest part at most $A$ and are polynomially
bounded. The rank shift is $r+u$. This pure-row identity also holds
for the cumulative condition $R\leq r$, with $P_{r+u}(n-w)$ in
place of the rank atom.

Exact values are obtained by finite differences:
\begin{equation}\label{dgp:sub:eq:exactdiff}
 \one_{Z_a=x_a}=\one_{Z_a\geq x_a}-\one_{Z_a\geq x_a+1},
 \qquad
 \one_{G_b=y_b}=\one_{G_b\geq y_b}-\one_{G_b\geq y_b+1}.
\end{equation}
Consequently every finite pattern of marked values at fixed rank
is a finite integer combination of shifted $N(r,n-w)$ plus
polynomially bounded complements. Pure-row cumulative patterns
reduce similarly to $P_r(n-w)$. Rank symmetry gives the exact
conversion, for every integer $r$,
\begin{equation}\label{dgp:sub:eq:cdf}
 P_r(n)=
 \begin{cases}
 \displaystyle\frac{p(n)}2+\frac{N(0,n)}2+\sum_{j=1}^{r}N(j,n),&r\geq0,\\[6pt]
 \displaystyle\frac{p(n)}2-\frac{N(0,n)}2-\sum_{j=1}^{-r-1}N(j,n),&r<0.
 \end{cases}
\end{equation}
Empty sums vanish. The factor $1/2$ is the only source of rational
rather than integer coefficients in this combinatorial reduction.

\section{The finite algorithm at every order}\label{dgp:sub:sec:algorithm}
We now prove the existence assertion in Theorem~\ref{dgp:sub:thm:allorders}.
Apply inclusion--exclusion to the finitely many bottom and top
events defining $E_n^{(M)}$.

In an intersection containing bottom indices $I$, let $A=\max I$.
The largest-index condition bounds $Z_1+\cdots+Z_A\leq A-1$;
together with $Z_1\geq1$ it leaves finitely many vectors $x$.
The other bottom conditions select a subset of this finite list.
If top indices $J$ occur, let $B=\max J-1$. Their ranks are
$-d$, where $0\leq d\leq\min J-2$. The largest-index condition
bounds $G_1+\cdots+G_B\leq B-1-d$, leaving finitely many vectors
$y$ for each of finitely many ranks.

For a bottom-only intersection, use the cumulative pure-row
reduction. For a top-only intersection, keep the exact gaps and
rank and impose the single lower bound $Z_1\geq1$ in
Lemma~\ref{dgp:sub:lem:mixed}; no enumeration over arbitrarily many ones
is needed. For a mixed intersection, use exact $x$ and $y$ at
fixed rank. Thus \eqref{dgp:sub:eq:exactdiff} and \eqref{dgp:sub:eq:cdf} produce
explicit finite lists with
\begin{equation}\label{dgp:sub:eq:finite}
 E_n^{(M)}=\sum_t\alpha_t p(n-w_t)
              +\sum_t\eta_t N(r_t,n-w_t)+e_M(n),
\end{equation}
where $\alpha_t,\eta_t$ are rational, all ranks and nonnegative
shifts are fixed, and $e_M(n)$ is an explicitly enumerable signed
combination of polynomially bounded complement counts.
The function $e_M$ itself is polynomially bounded, without being
asserted polynomial.

For a requested $R\geq0$, choose $M=R+2$. The tail
\eqref{dgp:sub:eq:tail} is then $O_R(\beta^{R+1}p(n))$.
The exceptional term in \eqref{dgp:sub:eq:finite}, divided by $p(n)$, is
smaller than every fixed power of $\beta$. Expand its finitely many
shifted partition counts by \eqref{dgp:sub:eq:H} and its finitely many
rank atoms by Zhou's theorem, or equivalently by its specialized
recipe \eqref{dgp:sub:eq:rankrecipe}. Take enough source terms to leave
an $O_R(\beta^{R+1})$ error. Finally subtract from
$P_1(n-1)/p(n)$ using \eqref{dgp:sub:eq:source} and \eqref{dgp:sub:eq:cdf}.
This proves \eqref{dgp:sub:eq:allorders}, apart from the explicit coefficients.

For reference, the terminating procedure is:
\begin{enumerate}
\item Fix $R$, set $M=R+2$, and list the endpoint events
\item Form their finite intersections by inclusion--exclusion
\item Enumerate each bounded multiplicity vector, gap vector, and rank
\item Replace exact coordinates by lower-bound differences
\item Apply the mixed or pure-row bijection, retaining its complements
\item Convert cumulative ranks to $p$ and finitely many rank atoms
\item Expand the fixed shifts with the published analytic formulas
      and collect powers of $\beta$
\end{enumerate}
Every operation is finite. Its naive running time may be large.
All formal operations involve rational numbers and powers of
$\pi^{-2}$, so the coefficient field is as claimed.

Finally, expansions derived with different endpoint cutoffs agree
to their common order. Otherwise subtraction and division by the
first differing power of $\beta$ would give a nonzero constant
converging to zero. The strict asymptotic scale
$1,\beta,\beta^2,\ldots$ therefore determines a single unique
formal series. This proves the all-orders part of the theorem.

\section{The first two corrections}\label{dgp:sub:sec:firsttwo}
The leading corrections can be isolated without enumerating the
full finite reduction. This also shows why one must include both
ends of the partition.

At fixed $r$, \eqref{dgp:sub:eq:rankcubic} begins
$N(r,n)/p(n)=\beta/4+a\beta^2+O_r(\beta^3)$.
For fixed positive $A,B$, deletion of small parts and
inclusion--exclusion give, for all sufficiently large $n$,
\begin{align}\label{dgp:sub:eq:rankbox}
 \#\{R=r,Z_1<A,Z_2<B\}={}&N(r,n)-N(r+A,n-A)\nonumber\\
 &-N(r+B,n-2B)+N(r+A+B,n-A-2B).
\end{align}
The only largest-part exceptions are partitions supported on
$\{1,2\}$; at fixed rank these have bounded size.
The leading terms and the shift-linear second coefficients cancel,
so \eqref{dgp:sub:eq:rankbox} is $O(\beta^3p(n))$.
Likewise, a fixed ordinary box in $Z_1,Z_2,Z_3$ has
$O(\beta^3p(n))$ members by \eqref{dgp:sub:eq:boxordinary}.

Use $M=4$ in the endpoint truncation. Bottom failures at $i=3,4$
lie in a bounded three-coordinate box. Top failures at $j=3,4,5$,
after conjugation, have rank in $\{0,1,2,3\}$ and two bounded
small-part multiplicities; \eqref{dgp:sub:eq:rankbox} applies. Therefore only
\[
 \mathcal B_n=\{R\leq0,Z_1=1,Z_2=0\},\qquad
 \mathcal T_n=\{R=0,\mu_1=\mu_2,Z_1\geq1\}
\]
survive through order two. Their overlap is contained in a fixed-rank
two-coordinate box and is $O(\beta^3p(n))$. Hence
\begin{equation}\label{dgp:sub:eq:Esecond}
 E_n=|\mathcal B_n|+|\mathcal T_n|+O(\beta^3p(n)).
\end{equation}
For all sufficiently large $n$, exact deletion identities are
\begin{align}
 |\mathcal B_n|&=P_1(n-1)-P_2(n-2)-P_2(n-3)+P_3(n-4),\label{dgp:sub:eq:Bexact}\\
 |\mathcal T_n|&=N(1,n-1)-N(0,n-2).\label{dgp:sub:eq:Texact}
\end{align}
For \eqref{dgp:sub:eq:Bexact}, first enforce exactly one $1$, then exclude
partitions having a $2$. For \eqref{dgp:sub:eq:Texact}, delete one $1$,
conjugate, and count rank $-1$ partitions without ones by deleting
a $1$ from the complementary class. The small largest-part
exceptions are absent at sufficiently large weight.

Since
\[
 \frac{P_r(n)}{p(n)}=\frac12+(r+1/2)(\beta/4+a\beta^2)
                       +O_r(\beta^3)\quad(r\geq0),
\]
substitution of the fixed-shift formulas yields
\[
 \frac{|\mathcal B_n|}{p(n)}=\frac{\beta^2}4+O(\beta^3),
 \qquad
 \frac{|\mathcal T_n|}{p(n)}=\frac{\beta^2}4+O(\beta^3).
\]
Thus $E_n/p(n)=\beta^2/2+O(\beta^3)$, while
\begin{equation}\label{dgp:sub:eq:Psource2}
 \frac{P_1(n-1)}{p(n)}=\frac12-\frac\beta8
                       +\left(\frac5{32}+\frac t8\right)\beta^2
                       +O(\beta^3).
\end{equation}
Subtracting proves $c_0,c_1,c_2$ as stated. In particular the first
coefficient comes from the cancellation of the one-part shift
$-\beta/2$ with the three half-rank contributions $3\beta/8$.
The two surviving endpoint obstructions first affect the second
coefficient.

\section{Six patterns determine the cubic coefficient}\label{dgp:sub:sec:cubic}
Set $M=5$, so the omitted endpoint tail is $O(\beta^4p(n))$.
Through cubic order, \eqref{dgp:sub:eq:Nshift} has total degree at most two
in the fixed rank and shift variables. The cumulative version for
nonnegative threshold has total degree at most three: the sum of
$j^2$ in \eqref{dgp:sub:eq:cdf} is cubic in its upper endpoint.
Each exact coordinate in \eqref{dgp:sub:eq:exactdiff} is a finite difference
translating the rank and shift by fixed constants, lowering total
polynomial degree by at least one. Consequently:
\begin{itemize}
\item a fixed-rank pattern with at least three exact coordinates
      contributes $O(\beta^4p(n))$
\item a pure top pattern with at least three exact gaps has that
      bound, even with the lower bound of one $1$
\item a pure bottom pattern with at least four exact multiplicities
      has that bound
\end{itemize}
The polynomially bounded rectangle complements remain negligible.
All cumulative thresholds in the pure-bottom reduction are
nonnegative, so the polynomial statement applies throughout.

Classify failures by their first failing bottom or top index. Only
the following first-failure patterns remain through degree three;
within each orientation they are disjoint:
\begin{center}
\begin{tabular}{c c c}
\toprule
& Bottom exact multiplicities & Top rank and exact gaps\\
\midrule
1 & $(Z_1,Z_2)=(1,0)$ & $d=0,\ (G_1)=(0)$\\
2 & $(Z_1,Z_2,Z_3)=(1,1,0)$ & $d=0,\ (G_1,G_2)=(1,0)$\\
3 & $(Z_1,Z_2,Z_3)=(2,0,0)$ & $d=1,\ (G_1,G_2)=(0,0)$\\
\bottomrule
\end{tabular}
\end{center}
Here the bottom rank condition is $R\leq0$, the top condition is
$R=-d$, and every pattern has $Z_1\geq1$. For example, first bottom
failure at $i=3$ forces $Z_1+Z_2=2$, $Z_3=0$, with $Z_1\geq1$;
this gives exactly patterns 2 and 3. The top classification follows
similarly from \eqref{dgp:sub:eq:topfailure}. Every bottom/top intersection
has at least two exact multiplicities and one exact gap at a fixed
rank, hence contributes only $O(\beta^4p(n))$.

Here are explicit finite formulas from which the coefficients can
be checked without any numerical fit. For a vector $v$, write
$|v|=\sum_a v_a$ and $\mathrm{wt}(v)=\sum_a av_a$.
For a bottom vector $x$ and top vector $y$, the relevant counts,
up to polynomially bounded complements, are
\begin{align}
 \mathcal B(x;n)&=\sum_{\epsilon\in\{0,1\}^{A}}(-1)^{|\epsilon|}
  P_{|x+\epsilon|}(n-\mathrm{wt}(x+\epsilon)),\label{dgp:sub:eq:Bpattern}\\
 \mathcal T(d,y;n)&=\sum_{\epsilon\in\{0,1\}^{B}}(-1)^{|\epsilon|}
  N(1-d-|y+\epsilon|,n-1-\mathrm{wt}(y+\epsilon)).\label{dgp:sub:eq:Tpattern}
\end{align}
For the latter, the single mandatory $1$ supplies the rank shift
$+1$ and weight shift $1$ in Lemma~\ref{dgp:sub:lem:mixed}.
Substitution of \eqref{dgp:sub:eq:pshift}, \eqref{dgp:sub:eq:Nshift}, and
\eqref{dgp:sub:eq:cdf} gives the same three normalized contributions on
each side:
\begin{equation}\label{dgp:sub:eq:three}
 \frac{\beta^2}4-\left(\frac14+\frac{5t}8\right)\beta^3,
 \qquad\frac{3\beta^3}8,
 \qquad\frac{3\beta^3}8,
 \qquad\text{each up to }O(\beta^4).
\end{equation}
For instance the first top expression is the difference
$N(1,n-1)-N(0,n-2)$; its cubic coefficient from
\eqref{dgp:sub:eq:Nshift} is $-1/4-5t/8$. The other five follow from the
same finite sums \eqref{dgp:sub:eq:Bpattern}--\eqref{dgp:sub:eq:Tpattern}.
They may equally be verified by exact rational polynomial arithmetic.

Adding both orientations and discarding their negligible
intersections gives
\begin{equation}\label{dgp:sub:eq:Ecubic}
 \frac{E_n}{p(n)}=\frac{\beta^2}2+\left(1-\frac{5t}4\right)\beta^3
                      +O(\beta^4).
\end{equation}
The source term is
\begin{align}\label{dgp:sub:eq:sourcecubic}
 \frac{P_1(n-1)}{p(n)}={}&\frac12-\frac\beta8
 +\left(\frac5{32}+\frac t8\right)\beta^2\nonumber\\
 &+\left(\frac{67}{384}-\frac{151t}{384}-\frac{t^2}{32}\right)\beta^3
 +O(\beta^4).
\end{align}
Subtracting \eqref{dgp:sub:eq:Ecubic} yields
\[
 c_3=-\frac{317}{384}+\frac{329t}{384}-\frac{t^2}{32}
     =-\frac{317}{384}+\frac{329}{64\pi^2}-\frac9{8\pi^4}.
\]
This completes the explicit part of Theorem~\ref{dgp:sub:thm:allorders}.
The finite full-cutoff computation offers a second organization of
the same algebra: at $M=5$ it uses 22 bottom patterns, 64 top patterns,
and 1408 mixed intersections, reducing to 201 nonzero shifted terms.
The six-pattern argument explains why most of those terms cancel.

\section{Smooth models and correctly rounded integer inverses}\label{dgp:sub:sec:inverse}
Prepending one $1$ to a subdiagonal partition preserves all its
inequalities and is injective. Thus $s(n)$ is nondecreasing.
Theorem~\ref{dgp:sub:thm:allorders} and the partition asymptotic show that
$s(n)$ is unbounded, so \eqref{dgp:sub:eq:T} exists for all large $y$.
Putting $z=c\sqrt{x}$ in $F(x)=y$ gives
$ze^{-z}=c/\sqrt{8\sqrt3\,y}$; the large solution is exactly
\eqref{dgp:sub:eq:F} on the $W_{-1}$ branch.

\subsection{A general model window}
Multiplying the complete ratio expansion by the classical partition
expansion produces, for every fixed $J\geq0$,
\begin{equation}\label{dgp:sub:eq:countmodel}
 s(n)=F(n)\left(1+\sum_{j=1}^{J}a_j\beta_n^j
                       +O_J(\beta_n^{J+1})\right),
 \qquad\beta_x=\frac c{\sqrt{x}}.
\end{equation}
Define the explicit smooth function
\begin{equation}\label{dgp:sub:eq:MJ}
 M_J(x)=F(x)\left(1+\sum_{j=1}^{J}a_j\beta_x^j\right).
\end{equation}
It is eventually positive and strictly increasing because
\begin{equation}\label{dgp:sub:eq:MJderiv}
 (\log M_J)'(x)=\frac c{\sqrt{x}}+O_J(x^{-1})>0.
\end{equation}
Only an explicit finite model is differentiated. Let $\nu_J(y)$ be
its unique sufficiently large solution of $M_J(\nu_J)=y$.
Then $\nu_J\sim x_0$.

\begin{lemma}[Ceiling window for the model root]\label{dgp:sub:lem:window}
For some constants $D_J,y_J>0$ and every $y\geq y_J$,
\begin{equation}\label{dgp:sub:eq:nuwindow}
 \ceil{\nu_J-D_J\nu_J^{-J/2}}\leq T(y)
 \leq\ceil{\nu_J+D_J\nu_J^{-J/2}}.
\end{equation}
\end{lemma}
\begin{proof}
Write $\nu=\nu_J$ and $\delta=D\nu^{-J/2}$. The logarithmic
counting error in \eqref{dgp:sub:eq:countmodel}, relative to $M_J$, is
$O_J(\nu^{-(J+1)/2})$ at integers in any fixed neighborhood of
$\nu$. Take the boundary integers
\[
 n_- =\ceil{\nu-\delta}-1,\qquad n_+=\ceil{\nu+\delta}.
\]
For fixed $D$ they lie in such a neighborhood, including when
$J=0$. Their distances below and above $\nu$ are at least
$\delta$. By \eqref{dgp:sub:eq:MJderiv}, the magnitude of the corresponding
change in $\log M_J$ is at least a fixed positive multiple of
$D\nu^{-(J+1)/2}$. Choose $D$ sufficiently large to dominate the
counting error. Then $s(n_-)<y$ and $s(n_+)\geq y$ for all
sufficiently large $y$. Monotonicity gives
$n_-+1\leq T(y)\leq n_+$, which is \eqref{dgp:sub:eq:nuwindow}.
The strict lower boundary is valid whether either real endpoint
is integral or nonintegral.
\end{proof}

\subsection{The first two explicit centers}
The smooth partition term gives
\begin{equation}\label{dgp:sub:eq:partrel}
 \frac{p(n)}{e^{2c\sqrt n}/(4\sqrt3\,n)}
 =1-\left(\frac1{24}+\frac t2\right)\beta
    +\left(\frac1{1152}+\frac t{16}\right)\beta^2+O(\beta^3).
\end{equation}
Combining it with the proved ratio expansion yields
\begin{equation}\label{dgp:sub:eq:ha2}
 a_1=-h,\qquad h=\frac7{24}+\frac t2,
 \qquad a_2=-\frac{779}{1152}+\frac{7t}{16}.
\end{equation}
Both the ordinary partition correction and the subdiagonal
correction contribute to $h$.
The first root therefore satisfies
$\nu_1=x_0+h+O(x_0^{-1/2})$. Lemma~\ref{dgp:sub:lem:window} gives
\[
 \ceil{x_0+h-Cx_0^{-1/2}}\leq T(y)
 \leq\ceil{x_0+h+Cx_0^{-1/2}}
\]
for some fixed $C$ and sufficiently large $y$.

For the next center put $x=x_0+h+k\beta_0$. Direct expansion of
the explicit functions gives
\begin{align}
 \log F(x)-\log F(x_0)&=h\beta_0+(k-ht)\beta_0^2+O(\beta_0^3),\nonumber\\
 \log(1-h\beta_x+a_2\beta_x^2)&=-h\beta_0
             +(a_2-h^2/2)\beta_0^2+O(\beta_0^3).\label{dgp:sub:eq:centercancel}
\end{align}
Thus
\[
 k=ht-a_2+h^2/2=\frac{23}{32}+\frac{5t^2}8
               =\frac{23}{32}+\frac{45}{2\pi^4}.
\]
The logarithmic residual of $M_2$ at this center is $O(\beta_0^3)$.
Dividing by its derivative scale $\asymp\beta_0$ gives
$\nu_2-X_2=O(\beta_0^2)=O(x_0^{-1})$. Absorbing this into the
$J=2$ model window proves
\begin{equation}\label{dgp:sub:eq:X2window}
 \ceil{X_2-C/x_0}\leq T(y)\leq\ceil{X_2+C/x_0}.
\end{equation}

\subsection{Every fixed order and the remaining rounding ambiguity}
For general $J\geq1$, substitute the ansatz \eqref{dgp:sub:eq:XJ} in
$\log M_J(X_J)-\log y$ and equate coefficients through degree $J$
in $\beta_0$. The new coefficient $b_r$ first enters with coefficient
one at degree $r+1$, because the leading logarithmic derivative is
$\beta_0+O(\beta_0^2)$. This is a triangular rational recursion.
The resulting logarithmic residual is $O_J(\beta_0^{J+1})$,
so the mean value theorem applied only to the explicit model gives
\[
 \nu_J-X_J=O_J(\beta_0^J)=O_J(x_0^{-J/2}).
\]
Substitution in Lemma~\ref{dgp:sub:lem:window} proves
Theorem~\ref{dgp:sub:thm:inverse}, including the coefficient field.
There is no uniformity claim with $J$ growing with $y$.
\begin{writenote}[6 October 2026]
The centre $x_0$ is Part~I's \eqref{dgp:conj:eq:inversecenter}, and the injection that prepends a~$1$ is Part~I's (Section~\ref{dgp:conj:sec:inverse}). Theorem~\ref{dgp:sub:thm:inverse} sharpens Part~I's Theorem~\ref{dgp:conj:thm:inverse} from a window of width $O(\log y/\log\log y)$ to two adjacent integers. In the volume's terms (front matter, ``The inverses''): $x_0$ is an instance of \file{p0:thm:lambert-core} after the reduction $\xi=\sqrt x$, and $X_J$ is an instance of the reversion \file{p0:thm:lambert-centered} for the model \eqref{dgp:sub:eq:MJ} (the write checked $b_0=h$ and $b_1=ht-a_2+h^2/2$ against the volume's recurrence); the ceiling windows of Lemma~\ref{dgp:sub:lem:window} are Part~II's own, not an instance of \file{p0:thm:staircase}.
\end{writenote}

The ceilings cannot in general be removed. The integer threshold
may stay an order-one distance from the smooth center even though
the pre-rounding uncertainty tends to zero. Nor do unspecified
$C_J,y_J$ certify a numerical bracket at a chosen finite input.
The exact threshold command in the companion instead searches
actual integer counts in a bounded range and reports explicitly
when the requested value has not been reached.

\section{Exact computation and reproducibility}\label{dgp:sub:sec:compute}
The source package contains a Python standard-library companion,
exact deterministic JSON data, regression tests, a clean PDF
builder, and a fixed-allowlist ZIP builder. The mathematical theorem
holds at every fixed order; the software deliberately caps
coefficient extraction at $R\leq3$.

The counting algorithm records length $k$, weight $n$, and largest
part $a$. If $D_k(n,a)$ counts increasing subdiagonal partitions
with these parameters, then
\[
 D_k(n,a)=\sum_{b\leq a}D_{k-1}(n-a,b),\qquad 1\leq a\leq k,
 \qquad D_0(0,0)=1.
\]
Appending $a$ enforces precisely the new inequality $a\leq k$.
Prefix sums give an exact positive dynamic program. An independent
recursive partition enumerator checks the resulting counts and
directly tests endpoint events, mixed insertion maps, and
exceptional-set identities over the recorded finite ranges.
Numeric OEIS prefixes are attributed to their original entry;
third-party journal PDFs and full OEIS entry text are not
redistributed.

The coefficient implementation uses exact rational polynomials in
$t=6/\pi^2$ and the source-checked shifted formulas
\eqref{dgp:sub:eq:pshift} and \eqref{dgp:sub:eq:Nshift}. It performs the finite
endpoint construction, including mixed intersections, and checks
agreement with the six-pattern cubic calculation. It neither fits
coefficients to numerical data nor asserts validity of a derivative
transfer rule beyond the checked order.

For orientation, exact counts are
\begin{center}
\begin{tabular}{r r r}
\toprule
$n$ & $p(n)$ & $s(n)$\\
\midrule
25 & 1958 & 898\\
40 & 37338 & 17319\\
50 & 204226 & 95502\\
100 & 190569292 & 91171308\\
200 & 3972999029388 & 1930845918142\\
\bottomrule
\end{tabular}
\end{center}
These finite values do not establish any asymptotic remainder
constant, onset, or novelty claim. The exact verification outputs
and the README state all workload limits and check counts.

Regression tests run normally and with Python optimization enabled;
validation does not depend on Python \texttt{assert} statements.
Output destinations must be fresh files in trusted directories.
The release helpers reject existing output paths and static symlink
components. Their guarantees assume no malicious same-user process
concurrently replacing the parent or source tree; they are not an
operating-system sandbox.

The PDF is built in an exclusive clean directory with isolated
LaTeX state, disabled shell escape, fixed metadata and epoch, and
references stabilized before the settled log is inspected. The ZIP
uses a fixed allowlist and sorted SHA-256 manifest, excluding
working files, review records, and other reports. Two clean PDF
builds and two ZIP replays are required to agree byte for byte
within the recorded toolchain. Reproducibility across arbitrary
TeX versions, fonts, or operating systems is not promised. The
README gives the exact replay commands and qualifications.
\begin{writenote}[6 October 2026]
The package is shipped with the prefix \file{159-subdiag-} (code under \file{code/}; the three result files, the OEIS prefix and the provenance record under \file{data/}; the PDF, README and \file{SHA256SUMS} are not shipped). At intake, on a Windows copy, the exact data, including \file{endpoint_expansion.json}, were regenerated byte for byte; 31 companion tests ran without failure apart from three errors and eight skips of POSIX-only output tests. The report's README gives the rerun route.
\end{writenote}

\section{Source record and limits}\label{dgp:sub:sec:sources}
The direct antecedent is Conjecture~5.3 in Archibald et
al.~\cite{abek}, asserting $s(n)/p(n)\to1/2$. Their definitions,
generating functions, and rank-zero injection are prior results.
Report~156~\cite{r156} supplies the earlier elementary quantitative
proof and its coarse inverse; the present treatment repeats the
needed localization estimates and does not require that report as
an unproved lemma.

The additional analytic inputs are Liu--Zhou's fixed-rank local
estimate, expressly valid at zero, and Zhou's complete expansion.
The cited preprint versions are arXiv:1903.00835v3 and
arXiv:2110.11174v3. The latter's Theorem~4.1 and Proposition~2.2
are the source of the general recipe and the direct cubic
calculation. These results, and classical Hardy--Ramanujan and
Rademacher theory, are not claimed as new.

A bounded catalogue and source check found no identical all-orders
application to this exact subdiagonal sequence. This cannot establish
first-ever priority or exclude uncatalogued work. The package's
provenance file records source identifiers, hashes where available,
and inspection scopes without redistributing third-party full text.

The results proved here are an asymptotic expansion and integer
localization at each fixed order. They do not establish convergence
of the series, uniformity at growing truncation order, efficient
high-order computation, explicit numerical error constants, or exact
rounding for every input. The general algorithm is justified by the
finite reduction and the published fixed-rank theorem. No coefficient
beyond $c_3$ is stated or computationally certified in this package.
\begin{writenote}[6 October 2026]
``Report~156'' is Part~I. The direct antecedent is quoted in the front matter. The open items of this section are collected in Section~\ref{dgp:sec:further}, item~\ref{dgp:q:series}; the analytic inputs are external theorems whose statements, not proofs, the write read (note in Section~\ref{dgp:sub:sec:analytic}). The abstract's ``Zhou's published complete fixed-rank expansion'' is such an input.
\end{writenote}

\clearpage
\part{Exact exponential growth of increasing superdiagonal partitions}\label{dgp:rate:part}
\partsource{Bundle Report~158, archive \file{Superdiagonal_Partitions_Exact_Exponential_Growth_and_Inverse_Source.zip}, delivered as \file{Report158.tex} (13 pages). Delivered title block ``Exact exponential growth of increasing superdiagonal partitions / An elementary proof and quantitative threshold inversion / Report 158 / 3 October 2026''. Printed as delivered, with \textbf{[write]} notes; the proof of Lemma~\ref{dgp:rate:lem:dyck} is replaced by a pointer to Part~I. Delivered Section~$k$ is Section~$k+19$ here, and its statement or equation $k.j$ is $(k+19).j$.}
\begin{center}
\small
\begin{tabular}{@{}>{\raggedright\arraybackslash}p{0.46\textwidth} >{\raggedright\arraybackslash}p{0.46\textwidth}@{}}
\toprule
Reading conventions of this Part & Elsewhere in the report\\
\midrule
$b=\log2$; $C=\pi^2/12+b^2$; $B=2\sqrt C$; $L(x)=x\log4-x^2/2+I(x)$ & Part~I: $b=\pi/\sqrt3$, $B_*<B$; Part~V: $L=\log4$\\
$t$ the Boltzmann variable; $M$, $h$ prefix length and height; $D_h(M)$, $\rho_h$, $c_h$ Dyck data & Part~I: the same $D_h(M)$, $\rho_h$, $c_h$\\
$p$ a geometric parameter; $r=(1-p)/p$; $K_0,\dots,K_3$ constants; $\mathbb P$ a Boltzmann law & Parts~I, II: $p(n)$, $\mathbb P_n$ uniform\\
\bottomrule
\end{tabular}
\end{center}
\begin{partabstract}
Let $A(n)$ count weakly increasing positive-integer partitions
$\lambda_1\leq\cdots\leq\lambda_k$ of $n$ satisfying $\lambda_i\geq i$,
with $A(0)=1$. This is OEIS A238873. We prove
\[
 \log A(n)=B\sqrt n+O(n^{1/3}),\qquad
 B=\sqrt{\frac{\pi^2}{3}+4\log^2 2}=2.282910464597606\ldots.
\]
The upper error is in fact $O(1)$. For
$F(t)=\sum_{n\geq0}A(n)e^{-tn}$ and
$C=B^2/4=\pi^2/12+\log^2 2$, we establish
\[
 \frac Ct-O(t^{-1/3})\leq\log F(t)\leq\frac Ct+O(1).
\]
An exponential tilt gives the upper bound. A bounded-height Dyck prefix
and a geometric tail with an explicit positive survival probability give
the matching lower bound. An exact spectral inequality permits a growing
strip height. Chernoff bounds and an elementary monotonicity injection
then transfer the estimate to every coefficient without a Tauberian
black box. The threshold inverse satisfies
\[
 \min\{n\geq0:A(n)\geq y\}
 =\left(\frac{\log y}{B}\right)^2+O((\log y)^{5/3}).
\]
The result determines the exponential rate. It does not determine a
multiplicative equivalent, power prefactor, sharp second term, Airy
correction, or all-orders expansion. Classical Dyck representations and
the earlier lower bound in Report 156 are explicitly credited.
\end{partabstract}
\noindent\textbf{Scope and provenance.}
All logarithms are natural. The increasing order in the definition is
essential: the weakly decreasing superdiagonal convention counts a
different sequence. The proof is self-contained apart from elementary
analysis, the Basel sum, and finite-dimensional linear algebra, all
applications of which are made explicit. Finite exact computations are
separate implementation checks, not evidence from which an asymptotic
statement is inferred. A bounded primary-source search and report-catalogue
check found no equivalent exact-rate result; this is not an exhaustive
claim of scholarly priority. Section~\ref{dgp:rate:sec:sources} records the scope.
\medskip

\section{The partition model and main results}\label{dgp:rate:sec:main}
A partition here is a finite weakly increasing list of positive integers.
Its size is $|\lambda|=\sum_i\lambda_i$. Call it superdiagonal if
\begin{equation}\label{dgp:rate:eq:model}
 \lambda_i\geq i\qquad(1\leq i\leq k).
\end{equation}
The empty partition is superdiagonal and has size zero. Write $A(n)$ for
the number of such partitions of size $n$. The first values are
\[
 A(0),\ldots,A(13)=1,1,1,2,3,3,5,7,9,11,14,19,25,31.
\]
This is the class of Archibald, Blecher, Elizalde and
Knopfmacher~\cite{abek}, Section 2, and OEIS A238873~\cite{oeisA238873}.
Applying the same indexed condition to a decreasing list instead produces
A003114; it is not the model of this report.

Set
\begin{equation}\label{dgp:rate:eq:constants}
 b=\log2,\qquad C=\frac{\pi^2}{12}+b^2,
 \qquad B=2\sqrt C=\sqrt{\frac{\pi^2}{3}+4b^2}.
\end{equation}
For orientation, $C=1.3029200473423146\ldots$ and
$B=2.2829104645976064\ldots$. The exact expressions, not their decimals,
define these constants.

\begin{theorem}[Generating function and coefficient estimates]\label{dgp:rate:thm:main}
Let $F(t)=\sum_{n\geq0}A(n)e^{-tn}$ for $t>0$. There are positive
constants $K_0,K_1,K_2,K_3$ and positive thresholds such that
\begin{align}
 \frac Ct-K_0t^{-1/3}&\leq\log F(t)\leq\frac Ct+K_1
   &&\text{for all sufficiently small }t>0,\label{dgp:rate:eq:gfmain}\\
 B\sqrt n-K_2n^{1/3}&\leq\log A(n)\leq B\sqrt n+K_3
   &&\text{for all sufficiently large integers }n.\label{dgp:rate:eq:coefmain}
\end{align}
In particular,
\begin{equation}\label{dgp:rate:eq:limits}
 \lim_{t\downarrow0}t\log F(t)=C,\qquad
 \lim_{n\to\infty}\frac{\log A(n)}{\sqrt n}=B,
 \qquad A(n)=\exp\{B\sqrt n+O(n^{1/3})\}.
\end{equation}
\end{theorem}

\begin{theorem}[Threshold inverse]\label{dgp:rate:thm:inverse}
For real $y>1$, put $T(y)=\min\{n\geq0:A(n)\geq y\}$. This minimum
exists, and as $y\to\infty$,
\begin{equation}\label{dgp:rate:eq:inverse}
 \left(\frac{\log y}{B}\right)^2-O(\log y)
 \leq T(y)\leq
 \left(\frac{\log y}{B}\right)^2+O((\log y)^{5/3}).
\end{equation}
Thus $T(y)=(\log y/B)^2+O((\log y)^{5/3})$.
\end{theorem}

The constants and onset in these $O$-terms are unspecified. In particular,
\eqref{dgp:rate:eq:inverse} is not a finite-input certified search bracket or an
exact rounding rule. Likewise \eqref{dgp:rate:eq:coefmain} does not imply
$A(n)\sim c n^\alpha e^{B\sqrt n}$ for any specified $c$ or $\alpha$.
The $O(n^{1/3})$ error is a rigorous bound, not an identified correction.
\begin{writenote}[6 October 2026]
Theorem~\ref{dgp:rate:thm:main} \emph{replaces} the lower rate $B_*$ of Part~I's Theorem~\ref{dgp:conj:thm:super}. Part~IV then shows that the lower error is exactly of order $n^{1/6}$: $B\sqrt n-\log A(n)=Dn^{1/6}(1+o(1))$ with $D=\zeta bC^{-1/6}=1.5507\ldots$ (Theorem~\ref{dgp:airy:thm:airy}). So the $O(n^{1/3})$ in \eqref{dgp:rate:eq:coefmain} is a valid but not optimal bound, and no bound $\log A(n)\ge B\sqrt n-o(n^{1/6})$ holds; the upper bound $B\sqrt n+O(1)$ stands. The approach is slow: $\log A(n)/\sqrt n=1.5488$, $1.9121$, $2.1177$ at $n=100$, $800$, $8000$ (intake), against $B=2.2829$. Part~V gives, conditionally, the full equivalent. Theorem~\ref{dgp:rate:thm:inverse} is refined by Part~IV's Corollary~\ref{dgp:airy:cor:inverse} (unconditional) and Part~V's Corollary~\ref{dgp:pre:cor:inverse} (conditional); it is not an instance of the transseries volume's inversion theorems.
\end{writenote}

\paragraph{Relation to prior results.}
The 2025 paper~\cite{abek} develops generating functions and Dyck and
ballot descriptions of this class, proves its density among all partitions
is zero, and asks for its asymptotic enumeration. Report 156~\cite{r156}
proves the distinct-partition density conjecture using the lower constant
\begin{equation}\label{dgp:rate:eq:previous}
 B_* =\sqrt{\frac{\pi^2}{3}+2\log^2 2}
      =2.061740566010393\ldots < B.
\end{equation}
That report does not assert an exact exponential rate. The present report
extends the prefix construction by a surviving geometric tail and adds a
matching tilted upper bound. All steps needed for the new theorem are
given below; no part of Report 156 is modified or silently strengthened.
\begin{writenote}[6 October 2026]
``Report~156'' \cite{r156} is Part~I; \eqref{dgp:rate:eq:previous} is Part~I's \eqref{dgp:conj:eq:constants}, and the bounded-height prefix construction is Part~I's Sections~\ref{dgp:conj:sec:prefix}--\ref{dgp:conj:sec:height}.
\end{writenote}

\section{Cumulative multiplicities and an integral identity}\label{dgp:rate:sec:cumulative}
Let $Z_j$ be the multiplicity of the part $j$, and put
$C_j=\sum_{i=1}^jZ_i$. Only finitely many $Z_j$ are nonzero for a partition.

\begin{lemma}[Equivalent barrier]\label{dgp:rate:lem:barrier}
Condition~\eqref{dgp:rate:eq:model} holds if and only if
\begin{equation}\label{dgp:rate:eq:barrier}
 C_j\leq j\qquad\text{for every integer }j\geq1.
\end{equation}
\end{lemma}
\begin{proof}
If $C_j=k>0$, then $\lambda_k\leq j$, while superdiagonality gives
$k\leq\lambda_k$, so $C_j\leq j$. The case $C_j=0$ is immediate.
Conversely, if $\lambda_i<i$, the value $j=\lambda_i$ has
$C_j\geq i>j$, violating~\eqref{dgp:rate:eq:barrier}. The empty partition satisfies
both conditions.
\end{proof}
\begin{writenote}[6 October 2026]
Printed once here. Part~IV (Section~\ref{dgp:airy:sec:results}) and Part~V (\eqref{dgp:pre:eq:barrier}) use the same equivalence; Part~IV's proof of it is not reprinted.
\end{writenote}

The unrestricted partition product converges at $e^{-t}$ for every $t>0$:
\[
 \prod_{j\geq1}(1-e^{-tj})^{-1}<\infty.
\]
Indeed, $-\log(1-u)\leq2u$ for $0\leq u\leq1/2$, so the logarithms
have a summable tail, and the remaining finite set causes no difficulty.
This product bounds $F(t)$ from above. Every generating sum below is thus
finite or the monotone limit of finite nonnegative sums.

Define
\begin{equation}\label{dgp:rate:eq:fI}
 f(x)=-\log(1-e^{-x}),\qquad
 I(a)=\int_a^\infty f(x)\,dx
     =\sum_{r\geq1}\frac{e^{-ar}}{r^2}=\Li_2(e^{-a}).
\end{equation}
The nonnegative series for $f$ justifies termwise integration. As $f$ is
decreasing, for every integer $M\geq0$ and $x=tM>0$,
\begin{equation}\label{dgp:rate:eq:riemann}
 \frac{I(x+t)}t\leq\sum_{j>M}f(tj)\leq\frac{I(x)}t.
\end{equation}
If $x$ stays in a fixed compact subinterval of $(0,\infty)$, either bound
differs from $I(x)/t$ by $O(1)$, since
$0\leq I(x)-I(x+t)\leq t f(x)$.

For later use, put
\begin{equation}\label{dgp:rate:eq:L}
 L(x)=x\log4-\frac{x^2}{2}+I(x).
\end{equation}
Then
\begin{equation}\label{dgp:rate:eq:Lb}
 L(b)=\frac32b^2+\Li_2(1/2)
     =\frac{\pi^2}{12}+b^2=C.
\end{equation}
Here the standard dilogarithm value needs only elementary calculus:
the derivative of
$\Li_2(u)+\Li_2(1-u)+\log u\log(1-u)$ vanishes on $(0,1)$.
The limit at zero is $\sum_{r\geq1}r^{-2}=\pi^2/6$; setting $u=1/2$
gives $\Li_2(1/2)=\pi^2/12-b^2/2$.
\begin{writenote}[6 October 2026]
This evaluation and the function $L$ are used again in Part~IV (\eqref{dgp:airy:eq:factor}, where $L(x)=2bx-x^2/2+I(x)$ is the same function) and in Part~V (\eqref{dgp:pre:eq:normalizer}); Part~IV's proof of the evaluation is not reprinted.
\end{writenote}

\section{A matching exponential tilt upper bound}\label{dgp:rate:sec:upper}
For sufficiently small $t>0$, take $M=\floor{b/t}\geq1$ and define
$\eta_j=b-tj$ for $1\leq j\leq M$. These multipliers are nonnegative
and nonincreasing. Summation by parts and~\eqref{dgp:rate:eq:barrier} give
\begin{align}
 \sum_{j=1}^M\eta_jZ_j
 &=\eta_MC_M+\sum_{j=1}^{M-1}(\eta_j-\eta_{j+1})C_j\notag\\
 &\leq\eta_MM+\sum_{j=1}^{M-1}(\eta_j-\eta_{j+1})j
   =\sum_{j=1}^M\eta_j.\label{dgp:rate:eq:tiltbarrier}
\end{align}
The signs of all coefficients matter: the cumulative upper bounds can be
substituted in precisely this direction.

Since $e^{-tjZ_j}=2^{-Z_j}e^{\eta_jZ_j}$ for $j\leq M$, bound the
exponential factor with~\eqref{dgp:rate:eq:tiltbarrier} and then discard every
remaining constraint. Summing the independent geometric series yields
\begin{equation}\label{dgp:rate:eq:exactupper}
 F(t)\leq
 \exp\!\left\{\sum_{j=1}^M(\log4-tj)\right\}
 \prod_{j>M}(1-e^{-tj})^{-1}.
\end{equation}
For the first $M$ variables, the extra $M\log2$ comes from
$\prod_{j=1}^M\sum_{z\geq0}2^{-z}=2^M$; this accounts for the factor
$\log4$ rather than $\log2$.

With $x=tM=b+O(t)$, equations~\eqref{dgp:rate:eq:riemann} and
\eqref{dgp:rate:eq:exactupper} imply
\begin{align}
 \log F(t)&\leq M\log4-\frac{tM(M+1)}2+\sum_{j>M}f(tj)\notag\\
 &\leq\frac{L(x)}t-\frac x2
  =\frac Ct+O(1).\label{dgp:rate:eq:upperfinal}
\end{align}
The last step uses smoothness of $L$ near $b$ and~\eqref{dgp:rate:eq:Lb}.
For completeness,
\[
 L'(x)=\log\!\bigl(4e^{-x}(1-e^{-x})\bigr),\qquad
 L'(b)=L''(b)=0,\quad L'''(b)=-2.
\]
These cancellations improve the cutoff discrepancy, but they are not
needed for the asserted upper bound.

\section{Bounded height prefixes and an exact spectral bound}\label{dgp:rate:sec:prefix}
For positive integers $M,h$, let $\mathcal P_{M,h}$ consist of weakly
increasing lists of exactly $M$ parts satisfying
\begin{equation}\label{dgp:rate:eq:prefix}
 i\leq\lambda_i\leq\min(M,i+h-1)\qquad(1\leq i\leq M).
\end{equation}
Every such list has all parts at most $M$ and weight at most
\begin{equation}\label{dgp:rate:eq:weight}
 W_h(M)=\frac{M(M+1)}2+(h-1)M.
\end{equation}
Write $D_h(M)=|\mathcal P_{M,h}|$.

\begin{lemma}[Dyck representation]\label{dgp:rate:lem:dyck}
$D_h(M)$ is the number of Dyck paths of semilength $M$ whose maximum
height is at most $h$.
\end{lemma}
\begin{proof}
\begin{writenote}[6 October 2026]
Report~158's proof is not reprinted: it is the argument of Part~I, step by step. Part~I's Section~\ref{dgp:conj:sec:prefix} maps a list to its reverse complement $\mu_j=M+1-\lambda_{M+1-j}$, a weakly increasing list with $1\le\mu_j\le j$, and places the $j$th north step of a path from $(0,0)$ to $(M,M)$ at horizontal coordinate $\mu_j-1$, which gives the Dyck paths of semilength $M$ bijectively (the inverse reads the coordinates). Part~I's Section~\ref{dgp:conj:sec:height} adds the height: the upper bound in \eqref{dgp:rate:eq:prefix} becomes $\mu_j\ge j-h+1$, the height just after the $j$th north step is $j-\mu_j+1$, between $1$ and $h$, and maxima occur just after north steps. The statement above is Report~158's.
\end{writenote}
\end{proof}

This is a specialized form of the classical path encodings discussed
in~\cite{abek}; those encodings and Catalan enumeration are prior inputs.
The bounded-height form was also used in Report 156~\cite{r156}.
We include the spectral calculation to specify exactly what remains
uniform when the height grows.

Let $T_h$ be the adjacency matrix of the path on vertices $0,1,\ldots,h$.
Reading Dyck heights gives
\begin{equation}\label{dgp:rate:eq:matrix}
 D_h(M)=(T_h^{2M})_{0,0}.
\end{equation}
Its orthonormal eigenvectors have coordinates
\[
 v_r(a)=\sqrt{\frac2{h+2}}\sin\frac{(a+1)r\pi}{h+2}
 \quad(0\leq a\leq h),\qquad 1\leq r\leq h+1,
\]
with eigenvalues $2\cos(r\pi/(h+2))$. The sine addition formula verifies
the eigenvalue relation, including the zero boundary coordinates;
finite sine orthogonality gives the normalization. Thus
\begin{equation}\label{dgp:rate:eq:spectral}
 D_h(M)=\frac2{h+2}\sum_{r=1}^{h+1}
 \sin^2\frac{r\pi}{h+2}
 \left(2\cos\frac{r\pi}{h+2}\right)^{2M}.
\end{equation}
Set
\begin{equation}\label{dgp:rate:eq:rho}
 \rho_h=4\cos^2\frac\pi{h+2},\qquad
 c_h=\frac4{h+2}\sin^2\frac\pi{h+2}.
\end{equation}
All summands in~\eqref{dgp:rate:eq:spectral} are nonnegative. Keeping the two
extreme eigenvalues proves the exact, uniform inequality
\begin{equation}\label{dgp:rate:eq:spectralbound}
 D_h(M)\geq c_h\rho_h^M\qquad(M,h\geq1).
\end{equation}
For fixed $h$, the remaining eigenvalues have smaller modulus, so
\begin{equation}\label{dgp:rate:eq:fixedh}
 \log D_h(M)=M\log\rho_h+O_h(1)\quad(M\to\infty),\qquad
 \rho_h\longrightarrow4\quad(h\to\infty).
\end{equation}
When $h=1$, the two extreme eigenvalues exhaust the spectrum and the
formula is already exact. We will use~\eqref{dgp:rate:eq:fixedh} only with fixed
height, and~\eqref{dgp:rate:eq:spectralbound} for growing height.

\section{A geometric tail with an explicit survival bound}\label{dgp:rate:sec:tail}
\begin{lemma}[Uniform survival]\label{dgp:rate:lem:survival}
Let $Y_1,Y_2,\ldots$ be independent with
$\PP(Y_i=z)=(1-p)p^z$ for $z=0,1,\ldots$, where $0<p<1/2$. Then
\begin{equation}\label{dgp:rate:eq:survival}
 \PP\left(\sum_{i=1}^jY_i\leq j\text{ for every }j\geq1\right)
 \geq1-2p>0.
\end{equation}
\end{lemma}
\begin{proof}
Put $S_n=\sum_{i=1}^n(Y_i-1)$, $S_0=0$, and $r=(1-p)/p>1$. Direct
summation gives
\[
 \EE r^{Y_i-1}=\frac{1-p}{r(1-pr)}=1.
\]
Therefore $r^{S_n}$ is a nonnegative mean-one martingale. If $\tau$ is
the first time $S_n\geq1$, finite-time stopping gives
$1=\EE r^{S_{\tau\wedge N}}\geq r\PP(\tau\leq N)$.
This finite identity also follows by expanding the stopped increments:
each increment multiplied by $\mathbf1_{\{\tau>j\}}$ has conditional
mean zero. Hence no infinite-time optional-stopping assertion is needed.
Letting $N\to\infty$ in the increasing events yields
\[
 \PP(S_n<1\text{ for every }n\geq1)\geq1-1/r.
\]
Now require $Y_1=0$ and impose this last event on the independent walk
formed from $Y_2,Y_3,\ldots$. Its partial sums are integers below $1$,
hence at most zero. The original partial sums are consequently at most
$-1$ from time one onward. This is a subset of the event in
\eqref{dgp:rate:eq:survival}, with probability at least
$(1-p)(1-1/r)=1-2p$.
\end{proof}
\begin{writenote}[6 October 2026]
Printed once here. Part~IV restates it as \eqref{dgp:airy:eq:martingale}--\eqref{dgp:airy:eq:survival} (its proof there is not reprinted) and uses the same martingale for the ruin bound \eqref{dgp:airy:eq:ruin}; Part~V uses it inside the proof of Lemma~\ref{dgp:pre:lem:cutoff}.
\end{writenote}

Fix $a>b$ and put $M=\floor{a/t}$. Consider independent tail
multiplicities $Z_{M+j}$ with parameters
$p_j=e^{-t(M+j)}$, $j\geq1$. The strict endpoint inequality
$t(M+1)>a$ implies
\[
 p_j<e^{-a}=p<1/2.
\]
A coordinatewise coupling gives $Z_{M+j}\leq Y_j$ for the iid geometrics
in Lemma~\ref{dgp:rate:lem:survival}: their tail probabilities are $p_j^k\leq p^k$
for every $k$. Since survival is decreasing in each multiplicity,
\begin{equation}\label{dgp:rate:eq:tailprob}
 \PP\left(\sum_{i=1}^j Z_{M+i}\leq j\text{ for every }j\geq1\right)
 \geq1-2e^{-a}.
\end{equation}
The product measure is supported on finite partitions. In fact
$\sum_{j>M}\EE Z_j=\sum_{j>M}e^{-tj}/(1-e^{-tj})<\infty$, which implies
a finite total number of parts almost surely.

Let $G_M(t)$ sum the weights $e^{-t|\lambda|}$ of tail partitions with
all parts greater than $M$ satisfying the barrier in~\eqref{dgp:rate:eq:tailprob}.
Normalizing the geometric product gives the precise inequality
\begin{equation}\label{dgp:rate:eq:taillower}
 G_M(t)\geq(1-2e^{-a})\prod_{j>M}(1-e^{-tj})^{-1}.
\end{equation}
The normalizing product is finite and positive. This argument may also
be obtained from finite products and monotone limits; it introduces no
infinite formal measure with a missing normalization.

\section{Concatenation and the exact generating function rate}\label{dgp:rate:sec:concat}
Append a tail counted by $G_M(t)$ to a prefix in $\mathcal P_{M,h}$.
The prefix has exactly $M$ parts, all at most $M$. For a value $M+j$,
the cumulative number of parts of the concatenated partition is at most
$M+j$ by the tail barrier; below $M$, the prefix already satisfies
Lemma~\ref{dgp:rate:lem:barrier}. The result is superdiagonal.
The cutoff at $M$ uniquely recovers both pieces, so this construction is
injective even when different prefixes have the same weight.

By~\eqref{dgp:rate:eq:weight} and~\eqref{dgp:rate:eq:taillower}, for $M=\floor{a/t}$,
\begin{equation}\label{dgp:rate:eq:concat}
 F(t)\geq D_h(M)e^{-tW_h(M)}(1-2e^{-a})
          \prod_{j>M}(1-e^{-tj})^{-1}.
\end{equation}
Keep $a>b$ and $h$ fixed while $t\downarrow0$. Multiplying logarithms by
$t$, and using $tM\to a$, \eqref{dgp:rate:eq:riemann} and~\eqref{dgp:rate:eq:fixedh}, gives
\begin{equation}\label{dgp:rate:eq:liminffixed}
 \liminf_{t\downarrow0}t\log F(t)
 \geq a\log\rho_h-\frac{a^2}2+I(a).
\end{equation}
The prefix excess weight contributes $t^2(h-1)M\to0$, and the survival
factor contributes a fixed logarithm times $t$, also tending to zero.
First let fixed $h\to\infty$, and then let $a\downarrow b$. The right
side tends to $L(b)=C$. Together with~\eqref{dgp:rate:eq:upperfinal}, this proves
\begin{equation}\label{dgp:rate:eq:gflimit}
 t\log F(t)\longrightarrow C.
\end{equation}
There is no interchange of a nonuniform pair of limits: each fixed
$(a,h)$ bounds the same liminf from below, and the supremum of these valid
bounds is at least $C$.

\section{Elementary transfer to every coefficient}\label{dgp:rate:sec:transfer}
\begin{lemma}[Monotonicity]\label{dgp:rate:lem:mono}
$A(n+1)\geq A(n)$ for every integer $n\geq0$.
\end{lemma}
\begin{proof}
For a nonempty partition increase one largest part by $1$. Its position
at the end remains valid, the indexed lower bounds are preserved, and
the resulting largest part is unique. Decreasing that unique largest
part recovers the original list on the image. Thus this map is injective.
For $n=0$, send the empty list to $(1)$.
\end{proof}
\begin{writenote}[6 October 2026]
Printed once here; Part~IV (Section~\ref{dgp:airy:sec:pointwise}) and Part~V (Section~\ref{dgp:pre:sec:inverse}) use the same injection, and their proofs of it are not reprinted. The inequality is not strict in general: $A(0)=A(1)=A(2)=1$ and $A(4)=A(5)=3$.
\end{writenote}

Define a random nonnegative integer $N_t$ by
\begin{equation}\label{dgp:rate:eq:Boltzmann}
 \PP(N_t=n)=\frac{A(n)e^{-tn}}{F(t)}.
\end{equation}
For now assume only~\eqref{dgp:rate:eq:gflimit}. Fix $0<\epsilon<1$ and choose
a small fixed $s\in(0,1)$. Exponential Markov bounds yield
\begin{align}
 \PP\left(N_t\geq\frac{(1+\epsilon)C}{t^2}\right)
 &\leq e^{-s(1+\epsilon)C/t}\frac{F((1-s)t)}{F(t)},\label{dgp:rate:eq:upperchernoff}\\
 \PP\left(N_t\leq\frac{(1-\epsilon)C}{t^2}\right)
 &\leq e^{s(1-\epsilon)C/t}\frac{F((1+s)t)}{F(t)}.\label{dgp:rate:eq:lowerchernoff}
\end{align}
Using~\eqref{dgp:rate:eq:gflimit}, their logarithms are at most, respectively,
\begin{align*}
 \frac{Cs}{t}\left(\frac{s}{1-s}-\epsilon\right)+o(1/t),\qquad
 \frac{Cs}{t}\left(\frac{s}{1+s}-\epsilon\right)+o(1/t).
\end{align*}
The choice $s=\epsilon/4$ makes both leading coefficients negative.
Thus the interval between $(1-\epsilon)C/t^2$ and
$(1+\epsilon)C/t^2$ has probability tending to one exponentially fast.
Real endpoints cause no difficulty: one sums the integer indices that
lie between them.

For an arbitrary sufficiently large integer $n$, set
$t=\sqrt{(1+\epsilon)C/n}$. The upper endpoint equals $n$. Monotonicity
and the lower endpoint bound on $N_t$ give
\begin{equation}\label{dgp:rate:eq:centralbound}
 (1-o(1))F(t)
 \leq\sum_{(1-\epsilon)C/t^2\leq m\leq n}A(m)e^{-tm}
 \leq(n+1)A(n)e^{-(1-\epsilon)C/t}.
\end{equation}
After taking logarithms,
\[
 \log A(n)\geq\frac{(2-\epsilon)C}{t}+o(1/t),\qquad
 \liminf_{n\to\infty}\frac{\log A(n)}{\sqrt n}
 \geq(2-\epsilon)\sqrt{\frac{C}{1+\epsilon}}.
\]
Let $\epsilon\downarrow0$. Conversely $A(n)e^{-tn}\leq F(t)$, and
$t=\sqrt{C/n}$ with~\eqref{dgp:rate:eq:upperfinal} gives
\begin{equation}\label{dgp:rate:eq:coefupper}
 \log A(n)\leq B\sqrt n+O(1).
\end{equation}
This proves the coefficient limit in~\eqref{dgp:rate:eq:limits}, pointwise for all
sufficiently large indices. No derivative estimate, local limit theorem,
complex-analytic saddle-point argument, or unstated Tauberian theorem is
being used.

\section{Quantitative errors and inverse localization}\label{dgp:rate:sec:errors}
\subsection{Growing height with a uniform lower bound}
For sufficiently small $t>0$, choose
\begin{equation}\label{dgp:rate:eq:growing}
 a=b+t,\qquad M=\floor{a/t}=\floor{b/t}+1,\qquad
 h=\ceil{t^{-1/3}},\qquad x=tM.
\end{equation}
Then $b<x\leq b+t$. The survival factor in~\eqref{dgp:rate:eq:concat} satisfies
\begin{equation}\label{dgp:rate:eq:delta}
 1-2e^{-a}=1-e^{-t}\geq t/2\qquad(0<t\leq1).
\end{equation}
For $h\geq2$, elementary trigonometric estimates imply
\begin{equation}\label{dgp:rate:eq:uniformspectral}
 c_h\geq\frac{16}{(h+2)^3},\qquad
 0\leq\log4-\log\rho_h\leq\frac{2\pi^2}{(h+2)^2}.
\end{equation}
Indeed, $\sin u\geq2u/\pi$ on $[0,\pi/2]$ gives the first bound.
For the second use $-\log\cos u\leq u^2$ on $[0,\pi/4]$, which follows
by integrating $\tan u\leq2u$ there. These inequalities apply uniformly
to the actual height in~\eqref{dgp:rate:eq:growing}.

Taking logarithms of~\eqref{dgp:rate:eq:concat}, using the exact inequality
\eqref{dgp:rate:eq:spectralbound}, and applying the lower Riemann bound yields
\begin{align}
 \log F(t)\geq{}&\frac{L(x)}t-
    O\!\left(\frac{M}{(h+2)^2}\right)-O(h)
    -O(\log(1/t))-O(1).\label{dgp:rate:eq:gflowerdetail}
\end{align}
Here $\log c_h$ costs $O(\log(h+2))$, the factor in~\eqref{dgp:rate:eq:delta}
costs $O(\log(1/t))$, and the prefix weight beyond $x^2/(2t)$ costs
$x/2+(h-1)x=O(h)$. Finally,
$I(x+t)/t\geq I(x)/t-f(x)$, with $f(x)=O(1)$ near $b$.
This accounts separately for every loss in~\eqref{dgp:rate:eq:gflowerdetail}.

Since $x=b+O(t)$, $L(x)/t=C/t+O(1)$. Moreover
$M/(h+2)^2=O(t^{-1/3})$, $h=O(t^{-1/3})$, and
$\log(1/t)=O(t^{-1/3})$. Thus
\begin{equation}\label{dgp:rate:eq:gferror}
 \log F(t)\geq C/t-O(t^{-1/3}).
\end{equation}
Together with~\eqref{dgp:rate:eq:upperfinal}, this proves~\eqref{dgp:rate:eq:gfmain}.
The argument uses an exact all-height inequality, not the fixed-height
asymptotic~\eqref{dgp:rate:eq:fixedh} with an uncontrolled height-dependent error.

\subsection{Quantitative coefficient transfer}
Write $\log F(t)=C/t+E(t)$, where
\begin{equation}\label{dgp:rate:eq:E}
 -Kt^{-1/3}\leq E(t)\leq K_1
\end{equation}
for all sufficiently small positive $t$. Take
$\epsilon=L_0t^{1/3}$ and $s=\epsilon/4$, with a fixed sufficiently
large $L_0>0$. For $\epsilon\leq1/2$, the leading terms of the logarithms
in~\eqref{dgp:rate:eq:upperchernoff} and~\eqref{dgp:rate:eq:lowerchernoff} are each at most
\[
 -\frac{C\epsilon^2}{8t}.
\]
For example, $s/(1-s)\leq\epsilon/2$ in this range, and multiplication
by $Cs/t=C\epsilon/(4t)$ gives the bound; the other sign is no larger.
The error in either logarithm is at most
\[
 E((1\pm s)t)-E(t)\leq Kt^{-1/3}+K_1.
\]
Choose $L_0$ so that $CL_0^2/8>K+1$. Then both tail probabilities are
at most $\exp(-ct^{-1/3})$ for some $c>0$ and all sufficiently small $t$.
The central probability is consequently bounded below by $1/2$ for
all sufficiently small $t$.

For each sufficiently large integer $n$, choose the positive solution of
\begin{equation}\label{dgp:rate:eq:tchoice}
 n=(1+L_0t^{1/3})C/t^2.
\end{equation}
Existence and uniqueness follow because the right side is the strictly
decreasing continuous function $Ct^{-2}+L_0Ct^{-5/3}$ with range
$(0,\infty)$. This solution tends to zero and satisfies
$t\asymp n^{-1/2}$. Apply the central-sum argument
\eqref{dgp:rate:eq:centralbound}, now with its probability at least $1/2$, to get
\begin{equation}\label{dgp:rate:eq:coeflowerinter}
 \log A(n)\geq\frac{(2-\epsilon)C}{t}
          -Kt^{-1/3}-\log(2(n+1)).
\end{equation}
Since $B\sqrt n=2C\sqrt{1+\epsilon}/t$ and
$\sqrt{1+\epsilon}\leq1+\epsilon/2$, the difference between this target
and the lower bound is
\[
 O(\epsilon/t+t^{-1/3}+\log n)
 =O(t^{-2/3})=O(n^{1/3}).
\]
This proves the lower half of~\eqref{dgp:rate:eq:coefmain}; the upper half is
\eqref{dgp:rate:eq:coefupper}. Theorem~\ref{dgp:rate:thm:main} is complete.

\subsection{Inverting the threshold}
By Lemma~\ref{dgp:rate:lem:mono} and the lower bound in~\eqref{dgp:rate:eq:coefmain},
$A(n)$ is nondecreasing and tends to infinity, so $T(y)$ exists. Put
$u=\log y$ and $x=(u/B)^2$. For large $y$, any index with $A(n)\geq y$
lies past the finite onset of the asymptotic estimates. The upper bound
$\log A(n)\leq B\sqrt n+K_3$ therefore gives
\begin{equation}\label{dgp:rate:eq:invlower}
 T(y)\geq\left(\frac{u-K_3}{B}\right)^2=x-O(u).
\end{equation}
For the other direction, choose a fixed $D>0$ and set
$n=\ceil{x+D u^{5/3}}$. Elementary Taylor expansion gives
\begin{align}
 B\sqrt n&=u+\frac{DB^2}{2}u^{2/3}+O(u^{1/3}),\label{dgp:rate:eq:invtaylor1}\\
 n^{1/3}&=B^{-2/3}u^{2/3}+O(u^{1/3}).\label{dgp:rate:eq:invtaylor2}
\end{align}
The ceiling contributes a still smaller error. Choose
$DB^2/2>K_2B^{-2/3}$. The lower bound in~\eqref{dgp:rate:eq:coefmain} and these
expansions then imply $\log A(n)\geq u$ for every sufficiently large $y$.
Consequently $T(y)\leq x+O(u^{5/3})$. Together with~\eqref{dgp:rate:eq:invlower}
this proves Theorem~\ref{dgp:rate:thm:inverse}. The permitted error grows, so no
adjacent-integer localization is asserted.

\section{Exact finite verification and reproducibility}\label{dgp:rate:sec:checks}
The standard-library companion supplied with this report performs exact
integer and rational checks. Its finite limits and output-safety contract
are specified in the package README. It is useful to record why the
algorithms correspond to the objects and inequalities in the proof.

\subsection{Capped multiplicity recurrence}
Fix a maximum weight $N$. After processing part values through $v$, let
$d_v(w,k)$ count multiplicity vectors of weight $w$ with $k$ parts,
satisfying every cumulative barrier so far. Starting from
$d_0(0,0)=1$ and zero elsewhere, the fresh-layer recurrence is
\begin{equation}\label{dgp:rate:eq:dp}
 d_v(w,k)=\mathbf1_{\{k\leq v\}}
 \sum_{z=0}^k d_{v-1}(w-vz,k-z),
\end{equation}
where entries with negative weight are zero. Choosing $z$ specifies the
multiplicity of the new value $v$ uniquely; the previous layer enforces
all earlier barriers, and $k\leq v$ enforces the new one. Therefore
$A(w)=\sum_kd_N(w,k)$ for $0\leq w\leq N$. No part larger than $N$ can
occur in a partition of these weights. The bound
$k(k+1)/2\leq w\leq N$ permits the finite cap
$k\leq\floor{(\sqrt{8N+1}-1)/2}$, evaluated with integer arithmetic.

For a tail cutoff $M$, process values $v>M$ and replace the barrier by
$k\leq v-M$. The same inductive proof counts exactly the tails in
Section~\ref{dgp:rate:sec:tail}, with weight still measured using their actual
part sizes. The stronger lower bound on a $k$-part tail is
$Mk+k(k+1)/2$; using the preceding ordinary cap merely retains extra zero
states and remains valid.

Direct enumeration of ascending partitions is independent of this
recurrence. For each list it compares the indexed and cumulative tests,
and checks that the largest-part injection is valid and collision-free.
Prefix enumeration independently compares its count with an exact
walk recurrence for $(T_h^{2M})_{0,0}$, checks the weight cap, and checks
the coefficientwise prefix-tail inequalities.

\subsection{Exact rational probability and product checks}
At a finite horizon $H$, retain states
$d=j-\sum_{i=1}^j Z_i\geq0$. From a state $d$, a new geometric
multiplicity $z$ is admissible precisely when $0\leq z\leq d+1$.
The transition has exact weight $(1-p)p^z$, and the new state is
$d+1-z$. Thus rational dynamic programming computes the finite survival
probability exactly when $p$ is rational. The companion compares it with
$1-2p$ and verifies the martingale normalization algebraically. These
finite-horizon tests do not themselves imply infinite survival; that
claim is proved in Lemma~\ref{dgp:rate:lem:survival}.

For rational $q\in(0,1)$ and a finite largest part $J$, another finite
multiplicity recurrence computes the total rational generating weight.
Choose $M$ such that $q^j\geq1/2$ for $j\leq M$ and
$q^{M+1}<1/2$. The finite version of~\eqref{dgp:rate:eq:exactupper} is
\begin{equation}\label{dgp:rate:eq:finiteupper}
 \sum_{\substack{\lambda\ \mathrm{superdiagonal}\\\lambda_k\leq J}}
 q^{|\lambda|}
 \leq4^M q^{M(M+1)/2}\prod_{j=M+1}^J(1-q^j)^{-1}.
\end{equation}
The sum includes the empty partition. For tails with cutoff $M$, the
corresponding lower bound is
\begin{equation}\label{dgp:rate:eq:finitelower}
 \sum_{\substack{\lambda\ \mathrm{tail\ admissible}\\\lambda_k\leq J}}
 q^{|\lambda|}
 \geq (1-2q^{M+1})\prod_{j=M+1}^J(1-q^j)^{-1}.
\end{equation}
It follows either by the same finite coupling or by applying the
infinite iid survival event to its first $J-M$ coordinates. The companion
compares these rational quantities exactly; decimal approximations are
not used for pass/fail decisions.

\subsection{Numerical values and scope}
The exact count table contains $A(n)$ for $0\leq n\leq200$. Selected
values are
\begin{center}
\begin{tabular}{rrrrrr}
\toprule
$n$&$20$&$40$&$60$&$100$&$200$\\
\midrule
$A(n)$&$135$&$3997$&$61221$&$5323937$&$25548856857$\\
\bottomrule
\end{tabular}
\end{center}
The attributed OEIS numerical prefix is compared with the computed table.
The package also supplies a bounded exact threshold command. If its
search bound does not reach a requested threshold, it reports that fact;
it does not substitute the asymptotic model for an exact answer.

Finite checks verify the selected integer and rational comparisons only.
They establish neither the theorem's asymptotic onset nor a sharp error
constant, and cannot prove a claim of priority. The proof of the theorem
is Sections~\ref{dgp:rate:sec:cumulative}--\ref{dgp:rate:sec:errors}. The source package
includes tests that remain effective under Python optimization, a clean
PDF builder, and fixed-allowlist deterministic archive construction.
Reproducibility is byte-for-byte within the recorded toolchain, not a
promise across different TeX versions or platforms.
\begin{writenote}[6 October 2026]
The package is shipped with the prefix \file{158-growth-} (code under \file{code/}; the two result files, the OEIS prefix and the provenance record under \file{data/}; the PDF, README and \file{SHA256SUMS} are not shipped). At intake, on a Windows copy, the exact data were regenerated byte for byte; four of the 40 companion tests fail there, all of them CLI or POSIX safe-output tests. The threshold $T(10^6)=84$ was reproduced at the write.
\end{writenote}

\section{Source assessment and limits of the result}\label{dgp:rate:sec:sources}
The full final 17-page 2025 article~\cite{abek} was inspected, including
its path preliminaries, Section 2 and the entire asymptotic Section 5.
Section 2 identifies the present increasing class and supplies exact
multivariate generating functions. Proposition 5.1 shows its density
among unrestricted partitions tends to zero. Conjecture 5.2 asks whether
distinct partitions have density zero within the superdiagonal class;
the closing discussion asks for asymptotic enumeration. The exact
constant~\eqref{dgp:rate:eq:constants} is not supplied there. The article's Dyck
and ballot encodings are credited rather than claimed as new.

The official OEIS A238873 page was checked on 3 October 2026. It states
the indexed condition, distinguishes the decreasing class, gives the
prefix through $n=60$, and cites~\cite{abek}; no exact exponential rate
was displayed. The author's primary publication list~\cite{elizalde},
including 2026 entries, and targeted searches for the sequence, the
partition phrase, logarithmic asymptotics, $\log2$, and the decimal rate
found no equivalent later result. Superdiagonal compositions are a
different ordered class and were not treated as evidence for this model.

The available report catalogues were refreshed on 3 October 2026 at
08:39 UTC; the relevant files were unchanged from the earlier inspection,
and the targeted repository searches found no matching result.
Report 156 is the explicitly distinguished antecedent~\eqref{dgp:rate:eq:previous}.
This bounded check supports proceeding with a separate exact-rate report,
but it does not establish exhaustive worldwide priority. The source
provenance file records hashes and inspection scopes without redistributing
full journal PDFs or full OEIS entry text.

The present conclusions are precisely the logarithmic estimates
\eqref{dgp:rate:eq:gfmain}--\eqref{dgp:rate:eq:coefmain} and the inverse localization
\eqref{dgp:rate:eq:inverse}. Determining a sharp second term, an amplitude, a
power prefactor, or an all-orders expansion remains outside these proofs.
No optimality claim is made for the error $O(n^{1/3})$. The growing-height
argument is only a sufficient lower-bound construction; its loss is not
identified with the true next-order asymptotics.
\begin{writenote}[6 October 2026]
\emph{Answered by Part~IV:} the next-order term is $-Dn^{1/6}$ (Theorem~\ref{dgp:airy:thm:airy}), so the error order $n^{1/3}$ is not optimal (Section~\ref{dgp:sec:further}, ``Answered inside this report''). The repository search of the preceding paragraph is still accurate: at the write (6~October 2026) no other report treats this sequence. A sharp amplitude, a power prefactor and all orders: Section~\ref{dgp:sec:further}, items~\ref{dgp:q:amplitude}--\ref{dgp:q:allorders}.
\end{writenote}

\clearpage
\part{Airy correction and limit shape for superdiagonal partitions}\label{dgp:airy:part}
\partsource{Bundle Report~160, archive \file{Superdiagonal_Partitions_Airy_Correction_and_Limit_Shape_Source.zip}, delivered as \file{Report160.tex} (17 pages). Delivered title block ``Airy correction and limit shape for superdiagonal partitions / Pointwise enumeration and macroscopic structure / Report 160 / 3 October 2026''. Printed as delivered, with \textbf{[write]} notes; four restated proofs (barrier, dilogarithm, survival, monotonicity) are replaced by pointers to Part~III. Delivered Section~$k$ is Section~$k+29$ here, and its statement or equation $k.j$ is $(k+29).j$.}
\begin{center}
\small
\begin{tabular}{@{}>{\raggedright\arraybackslash}p{0.46\textwidth} >{\raggedright\arraybackslash}p{0.46\textwidth}@{}}
\toprule
Reading conventions of this Part & Elsewhere in the report\\
\midrule
$-\zeta$ the largest zero of $\operatorname{Ai}$, $\zeta\approx2.3381>0$ & Part~V: $a=a_1=-\zeta$; Li: $\zeta$ an argument\\
$K$ the number of parts; $g$, $f$ the limit profiles; $h(x)=-\log(1-e^{-x})$; $m(x)$ mean profile & Part~V: $K$ an amplitude; Part~III: $f(x)=-\log(1-e^{-x})$\\
$\beta=d-2b-\log(1-e^{-d})$; $D_M(t)$, $D=\zeta bC^{-1/6}$; $\mathbb P_n$ uniform on superdiagonal partitions & Parts~I, II: $\beta=c/\sqrt n$; Part~V: the same $D$\\
\bottomrule
\end{tabular}
\end{center}
\begin{partabstract}
Let $A(n)$ count weakly increasing partitions
$1\leq\lambda_1\leq\cdots\leq\lambda_K$ of $n$ with
$\lambda_i\geq i$. Put $b=\log2$, $C=\pi^2/12+b^2$, and let
$-\zeta$ be the largest negative zero of the Airy function $\operatorname{Ai}$.
We prove the pointwise expansion
\[
 \log A(n)=2\sqrt{Cn}-\zeta b C^{-1/6}n^{1/6}+o(n^{1/6}).
\]
Its radial analogue is
$\log\sum_n A(n)e^{-tn}=C/t-\zeta b t^{-1/3}+o(t^{-1/3})$.
Mallein's published killed-walk area estimates supply the Airy input.
A separate exact-weight argument proves the coefficient lower bound:
it splits the geometric tail twice, creates a macroscopic buffer,
proves a uniform lattice local limit on all Fourier arcs, retilts at
a fixed cutoff, and subtracts the remaining ruin probability.
The ordinary threshold inverse has an explicit second term of order
$(\log y)^{4/3}$.

For a uniform partition of $n$, at scale $t=\sqrt{C/n}$, the scaled length $tK$
converges to $2b$. The cumulative multiplicity profile converges
uniformly on the whole nonnegative half-line to
$g(x)=x$ for $x\leq b$ and
$g(x)=2b+\log(1-e^{-x})$ for $x\geq b$.
Its inverse gives the row profile on compact subsets of $(0,2b)$.
All fixed-tolerance deviations have probability at most $e^{-c/t}$.
The initial half of the limiting row profile follows the diagonal;
this is not a statement about exact discrete contacts.
No power prefactor, multiplicative equivalent, further asymptotic
term, or fluctuation theorem is asserted.
\end{partabstract}
\noindent\textbf{Attribution and position.}
The class is OEIS A238873~\cite{oeisA238873}, studied by Archibald, Blecher,
Elizalde and Knopfmacher~\cite{abek}. Their generating functions and
ballot/Dyck encodings are prior results. Report~158~\cite{r158}
established the exact exponential rate and a coarse pointwise lower
bound. This report is a distinct refinement of its enumeration
precision, with an additional macroscopic structural consequence.
The shape proof needs only that earlier coarse bound; it does not
need the Airy correction. Reports~158 and 159 are unchanged.
A bounded literature and catalogue check found no duplicate of the
results stated here in the sources inspected. It is not exhaustive
scholarly priority clearance. The analytic Airy estimates are credited
to Mallein~\cite{mallein}; the local smoothing needed here is proved
explicitly rather than imported from a Tauberian assertion.
\medskip

\section{Definitions and main results}\label{dgp:airy:sec:results}
All partitions in this report are written in weakly increasing order.
Define
\[
 A(n)=\#\left\{(\lambda_1,\ldots,\lambda_K):
 1\leq\lambda_1\leq\cdots\leq\lambda_K,\quad
 \sum_i\lambda_i=n,\quad\lambda_i\geq i\right\},
 \qquad A(0)=1.
\]
The empty partition has length zero. This increasing-order convention
is essential: the weakly decreasing superdiagonal class is different.
Write
\[
 Z_j=\#\{i:\lambda_i=j\},\qquad C_j=\sum_{r=1}^jZ_r,\qquad C_0=0.
\]
Then the defining condition is equivalent to $C_j\leq j$ for every
$j\geq1$.
\begin{writenote}[6 October 2026]
This is Part~III's Lemma~\ref{dgp:rate:lem:barrier}. Report~160's two-line proof of it, the same argument, is not reprinted.
\end{writenote}
We use the constants
\begin{equation}\label{dgp:airy:eq:constants}
 b=\log2,\qquad C=\frac{\pi^2}{12}+b^2,\qquad
 \zeta\approx2.3381074105,
\end{equation}
where $\operatorname{Ai}(-\zeta)=0$ and no Airy zero is closer to zero
on the negative real axis; the displayed decimal is tabulated by NIST~\cite{dlmfairy}.
Generic positive constants will be denoted
$c_*,K_*$ or by indexed lowercase letters, never by the fixed $C$.
\begin{writenote}[6 October 2026]
Part~V writes $a=a_1=-\zeta$ for the same zero, and Li's preprint uses $\zeta$ for a scaled argument (front matter, ``Notation''). ``Report~158'' is Part~III and ``Report~159'' Part~II.
\end{writenote}
Set $F(t)=\sum_{n\geq0}A(n)e^{-tn}$ for $t>0$.
It is finite because it is bounded by the ordinary partition product.

\begin{theorem}[Radial and pointwise Airy corrections]\label{dgp:airy:thm:airy}
As $t\downarrow0$ and, respectively, as integer $n\to\infty$,
\begin{align}
 \log F(t)&=\frac Ct-\zeta b t^{-1/3}+o(t^{-1/3}),\label{dgp:airy:eq:radial}\\
 \log A(n)&=2\sqrt{Cn}-\zeta b C^{-1/6}n^{1/6}
                                      +o(n^{1/6}).\label{dgp:airy:eq:pointwise}
\end{align}
\end{theorem}
The second formula is pointwise in every sufficiently large integer
$n$, not only on a subsequence or after summation. The error term does
not claim an $O(\log n)$ bound. In particular it does not identify an
amplitude or a power of $n$ multiplying the exponential.
\begin{writenote}[6 October 2026]
Theorem~\ref{dgp:airy:thm:airy} is \emph{unconditional}: its inputs are elementary or published (Mallein's estimates, whose statements the write read; note in Section~\ref{dgp:airy:sec:source}). Part~V, conditionally on an unrefereed preprint, proves the stronger $A(n)\sim Kn^{-11/12}e^{B\sqrt n-Dn^{1/6}}$ and $\log F(t)=C/t+ba_1t^{-1/3}+\frac13\log t+O(1)$; its constants agree with these, since $ba_1=-\zeta b$ and Part~V's $D=-a_1bC^{-1/6}$ is this $\zeta bC^{-1/6}$. As long as Part~V is conditional, \eqref{dgp:airy:eq:radial}--\eqref{dgp:airy:eq:pointwise} are this report's unconditional statements of the Airy term. Here $2\sqrt C=B$ of Part~III.
\end{writenote}

\begin{corollary}[Ordinary threshold inverse]\label{dgp:airy:cor:inverse}
For $y>1$, let $T(y)=\min\{n\geq0:A(n)\geq y\}$ and put
\[
 x=\left(\frac{\log y}{2\sqrt C}\right)^2.
\]
Then
\begin{equation}\label{dgp:airy:eq:inverse}
 T(y)=x+\zeta b C^{-2/3}x^{2/3}+o(x^{2/3})
 \qquad(y\to\infty).
\end{equation}
\end{corollary}
This is an asymptotic localization of the integer threshold. It is
not an exact rounding rule or a certified finite-input interval.

For the structural theorem, let $\PP_n$ be the uniform law on the
partitions counted by $A(n)$, and set $t=t_n=\sqrt{C/n}$. Define
\begin{align}
 m(x)&=\begin{cases}1,&0\leq x\leq b,\\
 (e^x-1)^{-1},&x>b,\end{cases}\label{dgp:airy:eq:m}\\
 g(x)&=\int_0^x m(u)\dd u
 =\begin{cases}x,&0\leq x\leq b,\\
 2b+\log(1-e^{-x}),&x\geq b,\end{cases}\label{dgp:airy:eq:g}\\
 f(y)&=g^{-1}(y)
 =\begin{cases}y,&0\leq y\leq b,\\
 -\log(1-e^y/4),&b\leq y<2b.\end{cases}\label{dgp:airy:eq:f}
\end{align}
The function $g$ is strictly increasing and $1$-Lipschitz, with
$g(\infty)=2b$. Its inverse $f$ diverges at $2b$.

\begin{theorem}[Length and macroscopic shape]\label{dgp:airy:thm:shape}
For each fixed $\epsilon>0$, there is $c_\epsilon>0$ such that,
for all sufficiently large $n$,
\begin{align}
 \PP_n(|tK-2b|>\epsilon)&\leq e^{-c_\epsilon/t},\label{dgp:airy:eq:length}\\
 \PP_n\!\left(\sup_{x\geq0}
 |tC_{\floor{x/t}}-g(x)|>\epsilon\right)
 &\leq e^{-c_\epsilon/t}.\label{dgp:airy:eq:globalshape}
\end{align}
Consequently
\begin{equation}\label{dgp:airy:eq:lengthconstant}
 \frac K{\sqrt n}\ \longrightarrow\ \frac{2\log2}{\sqrt C}
       =1.214497355562513\ldots
\end{equation}
in probability and in every finite positive moment.
For any closed interval $I=[a,d_0]\subset(0,2b)$ and fixed
$\epsilon>0$, the probability that some index
$k_t(y)=\floor{y/t}$, $y\in I$, is outside $\{1,\ldots,K\}$,
or that
\[
 \sup_{y\in I}|t\lambda_{k_t(y)}-f(y)|>\epsilon,
\]
is at most $e^{-c_{I,\epsilon}/t}$ for all sufficiently large $n$.
The supremum is evaluated only if all the indicated rows exist;
a missing row is already a bad event.
\end{theorem}
The same interior row assertion holds for ceilings or any indexing
rule with uniform scaled discrepancy $O(t)$. No row-uniform claim is
made at the divergent endpoint $y=2b$.

\section{Exact geometric change of measure}\label{dgp:airy:sec:representation}
Fix an integer $M\geq1$. For $j\leq M$, take independent geometric
multiplicities with
\[
 \PP_0(Z_j=z)=2^{-z-1},\qquad z=0,1,\ldots,
 \qquad H_j=j-\sum_{i\leq j}Z_i,
 \qquad H_0=0.
\]
The increments $Y_j=H_j-H_{j-1}=1-Z_j$ have mean zero, variance $2$,
and a finite absolute exponential moment for every parameter below
$b$. Introduce the left-endpoint area and the prefix weight
\[
 \area_M=\sum_{j=0}^{M-1}H_j,\qquad W_P=\sum_{j=1}^M jZ_j,
 \qquad \eta_j=b-tj.
\]
Finite summation by parts gives the exact identities
\begin{align}
 \sum_{j=1}^M\eta_jZ_j
 &=\sum_{j=1}^M\eta_j-t\area_M-\eta_MH_M,\label{dgp:airy:eq:summation}\\
 W_P&=\frac{M(M+1)}2+\area_M-MH_M.\label{dgp:airy:eq:prefixweight}
\end{align}
They hold for every prefix multiplicity vector, with no barrier assumed.

For $j>M$, independently take
\[
 \PP_t(Z_j=z)=(1-e^{-tj})e^{-tjz}.
\]
The sum of the tail means is finite, so only finitely many tail parts
occur almost surely. Define
\[
 D_M(t)=\exp\!\left(2bM-\frac{tM(M+1)}2\right),\qquad
 P_M(t)=\prod_{j>M}(1-e^{-tj})^{-1}.
\]
Let $Q_{t,M}(h)$ be the tail probability that
\[
 \sum_{i=1}^k Z_{M+i}\leq k+h\quad\hbox{for every }k\geq1.
\]
Combining the prefix geometric mass $2^{-M-\sum Z_j}$ with
\eqref{dgp:airy:eq:summation} yields
\begin{equation}\label{dgp:airy:eq:walkrepresentation}
 F(t)=D_M(t)P_M(t)\,
 \EE_0\!\left[e^{-t\area_M-\eta_MH_M}
       \one_{\{H_j\geq0,\ 0\leq j\leq M\}}
       Q_{t,M}(H_M)\right].
\end{equation}
The factor is $4^M$ before the staircase weight, not $2^M$.
All products and sums here converge absolutely, so this is an identity
of ordinary positive quantities, not merely a formal series identity.

\subsection{The deterministic exponential factor}
Put
\[
 h(x)=-\log(1-e^{-x}),\qquad
 I(x)=\int_x^\infty h(u)\dd u=\operatorname{Li}_2(e^{-x}),
 \qquad L(x)=2bx-\frac{x^2}2+I(x).
\]
For $tM$ in a fixed positive compact interval, monotone Riemann bounds
give $\sum_{j>M}h(tj)=I(tM)/t+O(1)$. Consequently
\begin{equation}\label{dgp:airy:eq:factor}
 \log(D_MP_M)=\frac{L(tM)}t-\frac{tM}2+O(1)
             =\frac Ct+O(1)\quad\hbox{if }tM=b+O(t).
\end{equation}
Indeed,
\[
 L(b)=\frac32b^2+\operatorname{Li}_2(1/2)
      =\frac{\pi^2}{12}+b^2=C.
\]
\begin{writenote}[6 October 2026]
The evaluation is Part~III's \eqref{dgp:rate:eq:Lb}, and $L$ is Part~III's function \eqref{dgp:rate:eq:L} ($2b=\log4$). Report~160's two-sentence proof of the evaluation, by the dilogarithm reflection formula as in Part~III, is not reprinted.
\end{writenote}
Smoothness near $b$ suffices for \eqref{dgp:airy:eq:factor}; no derivative of
an unknown asymptotic remainder is involved.

\section{The published Airy input and its specialization}\label{dgp:airy:sec:source}
We use Mallein's killed-walk area estimates~\cite{mallein}.
The arXiv version~3 assumptions (2.12)--(2.13) and
Lemma~2.6, equation~(2.15), are on manuscript page~11;
Lemma~2.7, equation~(2.19), is on page~13. Both the arXiv and author
manuscripts were checked, including the displayed source equations.
The following is their homogeneous specialization.

For a centered independent-increment walk $S$ with variance $\sigma^2>0$
and an absolute exponential moment, set
\[
 \mathcal K_{m,q}(x,J)=
 \EE_x\!\left[\exp\!\left(-\frac qm\sum_{j=0}^{m-1}S_j\right);
        S_j\geq0\ (0\leq j\leq m),\ S_m\in J\right].
\]
For fixed $q>0$, the source gives
\begin{equation}\label{dgp:airy:eq:malleinupper}
 \limsup_{m\to\infty}m^{-1/3}
 \log\sup_{x\geq0}\mathcal K_{m,q}(x,[0,\infty))
 \leq-\zeta(\sigma^2/2)^{1/3}q^{2/3}.
\end{equation}
For fixed $0<a<c$ and $0<a'<c'$, it also gives
\begin{equation}\label{dgp:airy:eq:malleinlower}
 \liminf_{m\to\infty}m^{-1/3}
 \log\inf_{x\in[am^{1/3},cm^{1/3}]}
 \mathcal K_{m,q}(x,[a'm^{1/3},c'm^{1/3}])
 \geq-\zeta(\sigma^2/2)^{1/3}q^{2/3}.
\end{equation}
We require only integer starting states. The original assumptions
allow a centered independent triangular array with a positive
continuous variance profile and a uniform absolute exponential moment;
there is no symmetry or nonlattice requirement. Our homogeneous walk
$1-\operatorname{Geom}(1/2)$ satisfies every hypothesis, with
$\sigma^2=2$. The multiplier $(\sigma^2/2)^{1/3}$ is therefore exactly
one. The left-endpoint area and killing through time $m$ match
\eqref{dgp:airy:eq:walkrepresentation}.

The potential used below is $t=q_m/m$, with $q_m\to b$.
This does not require an unquoted uniform version of the source:
on nonnegative paths, sandwich the exponential between those with
fixed $q=b-\delta$ and $q=b+\delta$, apply
\eqref{dgp:airy:eq:malleinupper}--\eqref{dgp:airy:eq:malleinlower}, then let
$\delta\downarrow0$. The starting and endpoint intervals remain fixed
when the source is applied.
\begin{writenote}[6 October 2026]
The write read the two cited statements in arXiv:1307.4496v3 \cite{mallein} (journal version not compared). For a triangular array of independent centred variables with $\mathbf E[X_{n,k}^2]=\sigma^2_{k/n}$, $\sigma$ continuous and positive (2.12), and a uniform exponential moment (2.13), Lemma~2.6 (p.~11) bounds $\limsup n^{-1/3}\log\sup_{x\in\mathbb R}\mathbf E_x[e^{-\frac hn\sum_{j<n}S_j};S_j\ge0,j\le n]$ by $\frac{\alpha_1}{2^{1/3}}(h\underline\sigma)^{2/3}$, and Lemma~2.7 (p.~13) bounds the corresponding $\liminf$ of the infimum over starting points in $[an^{1/3},bn^{1/3}]$, with endpoint in $[a'n^{1/3},b'n^{1/3}]$, from below by $\frac{\alpha_1}{2^{1/3}}(h\overline\sigma)^{2/3}$; $\alpha_1\approx-2.3381$ is the largest zero of $\operatorname{Ai}$, $\underline\sigma$ and $\overline\sigma$ the minimum and maximum of $\sigma$. For a constant $\sigma$ both constants equal $-\zeta(\sigma^2/2)^{1/3}h^{2/3}$, as in \eqref{dgp:airy:eq:malleinupper}--\eqref{dgp:airy:eq:malleinlower} (whose upper supremum over $x\ge0$ is weaker than Mallein's over $x\in\mathbb R$). With $\sigma^2=2$ the factor is $1$. The write did not verify Mallein's proofs. Part~V obtains the same Airy constant by an independent route ($ba_1$ in \eqref{dgp:pre:eq:radial}), under its extra input.
\end{writenote}
\begin{writenote}[7 October 2026, batch~108 (reciprocal note)]
The same two inequalities of Mallein, (2.15) of Lemma~2.6 and (2.19) of Lemma~2.7, are the Airy input of the collection's report \file{a397711-bounded-indegree-dags} (labelled acyclic digraphs with every indegree at most two, OEIS A397711; its $n^{1/3}$ term $-3\zeta2^{-1/3}n^{1/3}$ in $\log a_n$). There they are applied to the walk $\mathrm{Poisson}(1)-1$, with $\sigma^2=1$ and hence the factor $(\sigma^2/2)^{1/3}=2^{-1/3}$, at fixed homogeneous tilts on the blocks of a geometric grid, to reach a potential of order $1/i$; here the walk is $1-\operatorname{Geom}(1/2)$, $\sigma^2=2$, with one tilt sequence $q_m\to b$. As here, only fixed tilts and fixed scaled intervals are used, and no strip estimate or time-inhomogeneous case is taken from Mallein. Its statement of the two lemmas agrees with the note above (its independent check re-read them on pp.~11 and~13 of arXiv:1307.4496v3); Mallein's proofs were verified in neither report. No object or theorem is shared.
\end{writenote}

\section{The radial theorem and a localized prefix}\label{dgp:airy:sec:radial}
For the upper bound, take $M=\floor{b/t}$. Then $\eta_M\geq0$ and
$Q_{t,M}\leq1$. Dropping both factors in \eqref{dgp:airy:eq:walkrepresentation}
and applying \eqref{dgp:airy:eq:malleinupper} gives
\[
 \limsup_{t\downarrow0}t^{1/3}
       \left(\log F(t)-\frac Ct\right)\leq-\zeta b.
\]
Here $t^{1/3}M^{1/3}\to b^{1/3}$, while the walk exponent on the
$M^{1/3}$ scale is $-\zeta b^{2/3}$.

\subsection{A tail survival bound}
For a geometric variable $Z$ with parameter $p<1/2$, put
$r=(1-p)/p>1$. Directly,
\begin{equation}\label{dgp:airy:eq:martingale}
 \EE r^{Z-1}=\frac{1-p}{r(1-pr)}=1.
\end{equation}
\begin{writenote}[6 October 2026]
Report~160 next proves, for iid such variables, the bound \eqref{dgp:airy:eq:survival} below. It is Part~III's Lemma~\ref{dgp:rate:lem:survival}, and \eqref{dgp:airy:eq:martingale} is the normalization in its proof. Report~160's proof (the nonnegative martingale $r^{\sum(Z_i-1)}$, the finite-horizon maximal inequality with monotone convergence, and a forced first zero) is Part~III's and is not reprinted. Report~160 remarks that no infinite-time optional-stopping equality is used and that upward jumps need not hit a prescribed level exactly; both hold for Part~III's proof as well.
\end{writenote}
\begin{equation}\label{dgp:airy:eq:survival}
 \PP\!\left(\sum_{i=1}^kZ_i\leq k\ \hbox{for all }k\geq1\right)
 \geq(1-p)(1-1/r)=1-2p.
\end{equation}
For the lower bound choose
\begin{equation}\label{dgp:airy:eq:lowercutoff}
 M=\ceil{b/t}+1.
\end{equation}
Then $-2t\leq\eta_M\leq-t$, and all tail geometric parameters are at
most $p=e^{-t}/2$. Coordinatewise coupling and \eqref{dgp:airy:eq:survival}
therefore give, uniformly for $h\geq0$,
\begin{equation}\label{dgp:airy:eq:tailsurvival}
 Q_{t,M}(h)\geq1-e^{-t}\geq t/2
\end{equation}
for sufficiently small $t$. A nonnegative initial buffer only helps.

\subsection{Repairing the start at zero}
The lower source estimate starts at positive height on the $M^{1/3}$
scale. Fix $0<\rho\leq1$ and force
$\ell=\ceil{\rho M^{1/3}}$ initial increments to be $1$.
The probability is $2^{-\ell}$, the path is a nonnegative ramp, its
area is $\ell(\ell-1)/2$, and it reaches height $\ell$.
Let $m=M-\ell$. For large $M$, the height lies in
$[(\rho/2)m^{1/3},2\rho m^{1/3}]$.
Apply \eqref{dgp:airy:eq:malleinlower} to the remaining $m$ steps with endpoint
interval $[m^{1/3},2m^{1/3}]$. The potential satisfies $tm\to b$.
The ramp area cost is $O(M^{-1/3})$, whereas
$\ell b/M^{1/3}\to\rho b$. Thus
\begin{equation}\label{dgp:airy:eq:zerostart}
 \liminf_{t\downarrow0}M^{-1/3}
 \log\EE_0[e^{-t\area_M};\ H_j\geq0\ (j\leq M)]
 \geq-\zeta b^{2/3}-\rho b.
\end{equation}
Taking the limit with $\rho$ fixed and then letting $\rho\downarrow0$
removes the boundary cost. In \eqref{dgp:airy:eq:walkrepresentation},
$e^{-\eta_MH_M}\geq1$. Equations \eqref{dgp:airy:eq:factor} and
\eqref{dgp:airy:eq:tailsurvival}, whose logarithmic loss is
$o(t^{-1/3})$, now prove the radial lower bound and hence
\eqref{dgp:airy:eq:radial}.

\subsection{A full fixed event for the pointwise argument}
Choose a fixed sufficiently large $K_0$ with
$K_0b>\zeta b^{2/3}+b+2$. At cutoff \eqref{dgp:airy:eq:lowercutoff}, define
\begin{equation}\label{dgp:airy:eq:fixedprefixevent}
 E_t=\{H_j\geq0\ (0\leq j\leq M),\quad
          0\leq H_M\leq2M^{1/3},\quad
          \area_M\leq K_0M^{4/3}\}.
\end{equation}
This is the entire three-constraint event, independent of the ramp
parameter $\rho$. The ramp subevents above end below $2M^{1/3}$.
On a nonnegative path the weighted mass discarded by imposing the
area constraint is at most
$\exp(-tK_0M^{4/3})$.
For every fixed $\rho\in(0,1]$, our choice of $K_0$ makes this
exponentially smaller than the ramp lower bound on the $M^{1/3}$
scale. Subtract it first, then let $\rho\downarrow0$.
The matching upper bound is \eqref{dgp:airy:eq:malleinupper}. We obtain
\begin{equation}\label{dgp:airy:eq:prefixmass}
 \log\EE_0[e^{-t\area_M};E_t]
          =-\zeta b t^{-1/3}+o(t^{-1/3}).
\end{equation}
Moreover, \eqref{dgp:airy:eq:prefixweight} gives the deterministic localization
\begin{equation}\label{dgp:airy:eq:prefixwindow}
 \left|W_P-\frac{M(M+1)}2\right|
 \leq(K_0+2)M^{4/3}=O(t^{-4/3})\qquad\hbox{on }E_t.
\end{equation}
The area restriction is an extra ingredient for pointwise extraction;
radial asymptotics alone do not supply it.

\section{Exact weight localization in the geometric tail}\label{dgp:airy:sec:tail}
Retain the cutoff \eqref{dgp:airy:eq:lowercutoff} and independent tail law
$\PP_t$. Put
\[
 W_T=\sum_{j>M}jZ_j,\quad \mu_T=\EE_tW_T,\quad
 B_k=k-\sum_{i=1}^kZ_{M+i}.
\]
\begin{lemma}[Uniform exact weight and survival]\label{dgp:airy:lem:localization}
For every fixed $K_1>0$ there are positive constants $c_1,K_2$ such
that, uniformly over integers $w$ with
$|w-\mu_T|\leq K_1t^{-4/3}$ and every $h\geq0$,
\begin{equation}\label{dgp:airy:eq:taillocal}
 \PP_t(W_T=w,\ B_k+h\geq0\ \hbox{for all }k\geq1)
 \geq c_1t^{5/2}\exp(-K_2t^{-1/6})
\end{equation}
for sufficiently small $t$.
\end{lemma}
Its logarithmic cost is $o(t^{-1/3})$, strictly smaller than the Airy
scale. We prove every component, including lattice smoothing.

\subsection{A second cutoff and a surviving head with a buffer}
Fix $d>b$, for example $d=b+1$, and set $J=\floor{d/t}$ and $N=J-M$.
Split the tail into
\[
 U=\sum_{j=M+1}^JjZ_j,\quad
 G=N-\sum_{j=M+1}^J Z_j,\quad
 V=\sum_{j>J}jZ_j,
 \qquad W_T=U+V.
\]
Write $\mu_U=\EE_tU$ and $\mu_V=\EE_tV$.
The stricter zero-buffer head barrier
$\sum_{i=1}^kZ_{M+i}\leq k$ for $1\leq k\leq N$
has probability at least $t/2$ by \eqref{dgp:airy:eq:tailsurvival}.
Riemann sums on $[b,d]$ give
\begin{equation}\label{dgp:airy:eq:buffermean}
 \EE_tG=\frac\beta t+O(1),\qquad
 \beta=(d-b)-\int_b^d\frac{\dd u}{e^u-1}
       =d-2b-\log(1-e^{-d})>0.
\end{equation}
Strict positivity follows from $(e^u-1)^{-1}<1$ for $u>b$.
Set
\begin{equation}\label{dgp:airy:eq:ascale}
 a(t)=t^{-19/12}.
\end{equation}
We claim
\begin{align}
 \PP_t(G<\beta/(2t))&\leq e^{-c_2/t},\label{dgp:airy:eq:bufferchernoff}\\
 \PP_t(|U-\mu_U|>a(t))&\leq2e^{-c_3t^{-1/6}}.\label{dgp:airy:eq:headchernoff}
\end{align}
There are $O(1/t)$ head variables, all with geometric parameter at
most $1/2$. Their unweighted centered logarithmic moment generating
functions sum to at most $K_3s^2/t$ on a fixed neighborhood of zero.
Since $\EE G\geq3\beta/(4t)$ eventually, a fixed sufficiently small
tilt proves \eqref{dgp:airy:eq:bufferchernoff}.
For the weighted head sum,
\[
 \frac{\dd^2}{\dd s^2}\log\EE_t e^{sU}
 =\sum_{j=M+1}^J\frac{j^2e^{-(t-s)j}}{(1-e^{-(t-s)j})^2}
 \leq K_4t^{-3},\qquad |s|\leq t/4.
\]
The scaled arguments stay away from zero and $j=O(1/t)$.
Taylor's theorem bounds the centered log moment generating function
by $K_4s^2t^{-3}/2$. Tilts of either sign, of magnitude a suitable
constant times $a(t)t^3=t^{17/12}=o(t)$, prove
\eqref{dgp:airy:eq:headchernoff} because $a(t)^2t^3=t^{-1/6}$.

Let $E_H$ require the zero-buffer head barrier,
$G\geq\beta/(2t)$, and $|U-\mu_U|\leq a(t)$.
A union bound gives
\begin{equation}\label{dgp:airy:eq:goodhead}
 \PP_t(E_H)\geq t/2-e^{-c_2/t}-2e^{-c_3t^{-1/6}}\geq t/3.
\end{equation}
This does not assume concentration after conditioning on survival.
The surviving head produces a buffer of order $1/t$ for the
independent far tail.

\subsection{The far tail local limit on every Fourier arc}\label{dgp:airy:sec:fourier}
Keep $J=\floor{d/t}$ fixed while varying a parameter $s>0$.
For $j>J$, take independent geometrically distributed variables with
parameter $e^{-sj}$. Define
\begin{align*}
 Z_J(s)&=\prod_{j>J}(1-e^{-sj})^{-1},& K_J(s)&=\log Z_J(s),\\
 \mu_J(s)&=\sum_{j>J}\frac j{e^{sj}-1}=-K_J'(s),&
 v_J(s)&=\sum_{j>J}\frac{j^2e^{-sj}}{(1-e^{-sj})^2}=K_J''(s).
\end{align*}
All differentiated sums converge locally uniformly for $s>0$.
Uniformly whenever $s/t\to1$, Riemann sums and exponential tails give
\begin{equation}\label{dgp:airy:eq:variance}
 v_J(s)\sim v_dt^{-3},\qquad
 v_d=\int_d^\infty\frac{x^2e^{-x}}{(1-e^{-x})^2}\dd x>0.
\end{equation}
The sum of the third-cumulant majorants is $O(t^{-4})$, uniformly in
a small relative neighborhood of $t$.

We prove that whenever $k=\mu_J(s)$ is an integer and $s/t\to1$,
\begin{equation}\label{dgp:airy:eq:llt}
 \PP_s(V=k)=\frac{1+o(1)}{\sqrt{2\pi v_J(s)}}.
\end{equation}
The characteristic function is the absolutely convergent product
\[
 \phi_s(\theta)=\prod_{j>J}
            \frac{1-e^{-sj}}{1-e^{-sj}e^{ij\theta}}.
\]
With $p_j=e^{-sj}$, its centered logarithm has the absolutely
convergent double series
\begin{equation}\label{dgp:airy:eq:doublelog}
 \log\phi_s(\theta)-i\mu_J(s)\theta
 =\sum_{j>J}\sum_{r\geq1}\frac{p_j^r}{r}
                  (e^{irj\theta}-1-irj\theta).
\end{equation}
The logarithm here is defined by this convergent expansion, so no
branch ambiguity is introduced. Taylor's elementary remainder bound
gives
\begin{align}
 \log\phi_s(\theta)-i\mu_J(s)\theta
 &=-\tfrac12v_J(s)\theta^2+O(t^{-4}|\theta|^3),\label{dgp:airy:eq:central}\\
 \sum_{j>J}\frac{j^3p_j(1+p_j)}{(1-p_j)^3}&=O(t^{-4}).\nonumber
\end{align}
The last bound follows by summing $j^3e^{-sj}$, since $sj$ is bounded
away from zero. It is uniform for the tilts under consideration.

For the consecutive block $j=J+1,\ldots,J+\floor{1/t}$, the parameters
$p_j$ lie in one fixed compact subinterval of $(0,1)$. The exact identity
\[
 \left|\frac{1-p}{1-pe^{iu}}\right|^2
 =\left(1+\frac{4p\sin^2(u/2)}{(1-p)^2}\right)^{-1}
\]
therefore implies
\begin{equation}\label{dgp:airy:eq:blockfourier}
 |\phi_s(\theta)|\leq
 \exp\!\left(-c_4\sum_{\rm block}\sin^2(j\theta/2)\right).
\end{equation}
We now cover the entire inversion interval $[-\pi,\pi]$.
\begin{enumerate}
\item Choose a fixed small $c_0>0$. If $|\theta|\leq c_0t$,
all block values $|j\theta|$ are below $\pi$, whence
\[
 \sum_{\rm block}\sin^2(j\theta/2)
       \geq c_5\theta^2\sum_{\rm block}j^2
       \geq c_6\theta^2t^{-3}.
\]
\item If $A_0t\leq|\theta|\leq\pi$, the finite geometric series gives
\[
 \sum_{\rm block}\sin^2(j\theta/2)
 \geq\frac{\floor{1/t}}2-\frac1{2|\sin(\theta/2)|}
 \geq c_7/t
\]
for sufficiently large fixed $A_0$, using
$|\sin(\theta/2)|\geq|\theta|/\pi$.
\item In the remaining region $c_0t\leq|\theta|\leq A_0t$, write
$u=\theta/t$. Uniform Riemann sums yield
\[
 t\sum_{\rm block}\sin^2(j\theta/2)
 \longrightarrow\int_d^{d+1}\sin^2(ux/2)\dd x
\]
uniformly on $c_0\leq|u|\leq A_0$. This continuous integral has a
strictly positive minimum there: a nonzero-frequency sine cannot
vanish throughout an interval. Thus the sum is again at least $c_8/t$.
\end{enumerate}
There are no unexamined secondary or rational arcs.
At $\theta=t^{3/2}x$, the remainder in \eqref{dgp:airy:eq:central} is
$O(t^{1/2}|x|^3)$ on bounded $x$-sets. Equations
\eqref{dgp:airy:eq:variance} and \eqref{dgp:airy:eq:blockfourier} give integrable
Gaussian domination on the inner arc after this rescaling; the total
outer-arc integral is $O(e^{-c_9/t})=o(t^{3/2})$.
Fourier inversion for the integer-valued $V$ and dominated convergence
prove \eqref{dgp:airy:eq:llt}. All estimates are uniform on shrinking relative
tilt bands $s/t\to1$. In particular they cover every tilt below.
Consecutive allowed weights also give span one; all integer targets
at least $J+1$ are represented by a single allowed part.
\begin{writenote}[6 October 2026]
Part~V, Section~\ref{dgp:pre:sec:local}, uses the same cutoff $J=\lfloor d/t\rfloor$, the same block of $\lfloor1/t\rfloor$ consecutive indices and the same three-arc cover of $[-\pi,\pi]$ (\eqref{dgp:pre:eq:smallarc}--\eqref{dgp:pre:eq:allotherarcs}), but assembles them differently: a moment-generating-function limit and a total-variation cutoff instead of the tilted local limit theorem with exact weights used here. Both are printed, as two routes.
\end{writenote}

\subsection{Moderate retilting with a fixed cutoff}
Suppose $k$ is an integer with
\begin{equation}\label{dgp:airy:eq:targetwindow}
 |k-\mu_J(t)|\leq2a(t).
\end{equation}
It is positive and of order $t^{-2}$, because $a(t)=o(t^{-2})$.
The mean $\mu_J(s)$ is continuous and strictly decreasing from infinity
to zero. In a small relative neighborhood of $t$, its negative
derivative $v_J(s)$ is bounded above and below by positive multiples
of $t^{-3}$. Integration of the lower bound brackets the unique root
of $\mu_J(s)=k$ and gives
\begin{equation}\label{dgp:airy:eq:tiltscale}
 |s-t|=O(a(t)t^3)=O(t^{17/12})=o(t).
\end{equation}
Here $J$ remains fixed; it is not redefined as $\floor{d/s}$.
The exact likelihood ratio is
\begin{equation}\label{dgp:airy:eq:likelihood}
 \PP_t(V=k)=\PP_s(V=k)
   \exp\{(s-t)k+K_J(s)-K_J(t)\}.
\end{equation}
Expanding $K_J(t)$ about $s$ and using $K_J'(s)=-k$ shows that the
exponent equals
\[
 -\tfrac12K_J''(\xi)(s-t)^2
 \geq-K_5 t^{-3}a(t)^2t^6=-K_5t^{-1/6}.
\]
Its sign is nonpositive, as convexity also confirms.
Combining this with \eqref{dgp:airy:eq:llt} proves the uniform bound
\begin{equation}\label{dgp:airy:eq:localmass}
 \PP_t(V=k)\geq c_{10}t^{3/2}e^{-K_5t^{-1/6}}
 \quad\hbox{throughout \eqref{dgp:airy:eq:targetwindow}}.
\end{equation}
No asymptotic remainder has been differentiated.

\subsection{Restoring the far tail barrier by subtraction}
Choose fixed $p_*\in(e^{-d},1/2)$ and let $r_*=(1-p_*)/p_*>1$.
The original $t$-law far-tail variables are dominated coordinatewise
by iid geometrics with parameter $p_*$. From \eqref{dgp:airy:eq:martingale}
and the same maximal inequality, ruin from a buffer $g\geq0$ obeys
\begin{equation}\label{dgp:airy:eq:ruin}
 \PP_t\!\left(\exists k\geq1:
       \sum_{i=1}^k(Z_{J+i}-1)>g\right)
 \leq r_*^{-\floor g-1}.
\end{equation}
On $E_H$, the buffer $G$ is at least $\beta/(2t)$, so ruin has
probability at most $e^{-c_{11}/t}$. An additional prefix buffer
$h\geq0$ cannot worsen the bound.
For a retained head weight $u$ and any target of
Lemma~\ref{dgp:airy:lem:localization},
\[
 |(w-u)-\mu_V|\leq K_1t^{-4/3}+a(t)\leq2a(t)
\]
for sufficiently small $t$. Apply \eqref{dgp:airy:eq:localmass} to this integer
weight and subtract \eqref{dgp:airy:eq:ruin}. The latter is negligible relative
to the exact local probability, not just relative to one. Uniformly
in the good head realization, the joint exact-weight and barrier
probability is at least
\[
 \frac{c_{10}}2t^{3/2}e^{-K_5t^{-1/6}}.
\]
Both probabilities being subtracted are under the original $t$ law.
Integrate over the independent head event of probability at least
$t/3$ in \eqref{dgp:airy:eq:goodhead}; this proves
\eqref{dgp:airy:eq:taillocal}. Neither weight and survival nor survival and
head concentration have been assumed independent.

\section{Pointwise extraction and integer inversion}\label{dgp:airy:sec:pointwise}
The radial expansion and coefficient positivity immediately give
$A(n)e^{-tn}\leq F(t)$. Taking $t=\sqrt{C/n}$ proves the upper half
of \eqref{dgp:airy:eq:pointwise}. For the lower half use the localized prefix
and Lemma~\ref{dgp:airy:lem:localization}.

An integration by parts gives
\[
 \int_b^\infty\frac{u}{e^u-1}\dd u
 =b\,h(b)+I(b)=b^2+\operatorname{Li}_2(1/2).
\]
Smooth Riemann sums away from zero and $tM=b+O(t)$ yield
\begin{align}
 \mu_T&=t^{-2}\bigl(b^2+\operatorname{Li}_2(1/2)\bigr)+O(t^{-1}),
                                                        \nonumber\\
 \frac{M(M+1)}2+\mu_T&=\frac{C}{t^2}+O(t^{-1})=n+O(t^{-1}).
                                                        \label{dgp:airy:eq:center}
\end{align}
By \eqref{dgp:airy:eq:prefixwindow}, each retained prefix therefore leaves an
integer residual target $w=n-W_P$ within a fixed $O(t^{-4/3})$
window of $\mu_T$. This target is positive for small $t$.

The exact coefficient version of \eqref{dgp:airy:eq:walkrepresentation} is
\begin{align}
 A(n)e^{-tn}=D_MP_M\,
 \EE_0\bigl[&e^{-t\area_M-\eta_MH_M}
      \one_{\{H_j\geq0,\ j\leq M\}}\nonumber\\[-2pt]
 &\times\PP_t(W_T=n-W_P,\ B_k+H_M\geq0\ \forall k\geq1)\bigr].
                                                        \label{dgp:airy:eq:coefficientexact}
\end{align}
Restrict to $E_t$. Since $\eta_M\leq0$, its terminal factor is at
least one. Apply \eqref{dgp:airy:eq:factor}, \eqref{dgp:airy:eq:prefixmass}, and
\eqref{dgp:airy:eq:taillocal}, uniformly over all retained prefixes.
The tail logarithmic loss satisfies
\[
 \log\bigl(c_1t^{5/2}e^{-K_2t^{-1/6}}\bigr)=o(t^{-1/3}).
\]
Thus
\[
 \log A(n)-tn\geq C/t-\zeta b t^{-1/3}+o(t^{-1/3}).
\]
At $t=\sqrt{C/n}$, $tn=C/t=\sqrt{Cn}$ and
$t^{-1/3}=C^{-1/6}n^{1/6}$. This proves the matching lower bound and
completes Theorem~\ref{dgp:airy:thm:airy}.
The compatible localization scales are worth recording:
\[
 t^{-4/3}\ll t^{-3/2}\ll t^{-19/12}\ll t^{-5/3},
 \qquad a(t)^2t^3=t^{-1/6}\ll t^{-1/3}.
\]
The proof has not inferred a pointwise coefficient lower bound from
a radial expansion by itself.

\subsection{Threshold inversion without differentiating an error}
\begin{writenote}[6 October 2026]
Report~160 first shows by the largest-part injection that $A(n)$ is nondecreasing. This is Part~III's Lemma~\ref{dgp:rate:lem:mono}, and Report~160's proof, the same injection, is not reprinted.
\end{writenote}
Thus $A(n)$ is nondecreasing; \eqref{dgp:airy:eq:pointwise} also shows it tends
to infinity.
Put $B=2\sqrt C$, $D=\zeta b C^{-1/6}$, $u=\log y$, and
$x=(u/B)^2$. For fixed real $k$ and any integer
$n=x+kx^{2/3}+O(1)$, expansion of the explicit terms gives
\[
 B\sqrt n-Dn^{1/6}=u+(Bk/2-D)x^{1/6}+o(x^{1/6}).
\]
The unspecified remainder in \eqref{dgp:airy:eq:pointwise} remains
$o(x^{1/6})$ on either such sequence because $n/x\to1$.
For fixed $\epsilon>0$, take
\[
 n_- =\floor{x+(2D/B-\epsilon)x^{2/3}},\qquad
 n_+ =\ceil{x+(2D/B+\epsilon)x^{2/3}}.
\]
Eventually $A(n_-)<y\leq A(n_+)$, whence
$n_-<T(y)\leq n_+$. Letting $\epsilon\downarrow0$ after the limit
proves Corollary~\ref{dgp:airy:cor:inverse}, since
$2D/B=\zeta b C^{-2/3}$. Integer rounding contributes only $O(1)$.
\begin{writenote}[6 October 2026]
Corollary~\ref{dgp:airy:cor:inverse} refines Part~III's Theorem~\ref{dgp:rate:thm:inverse}; Part~V's Corollary~\ref{dgp:pre:cor:inverse} extends it, conditionally, with the same coefficient $2D/B=1.3585577\ldots$. Because of the growing term $n^{1/6}$ it is not an instance of the transseries volume's theorems (front matter, ``The inverses'').
\end{writenote}

\section{Macroscopic concentration from a product reference law}\label{dgp:airy:sec:shapeproof}
We now prove Theorem~\ref{dgp:airy:thm:shape}. The only enumerative input
needed is the weaker statement
\begin{equation}\label{dgp:airy:eq:coarse}
 \log A(n)\geq2\sqrt{Cn}-K_6n^{1/3}
 \quad\hbox{for all sufficiently large integers }n.
\end{equation}
It was already established in Report~158 and also follows from
Theorem~\ref{dgp:airy:thm:airy}. Thus the structural argument does not depend
on the Airy correction or its local-limit proof. Even the exact
logarithmic rate with an $o(\sqrt n)$ error would suffice.
\begin{writenote}[6 October 2026]
``Report~158'' is Part~III; \eqref{dgp:airy:eq:coarse} is the lower half of its \eqref{dgp:rate:eq:coefmain}. Theorem~\ref{dgp:airy:thm:shape} is therefore unconditional and independent of Section~\ref{dgp:airy:sec:source}.
\end{writenote}

\subsection{Exact conditioning and its cost}
Take $M=\floor{b/t}$ and let $\QQ_t$ be the independent product law
with geometric parameters
\[
 p_j=\begin{cases}1/2,&j\leq M,\\ e^{-tj},&j>M.\end{cases}
\]
Every parameter is at most $1/2$, and the means are exactly
$\EE_{\QQ_t}Z_j=m(tj)$, including a possible equality $tM=b$.
Since $\sum_jm(tj)<\infty$, this law is supported on finite partitions.
Set
\[
 \mathcal U_M=D_MP_M,\qquad
 \mathcal W_t=e^{-t\area_M-\eta_MH_M},\qquad
 \mathcal B=\{H_j\geq0\ \hbox{for all }j\geq1\},\quad
 \mathcal N=\sum_{j\geq1}jZ_j.
\]
The pointwise change of measure is
\begin{equation}\label{dgp:airy:eq:pointwisecm}
 e^{-t|\lambda|}=\mathcal U_M\QQ_t(\{\lambda\})
                      \mathcal W_t(\lambda).
\end{equation}
On the full barrier $\mathcal B$, $0<\mathcal W_t\leq1$ because
$\eta_M\geq0$. For every event $E$ in the multiplicities,
including a $t$-dependent event,
\begin{align}
 \PP_n(E)&=\frac{\EE_{\QQ_t}[\one_E\one_{\mathcal B}
                          \one_{\{\mathcal N=n\}}\mathcal W_t]}{R_n(t)},
                                                        \label{dgp:airy:eq:decondition}\\
 R_n(t)&=\frac{A(n)e^{-tn}}{\mathcal U_M},\qquad
 \PP_n(E)\leq\frac{\QQ_t(E)}{R_n(t)}.\nonumber
\end{align}
The barrier includes every tail index.
At $t=\sqrt{C/n}$, \eqref{dgp:airy:eq:factor} and \eqref{dgp:airy:eq:coarse} give
\begin{equation}\label{dgp:airy:eq:conditioncost}
 \log R_n(t)\geq-K_7t^{-2/3}-K_8,\qquad
 \PP_n(E)\leq e^{K_7t^{-2/3}+K_8}\QQ_t(E).
\end{equation}
Therefore any product-law bound $A_*e^{-c_*/t}$ transfers to
$e^{-c_*/(2t)}$ for sufficiently small $t$.
No separate estimate of $\QQ_t(\mathcal N=n)$ or independence after
conditioning is required.

\subsection{Exponential moments uniformly over infinite subsets}
Put $s_0=b/2$. For a geometric variable of parameter $p\leq1/2$
and mean $\mu_p=p/(1-p)$, its centered log moment generating function
is
\[
 \ell_p(s)=\log(1-p)-\log(1-pe^s)-s\mu_p,\qquad |s|\leq s_0.
\]
It has value and derivative zero at zero and
\[
 \ell_p''(s)=\frac{pe^s}{(1-pe^s)^2}
 \leq\frac{e^{s_0}\mu_p}{(1-e^{s_0}/2)^2}.
\]
Taylor's theorem thus gives $\ell_p(s)\leq a_0s^2\mu_p$, where
$a_0=e^{s_0}/[2(1-e^{s_0}/2)^2]$.
Independence implies the same summed inequality over any finite
index set. It extends to an infinite set by increasing finite
subsets: use monotone convergence for uncentered exponentials when
$s\geq0$, and dominated convergence (the exponentials are at most
one) when $s<0$. The corresponding deterministic mean sums converge.
Since $m$ is nonincreasing,
\[
 t\sum_{j\geq1}m(tj)\leq\int_0^\infty m(x)\dd x=2b.
\]
With $V_0=2ba_0$, for every finite or infinite subset $I_0\subset\NN$,
\begin{equation}\label{dgp:airy:eq:subsetmgf}
 \log\EE_{\QQ_t}\exp\!\left(s\sum_{j\in I_0}(Z_j-m(tj))\right)
 \leq V_0s^2/t,\qquad |s|\leq s_0.
\end{equation}
For fixed $r>0$, choose $s_r=\min(s_0,r/(2V_0))$.
Chernoff's inequality with both signs yields
\begin{equation}\label{dgp:airy:eq:subsetchernoff}
 \QQ_t\!\left(t\left|\sum_{j\in I_0}(Z_j-m(tj))\right|\geq r\right)
 \leq2e^{-\kappa_r/t},\qquad \kappa_r=rs_r/2>0.
\end{equation}
The constants work simultaneously for every prefix and for the
infinite total length.

\subsection{Global Riemann bounds and length moments}
For $J_0=\floor{x/t}$, monotonicity of $m$ gives
\[
 0\leq\int_0^{tJ_0}m(u)\dd u-t\sum_{j=1}^{J_0}m(tj)
 \leq t\bigl(m(0)-m(tJ_0)\bigr)\leq t.
\]
The remaining interval has length below $t$ and $m\leq1$, so
\begin{equation}\label{dgp:airy:eq:globalriemann}
 0\leq g(x)-t\EE_{\QQ_t}C_{\floor{x/t}}\leq2t
 \quad\hbox{for every }x\geq0.
\end{equation}
This estimate is global, even for $x$ growing as $t$ decreases.
Letting the integer cutoff tend to infinity also gives
\begin{equation}\label{dgp:airy:eq:lengthmean}
 0\leq2b-t\EE_{\QQ_t}K\leq t.
\end{equation}
Equations \eqref{dgp:airy:eq:subsetchernoff}, \eqref{dgp:airy:eq:lengthmean} and
\eqref{dgp:airy:eq:conditioncost} prove \eqref{dgp:airy:eq:length}.
Under $\PP_n$, the deterministic constraint implies
$K(K+1)/2\leq n$, hence $K/\sqrt n\leq\sqrt2$.
Convergence in probability therefore implies convergence of every
finite positive moment, proving \eqref{dgp:airy:eq:lengthconstant} in the
stated sense.

\subsection{Uniform cumulative concentration}
Fix $L>0$ and $\epsilon>0$. Choose a finite grid of $[0,L]$ with mesh
at most $\epsilon/4$. If the deviation in \eqref{dgp:airy:eq:globalshape}
is at most $\epsilon/2$ at all grid points, monotonicity of the
cumulative count and the $1$-Lipschitz property of $g$ bound it by
$3\epsilon/4$ between adjacent points.
By \eqref{dgp:airy:eq:globalriemann}, a grid deviation exceeding
$\epsilon/2$ implies a centered deviation exceeding $\epsilon/4$
for sufficiently small $t$. A finite union bound in
\eqref{dgp:airy:eq:subsetchernoff} consequently proves
\[
 \QQ_t\!\left(\sup_{0\leq x\leq L}
       |tC_{\floor{x/t}}-g(x)|>\epsilon\right)
 \leq A_{L,\epsilon}e^{-c_{L,\epsilon}/t}.
\]
The grid is fixed after $L,\epsilon$ are fixed.
To reach the whole half-line, choose fixed $L$ with
$2b-g(L)<\epsilon/4$. On the event that the compact supremum is at
most $\epsilon/4$ and $|tK-2b|\leq\epsilon/4$, for $x\geq L$,
\begin{align*}
 tC_{\floor{x/t}}-g(x)&\leq tK-g(L)\leq\epsilon/2,\\
 g(x)-tC_{\floor{x/t}}&\leq2b-tC_{\floor{L/t}}\leq\epsilon/2.
\end{align*}
The complement has product-law probability at most
$A_\epsilon e^{-c_\epsilon/t}$. Transfer by
\eqref{dgp:airy:eq:conditioncost} proves \eqref{dgp:airy:eq:globalshape}.

\section{Row inversion and structural consequences}\label{dgp:airy:sec:rows}
Solving \eqref{dgp:airy:eq:g} for $x$ gives \eqref{dgp:airy:eq:f}; both pieces and their
first derivatives agree at $b$. We give the integer-index argument
for the claimed compact-interior row convergence.
Fix $I=[a,d_0]\subset(0,2b)$ and $\epsilon>0$. Put
$\delta=\min(\epsilon/2,a/2)>0$. Since $f(y)\geq y\geq a$,
$f(y)-\delta$ is nonnegative. The two continuous positive functions
\[
 y-g(f(y)-\delta),\qquad g(f(y)+\delta)-y
\]
have a common positive lower bound $\Delta$ on $I$. Set
$r=\Delta/3$ and suppose
\[
 \sup_{x\geq0}|tC_{\floor{x/t}}-g(x)|\leq r,
 \qquad t<\min(a/2,r).
\]
For $y\in I$ put $k=\floor{y/t}\geq1$ and
$x_\pm=f(y)\pm\delta$. As $y-t<tk\leq y$,
\[
 tC_{\floor{x_-/t}}\leq y-2r<tk,
 \qquad tC_{\floor{x_+/t}}\geq y+2r>tk.
\]
Thus $C_{\floor{x_-/t}}<k\leq C_{\floor{x_+/t}}$.
The right inequality proves the row exists. Since the list is sorted,
\[
 \floor{x_-/t}<\lambda_k\leq\floor{x_+/t},
 \qquad x_-<t\lambda_k\leq x_+.
\]
This holds simultaneously for all $y\in I$ and bounds the row
supremum by $\delta<\epsilon$. Its bad event is contained in the
global cumulative bad event at tolerance $r$, proving the row part
of Theorem~\ref{dgp:airy:thm:shape}. An $O(t)$ change of scaled indices is
absorbed by the same strict gaps.

\subsection{Half of the limiting row mass follows the diagonal}
At $M=\floor{b/t}$, \eqref{dgp:airy:eq:globalshape} gives $tC_M\to b$ while
\eqref{dgp:airy:eq:length} gives $tK\to2b$. Therefore
\begin{equation}\label{dgp:airy:eq:halfmass}
 \frac{C_{\floor{b/t}}}{K}\longrightarrow\frac12
\end{equation}
with exponentially small probability of every fixed-tolerance
deviation. This counts rows whose part sizes are at most the critical
scale $b/t$. Equivalently, the initial half of the limiting row curve
is $f(y)=y$. For every $0<a<d_0<b$,
\[
 \sup_{y\in[a,d_0]}t\bigl(\lambda_{\floor{y/t}}-\floor{y/t}\bigr)
 \longrightarrow0
\]
with the same kind of fixed-tolerance exponential bound.
These are macroscopic statements. They do not determine the number
or density of exact equalities $\lambda_i=i$.

\subsection{Area normalization and scope}
Changing variables $y=g(x)$ gives
\begin{align*}
 \int_0^{2b}f(y)\dd y
 &=\int_0^\infty x m(x)\dd x\\
 &=\frac{b^2}2+\int_b^\infty\frac{x}{e^x-1}\dd x
 =\frac32b^2+\operatorname{Li}_2(1/2)=C.
\end{align*}
This matches the deterministic rescaled area $t^2n=C$.
It is a consistency check, not a substitute for concentration or an
assertion of unproved extreme-row control.
The constants in the exponential probability bounds depend on the
fixed tolerance and, for rows, the fixed interior interval.
No optimal constant, central limit theorem, covariance process,
diagonal fluctuation scale, largest-part centering, or extreme-part
law is proved here.
\begin{writenote}[6 October 2026]
These open items are collected in Section~\ref{dgp:sec:further}, item~\ref{dgp:q:fluct}.
\end{writenote}

\begin{figure}[tbp]
\centering
\setlength{\unitlength}{1pt}
\begin{picture}(452,196)
\small
\put(22,26){\vector(1,0){181}}
\put(22,26){\vector(0,1){145}}
\put(109,184){\makebox(0,0){Cumulative profile}}
\put(15,18){\makebox(0,0){$0$}}
\put(252,26){\vector(1,0){181}}
\put(252,26){\vector(0,1){145}}
\put(339,184){\makebox(0,0){Row profile}}
\put(245,18){\makebox(0,0){$0$}}
\put(214,16){\makebox(0,0){$x$}}
\put(11,172){\makebox(0,0){$g(x)$}}
\put(52.325,24){\line(0,1){4}}
\put(52.325,15){\makebox(0,0){$b$}}
\put(109.500,24){\line(0,1){4}}
\put(109.500,15){\makebox(0,0){$2$}}
\put(197.000,24){\line(0,1){4}}
\put(197.000,15){\makebox(0,0){$4$}}
\put(20,88.845){\line(1,0){4}}
\put(11,88.845){\makebox(0,0){$b$}}
\put(20,151.691){\line(1,0){4}}
\put(11,151.691){\makebox(0,0){$2b$}}
\put(22,151.691){\line(1,0){2}}
\put(27,151.691){\line(1,0){2}}
\put(32,151.691){\line(1,0){2}}
\put(37,151.691){\line(1,0){2}}
\put(42,151.691){\line(1,0){2}}
\put(47,151.691){\line(1,0){2}}
\put(52,151.691){\line(1,0){2}}
\put(57,151.691){\line(1,0){2}}
\put(62,151.691){\line(1,0){2}}
\put(67,151.691){\line(1,0){2}}
\put(72,151.691){\line(1,0){2}}
\put(77,151.691){\line(1,0){2}}
\put(82,151.691){\line(1,0){2}}
\put(87,151.691){\line(1,0){2}}
\put(92,151.691){\line(1,0){2}}
\put(97,151.691){\line(1,0){2}}
\put(102,151.691){\line(1,0){2}}
\put(107,151.691){\line(1,0){2}}
\put(112,151.691){\line(1,0){2}}
\put(117,151.691){\line(1,0){2}}
\put(122,151.691){\line(1,0){2}}
\put(127,151.691){\line(1,0){2}}
\put(132,151.691){\line(1,0){2}}
\put(137,151.691){\line(1,0){2}}
\put(142,151.691){\line(1,0){2}}
\put(147,151.691){\line(1,0){2}}
\put(152,151.691){\line(1,0){2}}
\put(157,151.691){\line(1,0){2}}
\put(162,151.691){\line(1,0){2}}
\put(167,151.691){\line(1,0){2}}
\put(172,151.691){\line(1,0){2}}
\put(177,151.691){\line(1,0){2}}
\put(182,151.691){\line(1,0){2}}
\put(187,151.691){\line(1,0){2}}
\put(192,151.691){\line(1,0){2}}
\put(444,16){\makebox(0,0){$y$}}
\put(240,172){\makebox(0,0){$f(y)$}}
\put(339.500,24){\line(0,1){4}}
\put(339.500,15){\makebox(0,0){$b$}}
\put(427.000,24){\line(0,1){4}}
\put(427.000,15){\makebox(0,0){$2b$}}
\put(250,49.567){\line(1,0){4}}
\put(242,49.567){\makebox(0,0){$b$}}
\put(250,94.000){\line(1,0){4}}
\put(242,94.000){\makebox(0,0){$2$}}
\put(250,162.000){\line(1,0){4}}
\put(242,162.000){\makebox(0,0){$4$}}
\put(427,26){\line(0,1){2}}
\put(427,31){\line(0,1){2}}
\put(427,36){\line(0,1){2}}
\put(427,41){\line(0,1){2}}
\put(427,46){\line(0,1){2}}
\put(427,51){\line(0,1){2}}
\put(427,56){\line(0,1){2}}
\put(427,61){\line(0,1){2}}
\put(427,66){\line(0,1){2}}
\put(427,71){\line(0,1){2}}
\put(427,76){\line(0,1){2}}
\put(427,81){\line(0,1){2}}
\put(427,86){\line(0,1){2}}
\put(427,91){\line(0,1){2}}
\put(427,96){\line(0,1){2}}
\put(427,101){\line(0,1){2}}
\put(427,106){\line(0,1){2}}
\put(427,111){\line(0,1){2}}
\put(427,116){\line(0,1){2}}
\put(427,121){\line(0,1){2}}
\put(427,126){\line(0,1){2}}
\put(427,131){\line(0,1){2}}
\put(427,136){\line(0,1){2}}
\put(427,141){\line(0,1){2}}
\put(427,146){\line(0,1){2}}
\put(427,151){\line(0,1){2}}
\put(427,156){\line(0,1){2}}
\put(427,161){\line(0,1){2}}
\linethickness{0.9pt}
\qbezier(22.000,26.000)(22.875,27.813)(23.750,29.627)
\qbezier(23.750,29.627)(24.625,31.440)(25.500,33.253)
\qbezier(25.500,33.253)(26.375,35.067)(27.250,36.880)
\qbezier(27.250,36.880)(28.125,38.693)(29.000,40.507)
\qbezier(29.000,40.507)(29.875,42.320)(30.750,44.133)
\qbezier(30.750,44.133)(31.625,45.947)(32.500,47.760)
\qbezier(32.500,47.760)(33.375,49.573)(34.250,51.387)
\qbezier(34.250,51.387)(35.125,53.200)(36.000,55.013)
\qbezier(36.000,55.013)(36.875,56.827)(37.750,58.640)
\qbezier(37.750,58.640)(38.625,60.453)(39.500,62.267)
\qbezier(39.500,62.267)(40.375,64.080)(41.250,65.893)
\qbezier(41.250,65.893)(42.125,67.707)(43.000,69.520)
\qbezier(43.000,69.520)(43.875,71.333)(44.750,73.147)
\qbezier(44.750,73.147)(45.625,74.960)(46.500,76.773)
\qbezier(46.500,76.773)(47.375,78.587)(48.250,80.400)
\qbezier(48.250,80.400)(49.125,82.213)(50.000,84.027)
\qbezier(50.000,84.027)(50.875,85.840)(51.750,87.653)
\qbezier(51.750,87.653)(52.038,88.249)(52.325,88.845)
\qbezier(52.325,88.845)(52.913,90.031)(53.500,91.216)
\qbezier(53.500,91.216)(54.375,92.872)(55.250,94.527)
\qbezier(55.250,94.527)(56.125,96.062)(57.000,97.597)
\qbezier(57.000,97.597)(57.875,99.025)(58.750,100.453)
\qbezier(58.750,100.453)(59.625,101.783)(60.500,103.114)
\qbezier(60.500,103.114)(61.375,104.357)(62.250,105.599)
\qbezier(62.250,105.599)(63.125,106.762)(64.000,107.925)
\qbezier(64.000,107.925)(64.875,109.014)(65.750,110.104)
\qbezier(65.750,110.104)(66.625,111.127)(67.500,112.150)
\qbezier(67.500,112.150)(68.375,113.111)(69.250,114.073)
\qbezier(69.250,114.073)(70.125,114.978)(71.000,115.883)
\qbezier(71.000,115.883)(71.875,116.736)(72.750,117.588)
\qbezier(72.750,117.588)(73.625,118.393)(74.500,119.197)
\qbezier(74.500,119.197)(75.375,119.957)(76.250,120.717)
\qbezier(76.250,120.717)(77.125,121.435)(78.000,122.153)
\qbezier(78.000,122.153)(78.875,122.833)(79.750,123.512)
\qbezier(79.750,123.512)(80.625,124.155)(81.500,124.799)
\qbezier(81.500,124.799)(82.375,125.408)(83.250,126.018)
\qbezier(83.250,126.018)(84.125,126.596)(85.000,127.174)
\qbezier(85.000,127.174)(85.875,127.723)(86.750,128.271)
\qbezier(86.750,128.271)(87.625,128.792)(88.500,129.313)
\qbezier(88.500,129.313)(89.375,129.808)(90.250,130.303)
\qbezier(90.250,130.303)(91.125,130.773)(92.000,131.244)
\qbezier(92.000,131.244)(92.875,131.691)(93.750,132.139)
\qbezier(93.750,132.139)(94.625,132.564)(95.500,132.990)
\qbezier(95.500,132.990)(96.375,133.396)(97.250,133.801)
\qbezier(97.250,133.801)(98.125,134.187)(99.000,134.573)
\qbezier(99.000,134.573)(99.875,134.941)(100.750,135.309)
\qbezier(100.750,135.309)(101.625,135.660)(102.500,136.010)
\qbezier(102.500,136.010)(103.375,136.345)(104.250,136.679)
\qbezier(104.250,136.679)(105.125,136.998)(106.000,137.317)
\qbezier(106.000,137.317)(106.875,137.621)(107.750,137.926)
\qbezier(107.750,137.926)(108.625,138.216)(109.500,138.507)
\qbezier(109.500,138.507)(110.375,138.784)(111.250,139.061)
\qbezier(111.250,139.061)(112.125,139.326)(113.000,139.591)
\qbezier(113.000,139.591)(113.875,139.844)(114.750,140.097)
\qbezier(114.750,140.097)(115.625,140.339)(116.500,140.581)
\qbezier(116.500,140.581)(117.375,140.812)(118.250,141.043)
\qbezier(118.250,141.043)(119.125,141.264)(120.000,141.485)
\qbezier(120.000,141.485)(120.875,141.696)(121.750,141.908)
\qbezier(121.750,141.908)(122.625,142.110)(123.500,142.312)
\qbezier(123.500,142.312)(124.375,142.505)(125.250,142.698)
\qbezier(125.250,142.698)(126.125,142.883)(127.000,143.068)
\qbezier(127.000,143.068)(127.875,143.245)(128.750,143.422)
\qbezier(128.750,143.422)(129.625,143.592)(130.500,143.761)
\qbezier(130.500,143.761)(131.375,143.923)(132.250,144.085)
\qbezier(132.250,144.085)(133.125,144.241)(134.000,144.396)
\qbezier(134.000,144.396)(134.875,144.545)(135.750,144.693)
\qbezier(135.750,144.693)(136.625,144.836)(137.500,144.978)
\qbezier(137.500,144.978)(138.375,145.115)(139.250,145.251)
\qbezier(139.250,145.251)(140.125,145.382)(141.000,145.512)
\qbezier(141.000,145.512)(141.875,145.637)(142.750,145.763)
\qbezier(142.750,145.763)(143.625,145.883)(144.500,146.002)
\qbezier(144.500,146.002)(145.375,146.117)(146.250,146.232)
\qbezier(146.250,146.232)(147.125,146.343)(148.000,146.453)
\qbezier(148.000,146.453)(148.875,146.558)(149.750,146.664)
\qbezier(149.750,146.664)(150.625,146.765)(151.500,146.866)
\qbezier(151.500,146.866)(152.375,146.963)(153.250,147.060)
\qbezier(153.250,147.060)(154.125,147.153)(155.000,147.246)
\qbezier(155.000,147.246)(155.875,147.336)(156.750,147.425)
\qbezier(156.750,147.425)(157.625,147.510)(158.500,147.596)
\qbezier(158.500,147.596)(159.375,147.678)(160.250,147.760)
\qbezier(160.250,147.760)(161.125,147.839)(162.000,147.917)
\qbezier(162.000,147.917)(162.875,147.993)(163.750,148.068)
\qbezier(163.750,148.068)(164.625,148.141)(165.500,148.213)
\qbezier(165.500,148.213)(166.375,148.283)(167.250,148.352)
\qbezier(167.250,148.352)(168.125,148.419)(169.000,148.485)
\qbezier(169.000,148.485)(169.875,148.549)(170.750,148.613)
\qbezier(170.750,148.613)(171.625,148.675)(172.500,148.736)
\qbezier(172.500,148.736)(173.375,148.795)(174.250,148.854)
\qbezier(174.250,148.854)(175.125,148.910)(176.000,148.966)
\qbezier(176.000,148.966)(176.875,149.021)(177.750,149.075)
\qbezier(177.750,149.075)(178.625,149.127)(179.500,149.179)
\qbezier(179.500,149.179)(180.375,149.229)(181.250,149.279)
\qbezier(181.250,149.279)(182.125,149.327)(183.000,149.374)
\qbezier(183.000,149.374)(183.875,149.420)(184.750,149.466)
\qbezier(184.750,149.466)(185.625,149.511)(186.500,149.555)
\qbezier(186.500,149.555)(187.375,149.597)(188.250,149.639)
\qbezier(188.250,149.639)(189.125,149.680)(190.000,149.721)
\qbezier(190.000,149.721)(190.875,149.760)(191.750,149.799)
\qbezier(191.750,149.799)(192.625,149.836)(193.500,149.874)
\qbezier(193.500,149.874)(194.375,149.910)(195.250,149.946)
\qbezier(195.250,149.946)(196.125,149.980)(197.000,150.015)
\qbezier(252.000,26.000)(252.863,26.233)(253.727,26.465)
\qbezier(253.727,26.465)(254.590,26.698)(255.453,26.930)
\qbezier(255.453,26.930)(256.317,27.163)(257.180,27.395)
\qbezier(257.180,27.395)(258.043,27.628)(258.907,27.860)
\qbezier(258.907,27.860)(259.770,28.093)(260.633,28.325)
\qbezier(260.633,28.325)(261.497,28.558)(262.360,28.790)
\qbezier(262.360,28.790)(263.223,29.023)(264.087,29.255)
\qbezier(264.087,29.255)(264.950,29.488)(265.813,29.720)
\qbezier(265.813,29.720)(266.677,29.953)(267.540,30.185)
\qbezier(267.540,30.185)(268.403,30.418)(269.267,30.651)
\qbezier(269.267,30.651)(270.130,30.883)(270.993,31.116)
\qbezier(270.993,31.116)(271.857,31.348)(272.720,31.581)
\qbezier(272.720,31.581)(273.583,31.813)(274.447,32.046)
\qbezier(274.447,32.046)(275.310,32.278)(276.173,32.511)
\qbezier(276.173,32.511)(277.037,32.743)(277.900,32.976)
\qbezier(277.900,32.976)(278.763,33.208)(279.627,33.441)
\qbezier(279.627,33.441)(280.490,33.673)(281.353,33.906)
\qbezier(281.353,33.906)(282.217,34.138)(283.080,34.371)
\qbezier(283.080,34.371)(283.943,34.604)(284.807,34.836)
\qbezier(284.807,34.836)(285.670,35.069)(286.533,35.301)
\qbezier(286.533,35.301)(287.397,35.534)(288.260,35.766)
\qbezier(288.260,35.766)(289.123,35.999)(289.987,36.231)
\qbezier(289.987,36.231)(290.850,36.464)(291.713,36.696)
\qbezier(291.713,36.696)(292.577,36.929)(293.440,37.161)
\qbezier(293.440,37.161)(294.303,37.394)(295.167,37.626)
\qbezier(295.167,37.626)(296.030,37.859)(296.893,38.091)
\qbezier(296.893,38.091)(297.757,38.324)(298.620,38.556)
\qbezier(298.620,38.556)(299.483,38.789)(300.347,39.022)
\qbezier(300.347,39.022)(301.210,39.254)(302.073,39.487)
\qbezier(302.073,39.487)(302.937,39.719)(303.800,39.952)
\qbezier(303.800,39.952)(304.663,40.184)(305.527,40.417)
\qbezier(305.527,40.417)(306.390,40.649)(307.253,40.882)
\qbezier(307.253,40.882)(308.117,41.114)(308.980,41.347)
\qbezier(308.980,41.347)(309.843,41.579)(310.707,41.812)
\qbezier(310.707,41.812)(311.570,42.044)(312.433,42.277)
\qbezier(312.433,42.277)(313.297,42.509)(314.160,42.742)
\qbezier(314.160,42.742)(315.023,42.975)(315.887,43.207)
\qbezier(315.887,43.207)(316.750,43.440)(317.613,43.672)
\qbezier(317.613,43.672)(318.477,43.905)(319.340,44.137)
\qbezier(319.340,44.137)(320.203,44.370)(321.067,44.602)
\qbezier(321.067,44.602)(321.930,44.835)(322.793,45.067)
\qbezier(322.793,45.067)(323.657,45.300)(324.520,45.532)
\qbezier(324.520,45.532)(325.383,45.765)(326.247,45.997)
\qbezier(326.247,45.997)(327.110,46.230)(327.973,46.462)
\qbezier(327.973,46.462)(328.837,46.695)(329.700,46.927)
\qbezier(329.700,46.927)(330.563,47.160)(331.427,47.393)
\qbezier(331.427,47.393)(332.290,47.625)(333.153,47.858)
\qbezier(333.153,47.858)(334.017,48.090)(334.880,48.323)
\qbezier(334.880,48.323)(335.743,48.555)(336.607,48.788)
\qbezier(336.607,48.788)(337.470,49.020)(338.333,49.253)
\qbezier(338.333,49.253)(338.917,49.410)(339.500,49.567)
\qbezier(339.500,49.567)(339.780,49.643)(340.060,49.718)
\qbezier(340.060,49.718)(340.923,49.956)(341.787,50.194)
\qbezier(341.787,50.194)(342.650,50.439)(343.513,50.683)
\qbezier(343.513,50.683)(344.377,50.935)(345.240,51.187)
\qbezier(345.240,51.187)(346.103,51.446)(346.967,51.704)
\qbezier(346.967,51.704)(347.830,51.971)(348.693,52.238)
\qbezier(348.693,52.238)(349.557,52.512)(350.420,52.787)
\qbezier(350.420,52.787)(351.283,53.070)(352.147,53.353)
\qbezier(352.147,53.353)(353.010,53.645)(353.873,53.936)
\qbezier(353.873,53.936)(354.737,54.237)(355.600,54.538)
\qbezier(355.600,54.538)(356.463,54.849)(357.327,55.160)
\qbezier(357.327,55.160)(358.190,55.481)(359.053,55.801)
\qbezier(359.053,55.801)(359.917,56.133)(360.780,56.465)
\qbezier(360.780,56.465)(361.643,56.807)(362.507,57.150)
\qbezier(362.507,57.150)(363.370,57.505)(364.233,57.860)
\qbezier(364.233,57.860)(365.097,58.227)(365.960,58.595)
\qbezier(365.960,58.595)(366.823,58.975)(367.687,59.356)
\qbezier(367.687,59.356)(368.550,59.751)(369.413,60.146)
\qbezier(369.413,60.146)(370.277,60.556)(371.140,60.966)
\qbezier(371.140,60.966)(372.003,61.392)(372.867,61.818)
\qbezier(372.867,61.818)(373.730,62.261)(374.593,62.704)
\qbezier(374.593,62.704)(375.457,63.165)(376.320,63.626)
\qbezier(376.320,63.626)(377.183,64.107)(378.047,64.587)
\qbezier(378.047,64.587)(378.910,65.089)(379.773,65.590)
\qbezier(379.773,65.590)(380.637,66.114)(381.500,66.638)
\qbezier(381.500,66.638)(382.363,67.187)(383.227,67.735)
\qbezier(383.227,67.735)(384.090,68.309)(384.953,68.884)
\qbezier(384.953,68.884)(385.817,69.487)(386.680,70.090)
\qbezier(386.680,70.090)(387.543,70.724)(388.407,71.357)
\qbezier(388.407,71.357)(389.270,72.025)(390.133,72.693)
\qbezier(390.133,72.693)(390.997,73.397)(391.860,74.102)
\qbezier(391.860,74.102)(392.723,74.848)(393.587,75.593)
\qbezier(393.587,75.593)(394.450,76.384)(395.313,77.175)
\qbezier(395.313,77.175)(396.177,78.016)(397.040,78.857)
\qbezier(397.040,78.857)(397.903,79.754)(398.767,80.651)
\qbezier(398.767,80.651)(399.630,81.612)(400.493,82.573)
\qbezier(400.493,82.573)(401.357,83.605)(402.220,84.638)
\qbezier(402.220,84.638)(403.083,85.753)(403.946,86.869)
\qbezier(403.946,86.869)(404.810,88.079)(405.673,89.290)
\qbezier(405.673,89.290)(406.536,90.612)(407.400,91.934)
\qbezier(407.400,91.934)(408.263,93.389)(409.126,94.843)
\qbezier(409.126,94.843)(409.990,96.456)(410.853,98.070)
\qbezier(410.853,98.070)(411.716,99.879)(412.580,101.687)
\qbezier(412.580,101.687)(413.443,103.741)(414.306,105.795)
\qbezier(414.306,105.795)(415.170,108.166)(416.033,110.538)
\qbezier(416.033,110.538)(416.896,113.335)(417.760,116.133)
\qbezier(417.760,116.133)(418.623,119.534)(419.486,122.936)
\qbezier(419.486,122.936)(420.350,127.260)(421.213,131.583)
\qbezier(421.213,131.583)(422.076,137.492)(422.940,143.401)
\qbezier(422.940,143.401)(423.803,152.700)(424.666,162.000)
\end{picture}
\caption{The analytic limiting cumulative curve $g$ and its inverse row curve $f$. The solid curves follow the formulas in \eqref{dgp:airy:eq:g}--\eqref{dgp:airy:eq:f}; the junction is $b=\log2$. Dashed lines mark $g(\infty)=2b$ and the divergent endpoint of $f$. Numerical drawing coordinates are illustrative, not sampled data or probability certificates.}
\label{dgp:airy:fig:shape}
\end{figure}

\section{Reproducibility and evidentiary limits}\label{dgp:airy:sec:repro}
The accompanying source archive contains the editable \LaTeX,
a bounded standard-library Python companion, deterministic JSON
outputs, source provenance, tests, and exclusive clean PDF/ZIP
builders. Its exact finite checks test enumeration, the equivalence
of indexed and cumulative constraints, the largest-part injection,
and rational change-of-measure identities. Separate numerical
illustrations evaluate the explicit constants and shape curves.
They do not certify the Airy zero, an asymptotic remainder, a
probability bound, or an onset threshold.

The analytic proof depends on the explicitly stated Mallein estimates
and elementary arguments supplied here. The numerical computation
is independent supporting evidence, never a replacement for those
estimates. The same distinction applies to finite Fourier or
likelihood-ratio probes. No finite list of counts proves a
pointwise asymptotic theorem.

The package fixes a finite allowlist and SHA-256 manifest. It excludes
third-party papers, full OEIS text, working reviews, unrelated reports,
and unlisted files. The archive is reproducible byte for byte under
the recorded Python and \TeX\ toolchain; cross-toolchain identity is
not promised. Build commands disable shell escape and reject existing
output targets, but modified \TeX\ is not thereby made a security
sandbox. The README gives exact bounds, replay commands, and the
trusted-filesystem assumptions.

\subsection{What the refinement establishes}
The pointwise term $-\zeta b C^{-1/6}n^{1/6}$ strengthens the earlier
exact exponential rate; the radial term alone would not have justified
it. The explicit second-order inverse follows from monotonicity and
bracketing. The length and shape conclusions are structural
consequences already available from a coarse coefficient lower bound.
The diagonal row segment has half the limiting mass, while its exact
discrete contacts remain an open microscopic issue in this report.
The published partition class, path encodings, probability estimates,
and general tilting methods are not claimed as inventions here.
\begin{writenote}[6 October 2026]
The package is shipped with the prefix \file{160-airy-}: its companion README at the report root, the companion and its tests, the PDF builder, archive packer and release tools under \file{code/}, and the four result files, the OEIS prefix and the provenance record under \file{data/} (the PDF, README and \file{SHA256SUMS} are not shipped). At intake, on a Windows copy, the four result files were regenerated byte for byte; one of the 23 companion tests fails there (it relies on POSIX open flags). ``Reports~158 and~159'' are Parts~III and~II.
\end{writenote}

\clearpage
\part{Full prefactor and refined inverse for superdiagonal partitions (conditional)}\label{dgp:pre:part}
\partsource{Bundle Report~162, archive \file{Superdiagonal_Partitions_Full_Prefactor_and_Refined_Inverse_Source.zip}, delivered as \file{Report162.tex} (15 pages). Delivered title block ``Full prefactor and refined inverse for superdiagonal partitions / Exact coefficient asymptotics for A238873 / Report 162 / 3 October 2026''. \textbf{This Part's main theorem and corollary are conditional on the unrefereed preprint arXiv:2607.27504v1} \cite{li}; the unconditional superdiagonal statements of the report are those of Parts~III and~IV. Printed as delivered, with \textbf{[write]} notes and a ``conditional'' tag in the two statement headings; the proof of the monotonicity injection is replaced by a pointer to Part~III. Delivered Section~$k$ is Section~$k+38$ here, and its statement or equation $k.j$ is $(k+38).j$. Section~\ref{dgp:sec:further}, which closes the report, is added in the write.}
\begin{center}
\small
\begin{tabular}{@{}>{\raggedright\arraybackslash}p{0.46\textwidth} >{\raggedright\arraybackslash}p{0.46\textwidth}@{}}
\toprule
Reading conventions of this Part & Elsewhere in the report\\
\midrule
$a=a_1<0$ the largest zero of $\operatorname{Ai}$; $s$ the scaled argument (Li's $\zeta$); $L=\log4$ & Part~IV: $\zeta=-a$; Parts~III--IV: $L(x)$ a function\\
$H_q(y)$ Ramanujan's entire function (Li's $A_q(z)$); $\rho_j$ its zeros; $K$, $K_F$ amplitudes; $\alpha$ & Part~IV: $K$ the length; Part~I: $\rho_h$\\
$D=-abC^{-1/6}$; $D_M$; $d$ a cutoff and an inverse coefficient; $p$ a geometric parameter and $p=\frac{11}{12}$; $N$ the size & Part~IV: the same $D$, $D_M(t)$; Parts~I, II: $N(r,n)$\\
\bottomrule
\end{tabular}
\end{center}
\begin{partabstract}
Let $A(n)$ count weakly increasing positive-integer partitions
$\lambda_1\leq\cdots\leq\lambda_k$ with $\lambda_i\geq i$, including the empty
partition at $n=0$. Put $b=\log2$, $C=\pi^2/12+b^2$, and let $a<0$ be
the largest negative zero of $\Ai$. We prove
\[
 A(n)\sim K n^{-11/12}\exp\bigl(B\sqrt n-Dn^{1/6}\bigr),\qquad
 K=\frac{\sqrt2 C^{5/12}}{\sqrt\pi\,\Ai'(a)},
\]
where $B=2\sqrt C$ and $D=-abC^{-1/6}$.
The proof uses the 2025 exact partition generating function, the
2017 positive simple-zero product for Ramanujan's entire function,
and Yu-Tian Li's quantitative turning-point theorem in
\emph{arXiv:2607.27504v1, submitted 29 July 2026}.
The latter is an explicit preprint dependency; no journal publication
of it has been verified.

A positive first-pole estimate is uniform over the relevant moving
length saddle. It gives the complete radial amplitude and its first
relative correction. A separate, self-contained coefficient transfer
uses a real moment-generating-function limit, exponentially accurate
removal of far-tail constraints, and geometric Fourier damping on all
arcs. It does not assume complex-$q$ turning-point asymptotics.
The threshold inverse is obtained through the $\sqrt{x}\log x$ and
$\sqrt{x}$ terms, with $x=(\log y/B)^2$, by monotonicity and integer
bracketing. No pointwise first relative correction, effective finite
crossover, all-orders expansion, or exhaustive priority claim is made.
\end{partabstract}
\noindent\textbf{Scope and attribution.}
This report is self-contained after its explicitly stated source
inputs. The exact combinatorial identity is due to Archibald, Blecher,
Elizalde and Knopfmacher~\cite{abek}; the zero product is due to
\v{S}tampach and \v{S}\v{t}ov\'i\v{c}ek~\cite{ss}; the quantitative
Airy expansion is due to Li~\cite{li}. The coefficient local limit and
their assembly for this partition class are proved below. The two
source convention defects are repaired explicitly, so neither can
hide a normalization factor. Reports 156, 158 and 160 are unchanged
and are not used as proof dependencies. The companion performs bounded
finite checks; no numerical fit is used to determine an asymptotic
constant.
\medskip

\section{Definitions and main conclusions}\label{dgp:pre:sec:main}
Throughout, partitions are written in increasing order. Let
$\mathcal S_n$ be the finite set of lists
\[
 1\leq\lambda_1\leq\cdots\leq\lambda_k,\qquad
 \sum_{i=1}^k\lambda_i=n,\qquad \lambda_i\geq i.
\]
The empty list belongs to $\mathcal S_0$, and $A(n)=|\mathcal S_n|$.
This is the class associated with OEIS A238873~\cite{oeisA238873}.
Set
\begin{align}
 b&=\log2, & L&=2b=\log4, & C&=\frac{\pi^2}{12}+b^2,\label{dgp:pre:eq:constants1}\\
 a&=a_1<0, & B&=2\sqrt C, & D&=-abC^{-1/6},\label{dgp:pre:eq:constants2}\\
 K_F&=\frac{2\sqrt2}{\Ai'(a)},&
 K&=\frac{\sqrt2 C^{5/12}}{\sqrt\pi\,\Ai'(a)},&
 \alpha&=a^2\left(\frac14-\frac b{30}\right).
 \label{dgp:pre:eq:constants3}
\end{align}
Here $a_1>a_2>\cdots$ are the negative Airy zeros, and
$\Ai'(a_1)>0$. All constants in \eqref{dgp:pre:eq:constants2}--\eqref{dgp:pre:eq:constants3}
except $a$ are positive. The radial generating function is
\[
 F(t)=\sum_{n\geq0}A(n)e^{-tn},\qquad t>0.
\]
It is finite because $A(n)$ is bounded by the ordinary partition
number and $\prod_{j\geq1}(1-e^{-tj})^{-1}<\infty$.

\begin{theorem}[Full equivalent and radial correction; \wtag\ conditional on arXiv:2607.27504v1]\label{dgp:pre:thm:main}
Using the source inputs stated in Section~\ref{dgp:pre:sec:inputs}, as $t\downarrow0$,
\begin{equation}
 F(t)=K_F t^{1/3}\exp\left(\frac Ct+ba t^{-1/3}\right)
       \left(1+\alpha t^{1/3}+O(t^{2/3})\right).
 \label{dgp:pre:eq:radial}
\end{equation}
As integers $n\to\infty$,
\begin{equation}
 A(n)\sim K n^{-11/12}\exp\left(B\sqrt n-Dn^{1/6}\right).
 \label{dgp:pre:eq:equivalent}
\end{equation}
Equivalently,
\begin{equation}
 \log A(n)=B\sqrt n-Dn^{1/6}-\frac{11}{12}\log n+\log K+o(1).
 \label{dgp:pre:eq:logequivalent}
\end{equation}
\end{theorem}

\begin{corollary}[Refined threshold inverse; \wtag\ conditional on arXiv:2607.27504v1]\label{dgp:pre:cor:inverse}
For real $y>1$, let $T(y)=\min\{n\geq0:A(n)\geq y\}$ and put
$x=(\log y/B)^2$. Then, as $y\to\infty$,
\begin{equation}
 T(y)=x+\frac{2D}{B}x^{2/3}
       +\frac{11}{6B}\sqrt x\log x
       -\frac{2\log K}{B}\sqrt x+o(\sqrt x).
 \label{dgp:pre:eq:inverse}
\end{equation}
\end{corollary}

For orientation only, the constants have numerical values
\[
\begin{array}{lll}
 C\approx1.30292004734231,& B\approx2.28291046459761,&D\approx1.55073274812252,\\
 K_F\approx4.03363301463591,&K\approx1.27049761356030,&\alpha\approx1.24037790708076.
\end{array}
\]
These decimal evaluations are not certified intervals and play no role
in the proof. In particular \eqref{dgp:pre:eq:radial} does not by itself prove
\eqref{dgp:pre:eq:equivalent}; Section~\ref{dgp:pre:sec:local} supplies the local transfer.
Nor does its displayed correction establish a corresponding pointwise
relative correction. The local-limit error below is only $o(1)$.
\begin{writenote}[6 October 2026]
\emph{Conditional.} Theorem~\ref{dgp:pre:thm:main} and Corollary~\ref{dgp:pre:cor:inverse} use Theorem~1 and Corollary~2 of Li's preprint arXiv:2607.27504v1 \cite{li}, an unrefereed version~1 of 29~July 2026 (on 6~October 2026 arXiv lists no later version and no journal reference); they hold \emph{if} those two statements hold. The heading tags were added by the write. The input enters only through \eqref{dgp:pre:eq:Li}--\eqref{dgp:pre:eq:derivative}: the normalized expansion on compact $s$-sets with an $O(t)$ remainder, and the global indexing of the first two positive zeros. How the other steps depend on it: Lemma~\ref{dgp:pre:lem:pole} uses \eqref{dgp:pre:eq:Li}--\eqref{dgp:pre:eq:derivative} for \eqref{dgp:pre:eq:ratioG}; the proof of Lemma~\ref{dgp:pre:lem:cutoff} divides by \eqref{dgp:pre:eq:radial}, but only needs a lower bound $\log F(t)\ge C/t-O(t^{-1/3})$, which Part~III's \eqref{dgp:rate:eq:gfmain} gives unconditionally; Proposition~\ref{dgp:pre:prop:llt} uses \eqref{dgp:pre:eq:radial} at two nearby arguments to precision $o(1)$ (proof of \eqref{dgp:pre:eq:mgf}), which the unconditional Parts do not supply. The exact identity \eqref{dgp:pre:eq:exactF}, the product \eqref{dgp:pre:eq:product}, the Fourier bounds and the inverse argument of Section~\ref{dgp:pre:sec:inverse} do not use the preprint. Until the preprint is confirmed, this report's unconditional statements about $A(n)$ are Part~III's (the rate) and Part~IV's (the Airy term, the limit shape). The statements agree: \eqref{dgp:pre:eq:logequivalent} implies Part~IV's \eqref{dgp:airy:eq:pointwise}, since $D=-abC^{-1/6}=\zeta bC^{-1/6}$, and \eqref{dgp:pre:eq:radial} implies \eqref{dgp:airy:eq:radial}. Numerically the intake found $K$, $K_F$ and $\alpha$ consistent with exact counts to a few percent (front matter, ``What was checked'').
\end{writenote}

\section{Exact identity and carefully normalized source inputs}\label{dgp:pre:sec:inputs}
\subsection{The exact superdiagonal generating function}
For $0<q<1$, define
\begin{equation}
 H_q(y)=\sum_{j\geq0}\frac{(-1)^j q^{j^2}y^j}{(q;q)_j},
 \qquad (q;q)_j=\prod_{r=1}^j(1-q^r),\qquad
 c_k(q)=[y^k]\frac1{H_q(y)}.
 \label{dgp:pre:eq:H}
\end{equation}
The empty product is one. Theorem 2.1 and equation (6) of
\cite{abek}, with the empty term restored, give
\begin{equation}
 F(t)=\sum_{k\geq0}e^{-t k(k-1)/2}c_k(e^{-t}).
 \label{dgp:pre:eq:exactF}
\end{equation}
Here is the underlying combinatorial normalization. Subtract $i-1$
from the $i$th increasing part and reverse the resulting list. The
result is a positive composition $(\mu_1,\ldots,\mu_k)$ satisfying
$\mu_{i+1}\leq\mu_i+1$. This map is invertible and reduces the weight by
exactly $k(k-1)/2$. The paper's elementary functional equation for these
1-compositions gives the series $1/H_q(y)$ when the empty composition
is included. Thus $c_k(q)\geq0$ and \eqref{dgp:pre:eq:exactF} has exactly the
normalization $A(0)=1$.

The published display in \cite{abek} begins at $k=1$, although its
preceding generating series begins with 1. Starting at $k=0$ repairs
this empty-object convention. It does not multiply $F$ by any factor.
All ensuing summations are over nonnegative combinatorial weights;
finiteness was checked above.
\begin{writenote}[6 October 2026]
Confirmed in the paper: the display of Theorem~2.1 of \cite{abek} is $A(q,t)=\sum_{k\ge1}t^kq^{\binom k2}[y^k]\sum_{i\ge0}\bigl(\sum_{j\ge1}(-1)^{j-1}q^{j^2}y^j/(q;q)_j\bigr)^i$, starting at $k=1$; the inner sum is $1/H_q(y)$. Restoring the $k=0$ term adds the empty partition and gives $A(0)=1$, the convention of every Part of this report. This repairs a convention, not a mathematical error.
\end{writenote}

\subsection{The simple positive-zero product and a source typo}
For each $0<q<1$, the theorem and example in \cite[Theorem 5.2,
pp.~164--167, Example 5.2]{ss} give
\begin{equation}
 H_q(y)=\prod_{j\geq1}\left(1-\frac y{\rho_j}\right),
 \qquad 0<\rho_1<\rho_2<\cdots,\qquad
 \sum_{j\geq1}\rho_j^{-1}<\infty.
 \label{dgp:pre:eq:product}
\end{equation}
The zeros are simple, and the product converges locally uniformly.
For clarity, the relevant hypotheses and the exact normalization can
be checked directly. Take $x_k=q^{(2k-1)/4}$; then
$x_kx_{k+1}=q^k>0$ and $\sum_kx_kx_{k+1}<\infty$.
The paper's associated positive-off-diagonal Jacobi matrix is
self-adjoint and Hilbert--Schmidt. Its nonzero eigenvalues are real
and simple. The paired characteristic product in $w$ is a genus-zero
product in $y=w^2$, with no free exponential factor.

The intermediate identification in the published Example 5.2 prints
an extraneous factor $q$ in
$A_q(w^2)=q\,\mathfrak F(\{wq^{(2k-1)/4}\})$.
At $w=0$ that display would say $1=q$. The correct identity has no
factor $q$, as the defining independent-set sum shows: the coefficient
of $(-w^2)^m$ in $\mathfrak F$ is
\begin{equation}
 \sum_{\substack{1\leq k_1<\cdots<k_m\\k_{i+1}\geq k_i+2}}
       q^{k_1+\cdots+k_m}
 =q^{m^2}\sum_{0\leq r_1\leq\cdots\leq r_m}q^{r_1+\cdots+r_m}
 =\frac{q^{m^2}}{(q;q)_m}.
 \label{dgp:pre:eq:repair}
\end{equation}
We used $k_i=2i-1+r_i$ and the ordinary partition product with at most
$m$ parts. This is precisely \eqref{dgp:pre:eq:H}. The final unit-normalized
product in the source is already correct; \eqref{dgp:pre:eq:repair} repairs
the faulty intermediate display independently before we use it.
\begin{writenote}[6 October 2026]
\emph{Checked in the arXiv version.} The write read arXiv:1510.01454v1 \cite{ss} (the journal version was not available to it): there the factorization theorem is Theorem~15, the positive case is discussed after it, and the $q$-Airy example is Example~19 (p.~16), which defines $A_q(z)=\sum_{n\ge0}q^{n^2}(-z)^n/(q;q)_n$, i.e.\ $H_q$, and prints $A_q(w^2)=q\,\mathfrak F\bigl(\{wq^{(2k-1)/4}\}_{k=1}^\infty\bigr)$, with the extraneous factor $q$ described above, before the unit-normalized product $A_q(z)=\prod_k(1-z/\iota_k(q))$. The definition of $\mathfrak F$ there, $\mathfrak F(x)=1+\sum_{m\ge1}(-1)^m\sum x_{k_1}x_{k_1+1}\cdots x_{k_m}x_{k_m+1}$ over $k_{i+1}\ge k_i+2$, gives \eqref{dgp:pre:eq:repair}, which the write re-derived. Whether the journal's Example~5.2 has the same display was not checked; Part~V's identification of it with the journal numbering is the manuscript's.
\end{writenote}

\subsection{The quantitative turning-point theorem}
The newer input is Yu-Tian Li, \emph{Airy Turning-Point Asymptotics for
Ramanujan's Entire Function $A_q(z)$}, arXiv:2607.27504v1, submitted
29 July 2026, Theorem 1 and Corollary 2~\cite{li}. The exact defining
series is \eqref{dgp:pre:eq:H}; it is not a different $q$-Airy function.
The referenced version is a preprint. Its variable $q=e^{-t}$ has real
$t>0$. The complex uniformity below is in the scaled argument $s$, not
in $t$ or $q$.

Put $E_0=\pi^2/12-b^2$ and
\begin{align}
 y_t(s)&=\frac14\exp(-t/2-t^{2/3}s),\label{dgp:pre:eq:y}\\
 N_t(s)&=\frac{t^{1/6}}{2\sqrt\pi}
  \exp\left(\frac{E_0}{t}-bs t^{-1/3}-\frac{s^2t^{1/3}}4\right),
 \qquad U_t(s)=N_t(s)H_{e^{-t}}(y_t(s)).\label{dgp:pre:eq:N}
\end{align}
For every fixed compact set $\mathcal K\subset\mathbb C$, that theorem
states, uniformly for $s\in\mathcal K$,
\begin{equation}
 U_t(s)=\Ai(s)+\frac{t^{2/3}}{30}
          \bigl(4s\Ai(s)+s^2\Ai'(s)\bigr)+O_{\mathcal K}(t).
 \label{dgp:pre:eq:Li}
\end{equation}
Its corollary identifies the \emph{globally} first positive zero as
$\rho_1=y_t(s_t)$, where
\begin{equation}
 s_t=a-\frac{a^2}{30}t^{2/3}+O(t).
 \label{dgp:pre:eq:zero}
\end{equation}
The second positive zero has coordinate $a_2+O(t^{2/3})$.
Only these two fixed indices are needed. No growing-index
uniformity is inferred.

The remainder in \eqref{dgp:pre:eq:Li} is holomorphic in $s$. Cauchy's integral
formula on a slightly larger fixed disk therefore differentiates it
with the same $O(t)$ order on the smaller disk. Since
$\Ai''(a)=a\Ai(a)=0$,
\begin{equation}
 U_t'(s_t)=\Ai'(a)\left(1+\frac a5 t^{2/3}+O(t)\right).
 \label{dgp:pre:eq:derivative}
\end{equation}
Indeed the derivative of the correction at $a$ is
$[6a\Ai'(a)+a^2\Ai''(a)]/30=a\Ai'(a)/5$; the displacement of
$\Ai'(s_t)$ begins only at order $t^{4/3}$.

\begin{remark}[What the source input supplies]
The normalization \eqref{dgp:pre:eq:N}, the $O_{\mathcal K}(t)$ error, and global
zero indexing are essential. An unnormalized Airy approximation or a
fixed-argument oscillatory formula would not suffice. The preprint
obtains these through a Gaussian integral representation, pole-free
contour deformation, uniform product expansion, cubic descent and
zero counting. Its order-$t^{1/3}$ local density integrates to zero;
its next density gives the correction in \eqref{dgp:pre:eq:Li}. The proof here
uses the stated theorem, not a new proof of that turning-point result.
The older fixed-$x>1/4$ result of Li and Wong~\cite{liwong} is not a
substitute for its moving-turning-point error estimate.
\end{remark}
\begin{writenote}[6 October 2026]
At intake Theorem~1 and Corollary~2 of \cite{li} were read in the arXiv HTML text (5~October 2026): with $q=e^{-\varepsilon}$, $z=\frac{\sqrt q}4e^{-\varepsilon^{2/3}\zeta}$, $\zeta$ in a compact set and $\varepsilon\downarrow0$ real, $N_\varepsilon(\zeta)A_q(z)=\operatorname{Ai}(\zeta)+\frac{\varepsilon^{2/3}}{30}(4\zeta\operatorname{Ai}(\zeta)+\zeta^2\operatorname{Ai}'(\zeta))+O(\varepsilon)$ with $N_\varepsilon(\zeta)=\frac{\varepsilon^{1/6}}{2\sqrt\pi}\exp\{\frac{\pi^2-3L^2}{12\varepsilon}-\frac{L\zeta}{2\varepsilon^{1/3}}-\frac{\varepsilon^{1/3}\zeta^2}4\}$, $L=\log4$; and the $n$th positive zero satisfies $\rho_n(e^{-\varepsilon})=\frac14\exp(-a_n\varepsilon^{2/3}-\frac\varepsilon2+\frac{a_n^2\varepsilon^{4/3}}{30}+O_n(\varepsilon^{5/3}))$, globally indexed. With $\varepsilon=t$, $\zeta=s$ these are \eqref{dgp:pre:eq:y}--\eqref{dgp:pre:eq:zero}: $(\pi^2-3L^2)/12=\pi^2/12-b^2=E_0$, $L\zeta/2=bs$. The hypotheses are met (real $t>0$, compact $s$-sets, the fixed indices $1$ and $2$). Li's proof was not verified, at intake or at the write.
\end{writenote}

\section{Uniform positive first-pole transfer}\label{dgp:pre:sec:pole}
Write $\rho=\rho_1$ and remove the first zero from \eqref{dgp:pre:eq:product}:
\begin{equation}
 G_q(y)=\prod_{j\geq2}(1-y/\rho_j)^{-1}
        =\sum_{j\geq0}g_jy^j,
 \qquad S_q=G_q(\rho)=-\frac1{\rho H_q'(\rho)}>0.
 \label{dgp:pre:eq:G}
\end{equation}
The product convergence proves that $G_q$ is analytic for $|y|<\rho_2$,
that $S_q$ is finite, and that all $g_j\geq0$.
Convolution with $(1-y/\rho)^{-1}$ gives the exact identity and bound
\begin{equation}
 c_k(q)=\rho^{-k}\sum_{j=0}^k g_j\rho^j,
 \qquad 0\leq c_k(q)\leq S_q\rho^{-k}.
 \label{dgp:pre:eq:positive}
\end{equation}
For any $\rho<R<\rho_2$, the positive omitted tail satisfies
\begin{equation}
 0\leq1-\frac{c_k(q)}{S_q\rho^{-k}}
 \leq\left(\frac\rho R\right)^{k+1}\frac{G_q(R)}{S_q}.
 \label{dgp:pre:eq:tailpole}
\end{equation}
To see the last inequality, write each omitted term as
$g_jR^j(\rho/R)^j$ and use $j\geq k+1$.
This step avoids using a fixed-$q$ residue estimate nonuniformly.

Choose one fixed $s\in(a_2,a)$ and set $R=y_t(s)$. Since $y_t$ decreases
with its argument and the first two zero coordinates converge to
$a,a_2$, we have $\rho<R<\rho_2$ for small $t$.
At a zero, differentiation of \eqref{dgp:pre:eq:N} has no contribution from
$N_t'$, and $y_t'(s)=-t^{2/3}y_t(s)$. Thus the following identities are
exact:
\begin{equation}
 G_q(R)=\frac{1-R/\rho}{H_q(R)},\qquad
 S_q=\frac{t^{2/3}N_t(s_t)}{U_t'(s_t)}.
 \label{dgp:pre:eq:exactS}
\end{equation}
Let $\Delta_t=s_t-s\to a-s>0$. Equations
\eqref{dgp:pre:eq:Li}--\eqref{dgp:pre:eq:derivative} show
\begin{align}
 \log\frac{G_q(R)}{S_q}&=b\Delta_t t^{-1/3}+O(1),\label{dgp:pre:eq:ratioG}\\
 \log\frac\rho R&=-\Delta_t t^{2/3}.\label{dgp:pre:eq:rhoratio}
\end{align}
For \eqref{dgp:pre:eq:ratioG}, all factors not displayed in its exponent have
nonzero finite limits: $(1-R/\rho)/t^{2/3}\to s-a<0$,
$\Ai(s)<0$ on $(a_2,a)$, and $U_t'(s_t)\to\Ai'(a)>0$.
The $s^2t^{1/3}$ contribution to the normalization ratio is bounded.
The logarithms are taken of positive real quantities.

\begin{lemma}[Uniform dominance above the barrier scale]\label{dgp:pre:lem:pole}
For every fixed $\delta>0$, there is $c_\delta>0$ such that, uniformly
in integers $k\geq0$ with $kt\geq b+\delta$,
\begin{equation}
 c_k(e^{-t})=S_{e^{-t}}\rho^{-k}
       \left(1+O\bigl(e^{-c_\delta t^{-1/3}}\bigr)\right).
 \label{dgp:pre:eq:uniformpole}
\end{equation}
The upper bound in \eqref{dgp:pre:eq:positive} holds for every $k$.
\end{lemma}
\begin{proof}
The logarithm of the right side of \eqref{dgp:pre:eq:tailpole} is at most
\[
 -\bigl((k+1)t-b\bigr)\Delta_t t^{-1/3}+O(1).
\]
For the stated $k$, the first parenthesis is at least $\delta$ and
$\Delta_t$ is bounded away from zero. Decreasing a positive constant
absorbs the bounded error uniformly. The bound is one-sided before
being written in relative-error notation.
\end{proof}

The threshold $b$ is important. The length saddle will have
$kt\to2b$, so a fixed choice $0<\delta<b$ places its central region
strictly inside the uniform domain. We never need an infinite sum of
residues or asymptotics uniform over an increasing number of zeros.

\section{The Gaussian length sum and the radial amplitude}\label{dgp:pre:sec:radial}
Put
\[
 \lambda=-\log\rho=L+t/2+t^{2/3}s_t,\qquad
 m_t=\frac{\lambda+t/2}{t}=\frac{L+t+t^{2/3}s_t}{t}.
\]
Completing the square gives
\begin{equation}
 T_t:=\sum_{k\geq0}e^{-t k(k-1)/2}\rho^{-k}
 =e^{(\lambda+t/2)^2/(2t)}
       \sum_{k\geq0}e^{-t(k-m_t)^2/2}.
 \label{dgp:pre:eq:T}
\end{equation}
The elementary Gaussian Fourier transform and Poisson summation give
\[
 \sum_{k\in\ZZ}e^{-t(k-m_t)^2/2}
 =\sqrt{\frac{2\pi}t}
   \sum_{\ell\in\ZZ}e^{-2\pi^2\ell^2/t}e^{2\pi i\ell m_t}
 =\sqrt{\frac{2\pi}t}\left(1+O(e^{-2\pi^2/t})\right),
\]
uniformly in the fractional part of $m_t$. Since
$tm_t\to L>0$, the $k<0$ tail is exponentially small relative to the
full sum. Hence, for some $c>0$,
\begin{equation}
 T_t=\sqrt{\frac{2\pi}t}
        e^{(\lambda+t/2)^2/(2t)}(1+O(e^{-c/t})).
 \label{dgp:pre:eq:Tasympt}
\end{equation}
Fix $0<\delta<b$. The region $kt<b+\delta$ is a positive macroscopic
distance from the center $tm_t\to2b$; its Gaussian mass is also
$O(e^{-c/t})$ relative to \eqref{dgp:pre:eq:Tasympt}. On the complement use
Lemma~\ref{dgp:pre:lem:pole}. On the discarded part use the global upper bound
\eqref{dgp:pre:eq:positive}. These two-sided positive comparisons establish
\begin{equation}
 F(t)=S_{e^{-t}}T_t\left(1+O(e^{-c t^{-1/3}})\right).
 \label{dgp:pre:eq:FS}
\end{equation}
In particular there is no uncontrolled small-$k$ contribution.

The exact chain-rule expression \eqref{dgp:pre:eq:exactS} reads
\[
 S_{e^{-t}}=\frac{t^{5/6}}{2\sqrt\pi\,U_t'(s_t)}
       \exp\left(\frac{E_0}t-bs_t t^{-1/3}
                          -\frac{s_t^2t^{1/3}}4\right).
\]
Substitute this in \eqref{dgp:pre:eq:FS} and expand only the Gaussian square.
Since $E_0+L^2/2=C$, the result is
\begin{align}
 F(t)&=\frac{t^{1/3}}{\sqrt2\,U_t'(s_t)}
    \exp\left(\frac Ct+bs_t t^{-1/3}+L
        +\frac{s_t^2t^{1/3}}4+s_t t^{2/3}+\frac t2\right)
 \nonumber\\[-2pt]
 &\hspace{5em}\cdot\left(1+O(e^{-c t^{-1/3}})\right).
 \label{dgp:pre:eq:preexpand}
\end{align}
This identity at the asymptotic precision shown keeps every constant
factor visible. The algebraic powers are $t^{5/6}t^{-1/2}=t^{1/3}$.
The constant exponential factor is $e^L=4$. Thus the leading amplitude
is $4/(\sqrt2\Ai'(a))=K_F$.

For the next radial term, \eqref{dgp:pre:eq:zero} gives
\[
 bs_t t^{-1/3}=ba t^{-1/3}-\frac{ba^2}{30}t^{1/3}+O(t^{2/3}),
 \qquad \frac{s_t^2t^{1/3}}4=\frac{a^2}4t^{1/3}+O(t).
\]
The derivative correction in \eqref{dgp:pre:eq:derivative} and the term
$s_t t^{2/3}$ begin at order $t^{2/3}$. Exponentiating the remaining
order-$t^{1/3}$ term proves \eqref{dgp:pre:eq:radial}, including
$\alpha=a^2(1/4-b/30)$. Dropping the zero displacement would produce an
incorrect radial correction.

\section{A self contained exact coefficient local limit}\label{dgp:pre:sec:local}
This section proves the required transfer from the real radial
formula to an exact coefficient. Its ingredients are elementary
probability, an exact change of measure, and Fourier inversion.
No complex-$q$ version of \eqref{dgp:pre:eq:Li} is assumed.

For a partition, let $Z_j$ be the multiplicity of the part $j$ and
$C_j=\sum_{i\leq j}Z_i$. The indexed condition $\lambda_i\geq i$ is
equivalent to
\begin{equation}
 C_j\leq j\quad\text{for every }j\geq1.
 \label{dgp:pre:eq:barrier}
\end{equation}
Indeed, if $C_j=r>j$, then the $r$th part is at most $j<r$.
Conversely, if $\lambda_i<i$, take $j=\lambda_i$; then $C_j\geq i>j$.
\begin{writenote}[6 October 2026]
This is Part~III's Lemma~\ref{dgp:rate:lem:barrier}; the short proof is kept here, inside the setting of this section.
\end{writenote}
Define heights $H_j=j-C_j$, with $H_0=0$.
Under the Boltzmann law $\PP_t$, the probability of an admissible
partition is $e^{-tN}/F(t)$, where $N=\sum_{j\geq1}jZ_j$.

\subsection{The real moment generating function limit}
For fixed real $u$, $t-ut^{3/2}>0$ when $t$ is sufficiently small, and
\begin{equation}
 \EE_t\exp\left(u t^{3/2}(N-C/t^2)\right)
 =e^{-uCt^{-1/2}}\frac{F(t-ut^{3/2})}{F(t)}
 \longrightarrow e^{Cu^2}.
 \label{dgp:pre:eq:mgf}
\end{equation}
To verify the limit without differentiating a remainder, apply
\eqref{dgp:pre:eq:radial} at the two real arguments. The leading exponent gives
\[
 \frac{C}{t-ut^{3/2}}-\frac Ct-uCt^{-1/2}=Cu^2+O(t^{1/2}).
\]
The difference of the Airy exponents is $O(t^{1/6})$, and that of the
logarithmic powers is $O(t^{1/2})$. The radial relative errors tend to
zero. The mgfs are finite on every fixed real neighborhood of zero
for small $t$, so the mgf continuity theorem yields
\begin{equation}
 t^{3/2}(N-C/t^2)\ \Rightarrow\ \mathcal N(0,2C).
 \label{dgp:pre:eq:clt}
\end{equation}
Consequently the centered characteristic functions converge pointwise
to $e^{-Cu^2}$. Centering at $C/t^2$, rather than asserting an
unproved differentiated mean formula, is sufficient at this scale.

\subsection{Removing the constraints after a fixed cutoff}
Fix $d>b$ and put $M=\floor{b/t}$ and $J=\floor{d/t}$.
Let $F_J(t)$ sum the weights of all multiplicity vectors satisfying
\eqref{dgp:pre:eq:barrier} only for $j\leq J$. The size law normalized by $F_J$
is a convolution of a prefix law and the independent unrestricted tail
\begin{equation}
 V=\sum_{j>J}jZ_j,\qquad
 \PP(Z_j=z)=(1-e^{-tj})e^{-tjz}\quad(z=0,1,\ldots).
 \label{dgp:pre:eq:tail}
\end{equation}
Both its normalizer and its random tail are well-defined: for fixed
$t>0$, the ordinary partition product converges and
$\sum_j j\EE Z_j<\infty$.

\begin{lemma}[Exponentially accurate cutoff]\label{dgp:pre:lem:cutoff}
For fixed $d>b$, there exists $c>0$ such that, for all sufficiently
small $t>0$,
\begin{equation}
 0\leq\frac{F_J(t)}{F(t)}-1\leq e^{-c/t}.
 \label{dgp:pre:eq:cutoff}
\end{equation}
The total variation distance between the two size laws is at most
$e^{-c/t}$, after possibly decreasing $c$.
\end{lemma}
\begin{proof}
Use the independent reference law with $Z_j\sim\operatorname{Geom}(1/2)$
for $j\leq M$, and $Z_j\sim\operatorname{Geom}(e^{-tj})$ for $j>M$;
all geometrics have support $\{0,1,\ldots\}$ and parameter $p$ means
probability $(1-p)p^z$ at $z$. Set
\[
 \mathcal A_M=\sum_{j=0}^{M-1}H_j,\quad
 D_M=4^M e^{-tM(M+1)/2},\quad
 P_M=\prod_{j>M}(1-e^{-tj})^{-1},\quad
 \delta_M=b-tM\geq0.
\]
Summation by parts gives
\begin{equation}
 \sum_{j=1}^M jZ_j=\frac{M(M+1)}2-MH_M+\mathcal A_M.
 \label{dgp:pre:eq:sumbyparts}
\end{equation}
The ratio of the unnormalized Boltzmann weight to the reference
probability is therefore exactly
\begin{equation}
 D_MP_M\exp(-t\mathcal A_M-\delta_MH_M).
 \label{dgp:pre:eq:tilt}
\end{equation}
For example, the prefix ratio before simplification is
$2^{M+C_M}\exp(-t\sum_{j\leq M}jZ_j)$, and \eqref{dgp:pre:eq:sumbyparts} gives
\eqref{dgp:pre:eq:tilt}. On admissibility through $M$, both $\mathcal A_M$ and
$H_M$ are nonnegative, so the final exponential factor is at most one.

A monotone integral comparison for
$f(x)=-\log(1-e^{-x})$ gives
\[
 \log P_M=\frac1t\int_b^\infty f(x)\dd x+O(1)
         =\frac{\Li_2(1/2)}t+O(1).
\]
The lower limit differs from $tM$ by at most $t$ and $f$ is bounded
near $b$. Also $\log D_M=3b^2/(2t)+O(1)$.
Using $\Li_2(1/2)=\pi^2/12-b^2/2$, we obtain
\begin{equation}
 D_MP_M=\exp(C/t+O(1)).
 \label{dgp:pre:eq:normalizer}
\end{equation}

Let $G=\sum_{j=M+1}^J(1-Z_j)$. Under the reference product law,
\begin{align}
 \EE G&=\frac\beta t+O(1),\label{dgp:pre:eq:driftmean}\\
 \beta&=\int_b^d\left(1-\frac1{e^x-1}\right)\dd x
       =d-2b-\log(1-e^{-d})>0.
 \label{dgp:pre:eq:beta}
\end{align}
The integrand is strictly positive for $x>b$. All geometric parameters
in this block are at most $1/2$; their centered exponential moments
therefore have uniformly bounded second logarithmic derivatives on
a fixed neighborhood of zero. There are $O(1/t)$ variables, giving
\[
 \log\EE e^{s(G-\EE G)}\leq K_0s^2/t\qquad(|s|\leq s_0).
\]
Choose a fixed sufficiently small negative $s$. Markov's inequality,
\eqref{dgp:pre:eq:driftmean}, and $\beta>0$ give
\begin{equation}
 \PP\{G<\beta/(2t)\}\leq e^{-c_1/t}
 \label{dgp:pre:eq:chernoff}
\end{equation}
for some $c_1>0$ and all small $t$.

If the prefix through $J$ is admissible but $H_J<\beta/(2t)$, then
$H_M\geq0$ and $H_J=H_M+G$ force the event in \eqref{dgp:pre:eq:chernoff}.
Its tilted mass is at most $D_MP_Me^{-c_1/t}$ by \eqref{dgp:pre:eq:tilt}.
Otherwise $H_J\geq\beta/(2t)$. Choose
$e^{-d}<p_*<1/2$. For small $t$, every parameter after $J$ is at most
$p_*$. Monotone coupling bounds subsequent ruin by that for iid
$Z\sim\operatorname{Geom}(p_*)$. For $r=(1-p_*)/p_*>1$,
\[
 \EE r^{Z-1}=\frac{1-p_*}{r(1-p_*r)}=1.
\]
Thus $r^{\sum_{i=1}^\ell(Z_i-1)}$ is a nonnegative martingale of mean
one. The maximal inequality, first to a finite horizon and then by
monotone convergence, bounds ruin from height $h$ by $r^{-h}$
(and $r^{-\floor h}$ is a harmless weaker bound). From the stated
buffer this is at most $e^{-c_2/t}$ after reducing $c_2>0$.
No optional-stopping equality is required.
Conditional on the prefix, the tail is independent; integrate the
ruin bound and use again that the factor in \eqref{dgp:pre:eq:tilt} is at most
one on the admissible prefix. The lost mass is at most
$D_MP_Me^{-c_2/t}$.

Combining the two cases with \eqref{dgp:pre:eq:normalizer} gives
\[
 0\leq F_J-F\leq \exp(C/t+O(1))e^{-c_3/t}
\]
for a positive $c_3$, with a constant factor absorbed. Divide by
\eqref{dgp:pre:eq:radial}; the loss is only a power of $t$ and
$\exp(O(t^{-1/3}))$. A smaller positive $c$ proves \eqref{dgp:pre:eq:cutoff}.
The true partition law is exactly the cutoff law conditioned on all
remaining constraints, an event of probability $F/F_J$.
Total variation decreases under the size map, and is at most
$1-F/F_J\leq F_J/F-1$, proving the final assertion.
\end{proof}
\begin{writenote}[6 October 2026]
Restated material in this proof: \eqref{dgp:pre:eq:sumbyparts} is Part~IV's \eqref{dgp:airy:eq:prefixweight}, and \eqref{dgp:pre:eq:tilt} is the density in Part~IV's change of measure \eqref{dgp:airy:eq:walkrepresentation} (Part~IV's $\eta_M$ is $\delta_M$ here, and its $D_M(t)=e^{2bM-tM(M+1)/2}$ is this $D_M$); \eqref{dgp:pre:eq:normalizer} is Part~IV's \eqref{dgp:airy:eq:factor} with Part~III's \eqref{dgp:rate:eq:Lb}; the buffer \eqref{dgp:pre:eq:beta} is Part~IV's \eqref{dgp:airy:eq:buffermean}; the martingale ruin bound is Part~III's Lemma~\ref{dgp:rate:lem:survival} and Part~IV's \eqref{dgp:airy:eq:ruin}. They are kept, because the proof is assembled from them in its own order.
\end{writenote}

\subsection{Fourier bounds on the full circle}
Let $\phi_t(\theta)=\EE_t e^{i\theta N}$. The exact convolution in
\eqref{dgp:pre:eq:tail} and Lemma~\ref{dgp:pre:lem:cutoff} imply
\begin{equation}
 |\phi_t(\theta)|\leq|\EE e^{i\theta V}|+2e^{-c/t}.
 \label{dgp:pre:eq:phicutoff}
\end{equation}
The factor 2 is valid for the convention
$d_{\rm TV}(P,Q)=\sup_E|P(E)-Q(E)|$.
For a geometric variable of parameter $p=e^{-tj}$,
\begin{equation}
 \left|\frac{1-p}{1-pe^{ij\theta}}\right|
 =\left(1+\frac{4p}{(1-p)^2}\sin^2(j\theta/2)\right)^{-1/2}.
 \label{dgp:pre:eq:geomchar}
\end{equation}
On the consecutive block $\mathcal B_t=\{J+1,\ldots,J+\floor{1/t}\}$,
the parameters stay in a fixed compact subinterval of $(0,1)$.
Use $\log(1+v)\geq c v$ for bounded $v\geq0$ in \eqref{dgp:pre:eq:geomchar};
all other characteristic factors have modulus at most one. This yields
\begin{equation}
 |\EE e^{i\theta V}|
 \leq\exp\left(-c_4\sum_{j\in\mathcal B_t}\sin^2(j\theta/2)\right).
 \label{dgp:pre:eq:sinbound}
\end{equation}
We now cover every angle in $[-\pi,\pi]$ explicitly.

\emph{Small arc.} Choose $c_0>0$ sufficiently small depending on $d$.
For $|\theta|\leq c_0t$ and all $j\in\mathcal B_t$, the quantities
$|j\theta/2|$ lie in a fixed interval where $|\sin v|\geq c|v|$.
Since $\sum_{j\in\mathcal B_t}j^2\asymp t^{-3}$,
\begin{equation}
 |\EE e^{i\theta V}|\leq e^{-c_5\theta^2t^{-3}}
       \qquad(|\theta|\leq c_0t).
 \label{dgp:pre:eq:smallarc}
\end{equation}

\emph{Angles separated from the small arc.} Put $m=\floor{1/t}$.
The finite geometric-sum formula gives
\begin{equation}
 \sum_{j\in\mathcal B_t}\sin^2(j\theta/2)
 \geq\frac m2-\frac1{2|\sin(\theta/2)|}.
 \label{dgp:pre:eq:farsum}
\end{equation}
For $At\leq|\theta|\leq\pi$, use
$|\sin(\theta/2)|\geq|\theta|/\pi$. For a fixed sufficiently large
$A$ and sufficiently small $t$, \eqref{dgp:pre:eq:farsum} is at least $c/t$.

\emph{The intervening compact range.} For $\theta=ut$ with
$c_0\leq|u|\leq A$, the Riemann sums converge uniformly:
\begin{equation}
 t\sum_{j\in\mathcal B_t}\sin^2(utj/2)
 \longrightarrow \int_d^{d+1}\sin^2(ux/2)\dd x.
 \label{dgp:pre:eq:middlearc}
\end{equation}
The integrand is nonnegative; for $u\ne0$ its integral over a
nondegenerate interval is strictly positive. Continuity and compactness
give a positive minimum on this $u$-range. Combining the three ranges,
for some $c_6>0$,
\begin{equation}
 |\EE e^{i\theta V}|\leq e^{-c_6/t}
       \qquad(c_0t\leq|\theta|\leq\pi).
 \label{dgp:pre:eq:allotherarcs}
\end{equation}
This includes rational and secondary arcs. The consecutive tail
indices also rule out a residual nontrivial lattice span.
\begin{writenote}[6 October 2026]
The same three-arc cover appears in Part~IV, Section~\ref{dgp:airy:sec:fourier}, in a different assembly (note there).
\end{writenote}

\subsection{Fourier inversion and the exact coefficient}
\begin{proposition}[Local limit]\label{dgp:pre:prop:llt}
For any integer targets $m=m(t)$ such that
$(m-C/t^2)t^{3/2}\to z\in\mathbb R$,
\begin{equation}
 \PP_t\{N=m\}\sim\frac{t^{3/2}}{\sqrt{4\pi C}}
                           \exp\left(-\frac{z^2}{4C}\right).
 \label{dgp:pre:eq:llt}
\end{equation}
\end{proposition}
\begin{proof}
Fourier inversion on the integer lattice gives
\[
 \PP_t\{N=m\}=\frac1{2\pi}\int_{-\pi}^{\pi}
                   e^{-im\theta}\phi_t(\theta)\dd\theta.
\]
By \eqref{dgp:pre:eq:phicutoff} and \eqref{dgp:pre:eq:allotherarcs}, the contribution
from $|\theta|\geq c_0t$ is $O(e^{-c_7/t})=o(t^{3/2})$.
On the central arc set $\theta=t^{3/2}u$. The centered characteristic
function converges to $e^{-Cu^2}$ by \eqref{dgp:pre:eq:clt}; multiplying by
the target phase gives pointwise limit $e^{-izu-Cu^2}$.

For a literal integrable majorant, retain the additive error in
\eqref{dgp:pre:eq:phicutoff}. In this central range $|u|\leq c_0t^{-1/2}$,
\[
 e^{-c/t}\leq e^{-(c/c_0^2)u^2}.
\]
Together with \eqref{dgp:pre:eq:smallarc}, this bounds the absolute value of
the rescaled integrand, extended by zero outside the central range,
by $3e^{-c_8u^2}$ for a positive $c_8$ independent of small $t$.
Alternatively the additive error can be integrated separately over
the circle; its total is $o(t^{3/2})$. Dominated convergence is now
valid with no unseparated constant on the whole real line. It gives
\[
 t^{-3/2}\PP_t\{N=m\}\longrightarrow
 \frac1{2\pi}\int_{\mathbb R}e^{-izu-Cu^2}\dd u
 =\frac1{\sqrt{4\pi C}}e^{-z^2/(4C)}.
\]
The limiting density is positive for finite $z$, yielding the stated
equivalent.
\end{proof}

Take $t=\sqrt{C/n}$ and $m=n$, so $m=C/t^2$ exactly and $z=0$.
The defining Boltzmann law gives
\begin{equation}
 A(n)=e^{tn}F(t)\PP_t\{N=n\}.
 \label{dgp:pre:eq:extract}
\end{equation}
Using \eqref{dgp:pre:eq:radial} and \eqref{dgp:pre:eq:llt}, the power of $t$ in
\eqref{dgp:pre:eq:extract} is $1/3+3/2=11/6$. Hence
\[
 A(n)\sim\frac{K_F C^{11/12}}{\sqrt{4\pi C}}
       n^{-11/12}
       \exp\left(2\sqrt{Cn}+baC^{-1/6}n^{1/6}\right).
\]
The coefficient in front is precisely $K$ in
\eqref{dgp:pre:eq:constants3}. This finishes Theorem~\ref{dgp:pre:thm:main}.
Neither this argument nor the value of $K$ uses fitted count data.

\section{The refined threshold inverse}\label{dgp:pre:sec:inverse}
First note that $A(n)$ is nondecreasing.
\begin{writenote}[6 October 2026]
This is Part~III's Lemma~\ref{dgp:rate:lem:mono}. Report~162's proof, the same largest-part injection from $\mathcal S_n$ to $\mathcal S_{n+1}$, is not reprinted.
\end{writenote}
Theorem~\ref{dgp:pre:thm:main} also shows $A(n)\to\infty$, so $T(y)$ is finite.

Let $p=11/12$, $u=\log y$, and $x=(u/B)^2$. For fixed real $c,d,e$,
consider an integer target with
\begin{equation}
 n=x+cx^{2/3}+d\sqrt x\log x+e\sqrt x+O(1).
 \label{dgp:pre:eq:trialinverse}
\end{equation}
Since $(n-x)/x=O(x^{-1/3})$, elementary Taylor expansions in
\eqref{dgp:pre:eq:logequivalent} give
\begin{align}
 \log A(n)-u
 &=\left(\frac{Bc}2-D\right)x^{1/6}
   +\left(\frac{Bd}2-p\right)\log x
   +\frac{Be}2+\log K+o(1).
 \label{dgp:pre:eq:inverseexpand}
\end{align}
For example, the square-root quadratic error is
$O((n-x)^2/x^{3/2})=O(x^{-1/6})$ and tends to zero.
The variation of $n^{1/6}$ is $O(x^{-1/6})$, and
$\log n-\log x=O(x^{-1/3})$.
Thus no unreported constant is concealed in the cross terms.

Choose
\begin{equation}
 c=\frac{2D}B,\qquad d=\frac{2p}B=\frac{11}{6B},\qquad
 e=-\frac{2\log K}B.
 \label{dgp:pre:eq:inversecoeffs}
\end{equation}
For any fixed $\varepsilon>0$, let $n_-(y)$ be the floor of the right
side of \eqref{dgp:pre:eq:trialinverse} without $O(1)$, replacing $e$ by
$e-\varepsilon$; let $n_+(y)$ be its ceiling with $e+\varepsilon$.
Equation~\eqref{dgp:pre:eq:inverseexpand} gives
\[
 \log A(n_-(y))-u\to-B\varepsilon/2<0,
 \qquad
 \log A(n_+(y))-u\to B\varepsilon/2>0.
\]
For all sufficiently large $y$, monotonicity therefore brackets the
threshold as $n_-(y)<T(y)\leq n_+(y)$.
Divide the difference from the central expression by $\sqrt x$,
let $y\to\infty$, and then let $\varepsilon\downarrow0$.
This proves Corollary~\ref{dgp:pre:cor:inverse}. Rounding contributes only $O(1)$;
there is no differentiation of the unknown $o(1)$ remainder.

It is useful to identify the first omitted deterministic cross term,
without claiming a further inverse term. With \eqref{dgp:pre:eq:inversecoeffs},
the square-root quadratic contribution and the first Airy variation
sum to
\begin{equation}
 -\left(\frac{Bc^2}8+\frac{Dc}6\right)x^{-1/6}
 =-\frac{5D^2}{6B}x^{-1/6}=o(1).
 \label{dgp:pre:eq:crossterm}
\end{equation}
The logarithmic-shift cross terms are smaller still. A formal inverse
term of scale $x^{1/3}$ might compensate \eqref{dgp:pre:eq:crossterm}, but the
available $o(1)$ coefficient-log error can induce an error as large as
$o(\sqrt x)$. It does not certify such an $x^{1/3}$ term, an exact
rounding rule, or a numerical finite-input threshold window.
\begin{writenote}[6 October 2026]
Corollary~\ref{dgp:pre:cor:inverse} is conditional, like Theorem~\ref{dgp:pre:thm:main}; its first two terms are Part~IV's unconditional Corollary~\ref{dgp:airy:cor:inverse}, with the same $2D/B=\zeta bC^{-2/3}=1.3585577\ldots$. It is not an instance of the transseries volume's theorems (front matter, ``The inverses''). The $x^{1/3}$ term: Section~\ref{dgp:sec:further}, item~\ref{dgp:q:inverse}.
\end{writenote}

\section{Finite checks and reproducible companion}\label{dgp:pre:sec:companion}
The companion uses bounded standard-library Python and exact integers
or rational arithmetic for its combinatorial and algebraic checks.
Its results are implementation diagnostics, not proofs of the
asymptotic estimates or of a finite eventual-error bound.

\subsection{Two exact counting routes and direct enumeration}
The main recurrence stores $d(k,w)$ after processing the part sizes
$1,\ldots,v$. The entry counts weight-$w$, length-$k$ multiplicity
vectors with every processed cumulative barrier satisfied.
Initially $d(0,0)=1$. At part size $v$, update in increasing weight
order
\[
 d(k,w)\mathrel{+}=d(k-1,w-v)\qquad
   (w\geq v,\ 1\leq k\leq\min(v,K_N)),
\]
where $K_N=\max\{k:k(k+1)/2\leq N\}$.
The increasing order allows an arbitrary multiplicity of $v$.
The restriction $k\leq v$ enforces the new cumulative constraint;
earlier ones are unchanged. Summing the final $d(k,n)$ over $k$ gives
$A(n)$. Exact Python integers avoid a fixed-limb overflow assumption.
The workspace is $O(N\sqrt N)$ cells and the arithmetic operation
count is $O(N^2\sqrt N)$; the public bound is $N\leq400$.

An independent route truncates \eqref{dgp:pre:eq:H} as a formal $q$-series,
forms $c_0=1$ and $c_k=-\sum_{j=1}^k h_jc_{k-j}$, and applies the
weight shift $k(k-1)/2$ in \eqref{dgp:pre:eq:exactF}. All coefficients are
integers. The supplied verification compares this route with the
positive recurrence through $n=80$. Direct recursive enumeration of
ordinary increasing partitions through weight 28 checks the indexed
and cumulative conditions, the resulting counts, and the largest-part
injection independently.

The data file contains the exact recurrence output through 400.
The threshold command searches these actual counts for the first
crossing. Value one has threshold zero. If the requested bound has
not reached a target, the output explicitly says so and does not
substitute an asymptotic estimate.

\subsection{Algebra and numerical illustrations}
Exact rational polynomial arithmetic checks the two Airy-density
integrals mentioned by the source, reducing derivatives with
$\Ai''(s)=s\Ai(s)$. The densities are
\begin{align*}
 Q_1(s,u)&=s(u^2-s)/2,\\
 Q_2(s,u)&=-u^5/30+s^2u^4/8-s^3u^2/4+s^2u/4+s^4/8.
\end{align*}
Under $u^j\mapsto(-1)^j\Ai^{(j)}(s)$, their integrals are respectively
zero and $(4s\Ai(s)+s^2\Ai'(s))/30$.
Finite multiplicity vectors also check \eqref{dgp:pre:eq:sumbyparts} and the
exact tilt \eqref{dgp:pre:eq:tilt} at rational $q$. These tests verify algebra
and conventions, not Li's analytic uniform remainder.

The illustration output uses ordinary floating point and explicitly
recorded approximate Airy constants. It evaluates $B,D,K_F,K,\alpha$
and selected finite normalized log residuals
\[
 \log A(n)-B\sqrt n+Dn^{1/6}+\frac{11}{12}\log n-\log K.
\]
No fitting, empirical extrapolation, or certified rounding claim is
made. Slowly changing finite residuals do not establish a monotone
approach or locate an eventual crossover.

\subsection{Bounds and local release contract}
All public computational integer inputs reject booleans, nonintegers
and out-of-range values. CLI integers use canonical nonnegative ASCII
decimal notation. Counts are capped at 400, the independent formal
series at 100, direct enumeration at 32, and requested exact threshold
values at $10^{200}$. The verification command uses fixed smaller
workloads. Tests run both normally and with Python optimization so
checks cannot disappear with \texttt{-O}. The companion emits JSON to
standard output, has no file-writing API, and performs no network I/O.

The release includes the editable LaTeX source, deterministic PDF
builder, companion, exact and approximate JSON outputs, attribution
record, tests, and a sorted SHA-256 manifest. A fixed allowlist controls
the archive; internal reviews, source-paper PDFs, temporary renders,
build logs and previous reports are excluded. Two clean PDF builds,
two archives and exact JSON replays are compared byte for byte in the
recorded toolchain. PDF pages are rendered and visually inspected.
The README gives the exact commands and the platform qualification.
\begin{writenote}[6 October 2026]
The package is shipped with the prefix \file{162-prefactor-}: its companion README at the report root, the companion, its tests, the PDF builder, archive packer and release tools under \file{code/}, and the four result files and the provenance record under \file{data/} (the PDF, README and \file{SHA256SUMS} are not shipped). At the write, on a Windows copy of the arrival archive, the 14 tests passed and \texttt{counts} regenerated \file{data/counts.json}.
\end{writenote}

Release outputs are exclusive: existing files, links, directories and
special targets are refused. Parent traversal and symlink parents are
rejected with descriptor-relative POSIX operations. Input reads and
archive payloads have explicit size bounds. These tools assume trusted,
nonadversarial source trees and output parents; they are not a sandbox
against malicious same-user path replacement or modified TeX. Shell
escape is disabled, but arbitrary TeX is not made safe by that alone.
There is no upload, publication, external repository change or OEIS edit.

\section{Source status and limits of the result}\label{dgp:pre:sec:limits}
The analytic constants in this report are derived from the exact
normalizations \eqref{dgp:pre:eq:H}, \eqref{dgp:pre:eq:product} and \eqref{dgp:pre:eq:N}.
The preprint theorem is a genuine explicit dependency, even though
its needed normalization, contour orientation, error order and
zero-indexing argument have been independently checked in preparing
this work. Such checking is neither formal proof verification nor
peer review. The report supplies its own full coefficient transfer
rather than treating real radial asymptotics as sufficient.

The 2025 partition paper supplies exact generating functions and
identifies the superdiagonal asymptotics as an open direction in that
source. The 2017 paper supplies the entire-function product; Li's
July 2026 preprint concerns the Ramanujan function itself. A bounded
primary-source search on 3 October 2026 did not find the same full
partition prefactor and inverse in the inspected sources. That is
limited evidence for an application, not exhaustive novelty or priority
certification. Attribution and verified source hashes are retained in
\texttt{SOURCE\_PROVENANCE.json} without redistributing the full papers.

The conclusions stop at \eqref{dgp:pre:eq:equivalent}, the radial correction
\eqref{dgp:pre:eq:radial}, and the threshold precision \eqref{dgp:pre:eq:inverse}.
They do not give a pointwise first relative correction, numerical
constants in an eventual relative-error inequality, a certified finite
crossover, an all-orders expansion, or complex-$q$ turning-point
asymptotics. The central and local probability limits proved here are
for the Boltzmann size variable used in coefficient extraction; no
unmentioned fluctuation theorem for a uniformly sampled fixed-size
partition is inferred.
\begin{writenote}[6 October 2026]
Conditional status: note after Corollary~\ref{dgp:pre:cor:inverse}'s statement. The open items of this section, and an unconditional proof of the amplitude, are collected in Section~\ref{dgp:sec:further}, items~\ref{dgp:q:amplitude}--\ref{dgp:q:complex}.
\end{writenote}

\clearpage
\section{Further questions and research (all Parts)}\label{dgp:sec:further}
\begin{writenote}[6 October 2026]
This section is the write's. Under Vladimir's standing rule of 4~October 2026 it gathers every question left open by the five manuscripts, and every non-claim that marks a gap, with its source, the argument offered and what is missing. The manuscripts' own question and limits sections (Sections~\ref{dgp:conj:sec:questions}, \ref{dgp:sub:sec:sources}, \ref{dgp:rate:sec:sources}, \ref{dgp:airy:sec:repro} and~\ref{dgp:pre:sec:limits}) stay as delivered. No claim of any manuscript was found false, so nothing is refuted here.
\end{writenote}

\subsection*{Answered inside this report}
\begin{itemize}
\item Part~I, Section~\ref{dgp:conj:sec:questions}, ``The superdiagonal exponential constant'': the limit of $\log A(n)/\sqrt n$ exists and equals $B$ (Part~III, Theorem~\ref{dgp:rate:thm:main}), with the matching upper bound $\log A(n)\le B\sqrt n+O(1)$. The ``leading amplitude, correction series, and precise inverse'' it names are answered \emph{in part}: the next term $-Dn^{1/6}$ and the inverse to $o(x^{2/3})$ unconditionally (Part~IV), the amplitude $K$, the power $n^{-11/12}$ and the inverse to $o(\sqrt x)$ conditionally on \cite{li} (Part~V); the correction series stays open (item~\ref{dgp:q:allorders}).
\item Part~I, Section~\ref{dgp:conj:sec:questions}, ``The subdiagonal defect'': answered at $r=0$. By Part~II the defect is $\frac\beta4+(\frac7{16}-\frac3{2\pi^2})\beta^2+O(\beta^3)$ (the write's \eqref{dgp:conj:eq:defect}), so the scale is $n^{-1/2}$, not $1/\log n$, and $s_0(n)$ has an expansion to every order. Re-scoped in item~\ref{dgp:q:defect}.
\item Part~III, Section~\ref{dgp:rate:sec:sources}: ``No optimality claim is made for the error $O(n^{1/3})$''. By Part~IV the true lower error is $B\sqrt n-\log A(n)=Dn^{1/6}(1+o(1))$, so the order $n^{1/3}$ is not optimal, $n^{1/6}$ is, and no bound $\log A(n)\ge B\sqrt n-o(n^{1/6})$ holds (dated note in Section~\ref{dgp:rate:sec:main}).
\end{itemize}

\subsection*{Open}
\begin{enumerate}
\item\label{dgp:q:defect} \emph{The defect at fixed $r\ge1$ and at growing $r$} (Part~I, ``The subdiagonal defect'' and ``Sharper shifts''; re-scoped). Part~I proves $0\le\Prb_n(R\le r)-s_r(n)/p(n)=O(1/\log n)$ uniformly in $r\ge0$; Part~II's endpoint reduction is carried out only for the barrier $\lambda_i\le i$. Missing: the analogue of Part~II's localization and mixed bijection for the shifted barrier $\lambda_i\le i+r$ (first failure patterns at both ends with the shift), and, for growing $r$, a reduction uniform in $r$. Part~II evaluates rank atoms only at fixed ranks; the rank asymptotics it cites are uniform in a growing range (Liu--Zhou's Theorem~1.3 for $m=o(n^{3/4})$, Zhou's Theorem~4.1 uniformly in $m$ with $m^2\beta^3=o(1)$), but no reduction uniform in $r$ has been written. Part~I also asks for an error sensitive to $r/\sqrt n$, smaller where the logistic tail is small.
\item\label{dgp:q:onset} \emph{Explicit onsets and certified thresholds} (all Parts: Part~I, ``Explicit numerical onset''; Part~II, constants $C_J,y_J$ not explicit; Part~III, unspecified $K_0,\dots,K_3$; Part~IV, no rate in the $o$-terms; Part~V, ``no effective crossover''). Every asymptotic statement here has unspecified constants and onset, so no finite search interval or rounding rule is certified. Missing: explicit Rademacher and rank remainders (Parts~I, II), explicit spectral, Chernoff and Mallein-type constants (Parts~III, IV), an effective version of Li's theorem (Part~V). The companions' threshold commands search actual counts instead.
\item\label{dgp:q:series} \emph{The subdiagonal series} (Part~II, Section~\ref{dgp:sub:sec:sources}). Convergence of $\sum c_j\beta^j$, uniformity at growing order, efficient computation, and the coefficients beyond $c_3$ (the algorithm is proved at every order; the companion stops at order three). Also whether the ceilings in Theorem~\ref{dgp:sub:thm:inverse} can be removed: the integer threshold may stay at distance of order one from the smooth centre, so two adjacent candidates can remain.
\item\label{dgp:q:fluct} \emph{Fluctuations and microscopic structure of superdiagonal partitions} (Part~IV, Sections~\ref{dgp:airy:sec:rows} and~\ref{dgp:airy:sec:repro}). Theorem~\ref{dgp:airy:thm:shape} gives laws of large numbers with exponential deviation bounds. Open: a central limit theorem or covariance process for the profile, the diagonal fluctuation scale, the number of exact contacts $\lambda_i=i$ (the theorem is about the limiting row curve only), the centring and law of the largest part, row control at the divergent endpoint $y=2b$, and optimal deviation constants. Part~V's central and local limits concern the Boltzmann size variable, not a uniform partition of fixed size.
\item\label{dgp:q:amplitude} \emph{An unconditional amplitude} (Part~V, and the intake record). Part~V's equivalent \eqref{dgp:pre:eq:equivalent} rests on Li's unrefereed preprint \cite{li}, used for the normalized turning-point expansion \eqref{dgp:pre:eq:Li} on compact $s$-sets with an $O(t)$ remainder and for the global indexing of the first two zeros \eqref{dgp:pre:eq:zero}. Open: a refereed or independent proof of these two inputs, or a different route to $K$; meanwhile the unconditional statements are Part~III's and Part~IV's. Part~IV states that it proves no amplitude and no power prefactor; Part~V supplies both, conditionally.
\item\label{dgp:q:relative} \emph{A pointwise relative correction} (Part~V). Part~V's radial expansion has the relative correction $\alpha t^{1/3}$, but its local limit has only an $o(1)$ error, so no pointwise correction $A(n)=Kn^{-11/12}e^{\cdots}(1+\alpha'n^{-1/6}+\cdots)$ is proved. The intake's pointwise residual ($0.155$--$0.158$ for $n=3200$--$8000$, and least-squares $n^{-1/6}$ coefficients $1.68$ and $1.52$ in the fits described in the front matter) suggests such a term; it proves nothing. Missing: a local limit theorem with a rate, or an Edgeworth correction, uniform in the tilt.
\item\label{dgp:q:allorders} \emph{All orders for $A(n)$} (Parts~III--V). No all-orders expansion, radial or pointwise, is proved; each of the three Parts says so. A radial expansion would need Li's turning-point expansion \eqref{dgp:pre:eq:Li} to higher order, which Part~V does not quote, and a pointwise one would need item~\ref{dgp:q:relative} at every order.
\item\label{dgp:q:inverse} \emph{Finer superdiagonal inverses} (Part~V, remark at \eqref{dgp:pre:eq:crossterm}). The first omitted deterministic cross term $-\frac{5D^2}{6B}x^{-1/6}$ suggests an inverse term of order $x^{1/3}$, but the available $o(1)$ error in $\log A(n)$ allows an inverse error as large as $o(\sqrt x)$, so no such term is certified. Exact rounding is open for all five inverses.
\item\label{dgp:q:complex} \emph{Complex-$q$ turning points} (Part~V, Section~\ref{dgp:pre:sec:limits}). Part~V uses Li's theorem for real $t$ only and avoids complex-$q$ asymptotics by its real local-limit argument. A complex-$q$ version, which would allow a saddle-point proof and possibly the corrections of items~\ref{dgp:q:relative}--\ref{dgp:q:allorders}, is open.
\item\label{dgp:q:literature} \emph{Literature and priority} (every Part). Each manuscript made a bounded search and disclaims exhaustive priority; Part~II calls its expansion an application of Zhou's published theorem; Part~V calls its result ``limited evidence for an application''. A full literature comparison, in particular for the subdiagonal expansion and for the superdiagonal rate, was not made at the write either.
\end{enumerate}
\begin{writenote}[Independent check of the write, 6~October 2026]
An adversarial check made by the intake after the write (commit \file{0ca2804be}) re-read the proofs of Theorems~\ref{dgp:conj:thm:coupling} and~\ref{dgp:conj:thm:super}, Corollary~\ref{dgp:conj:cor:logistic} and Theorem~\ref{dgp:conj:thm:inverse} step by step and found no gap; it recomputed $b$, $B_0$, $B_*=2.0617406$, $B=2.2829105$, $B-\pi/\sqrt3=0.46911$ and $B_*-\pi/\sqrt3=0.24794$. \emph{Counts.} Its own dynamic programs for $A(n)$ and $s(n)$, cross-checked by enumeration for $n\le20$, agree with the OEIS data (61 terms of A238873, 57 of A238875, 49 of A387118 as $p(n)-A(n)$); $A(200)=25548856857$ and $s(200)=1930845918142$. The ratio $q(n)/A(n)$ decreases strictly for $16\le n\le400$ ($0.0021$ at $n=400$, where $\log(A/q)/\sqrt n=0.309$, approaching $B-\pi/\sqrt3$ slowly, as the term $-Dn^{1/6}$ predicts). \emph{Proposition~\ref{dgp:conj:prop:choosable}.} An independent bipartite matching on the intervals $\{1,\dots,y_i\}$, with the parts in decreasing and in shuffled order, agrees with superdiagonality partition by partition for every $n\le24$. The OEIS records were fetched again: A387112 is revision \#19 of 7~September 2026, and Clément Garrot's comment, signed 31~August 2026, gives the same proof and names A238873, A387118 and A388711, so the credit, date and revision are exactly as printed; A238873 (\#36), A238875 (\#33) and A387118 (\#13) are quoted faithfully. \emph{The defect.} The check re-derived $\frac a2-c_2=\frac7{16}-\frac3{2\pi^2}=0.285518224536$, reproduced the six numbers of the note after \eqref{dgp:conj:eq:defect}, and extended them: at $n=2000$ the defect divided by $\beta$ is $0.258505$ and $(\text{defect}-\beta/4)/\beta^2=0.296561$; the cubic residual $(\text{defect}-\frac\beta4-(\frac7{16}-\frac3{2\pi^2})\beta^2)/\beta^3$ is $0.553$, $0.526$, $0.420$, $0.397$, $0.385$ at $n=100$, $400$, $1000$, $1500$, $2000$, drifting towards the next coefficient $\frac12B(0)-c_3=0.3249$ (from \eqref{dgp:sub:eq:ab} and \eqref{dgp:sub:eq:cubic}), which corroborates the $\beta^2$ and $\beta^3$ coefficients jointly. $s(n)/p(n)<\frac12$ for every $10\le n\le2000$ (exact to $400$). \emph{Part~V.} Its conditional status is stated wherever its results appear (title, abstract, Guide, status table, provenance, Part head, both statement headings, the dated notes in Parts~I and~III--V, Section~\ref{dgp:sec:further}, README), and no sentence states its equivalent, $K$ or its inverse unconditionally; the extraneous factor $q$ in Štampach--Šťovíček's Example~19 (arXiv:1510.01454v1) was confirmed. \emph{The inverses.} The instance of the Lambert core for $x_0$ was re-derived, and Part~II's $b_0=\frac7{24}+\frac3{\pi^2}$ and $b_1=\frac{23}{32}+\frac{45}{2\pi^4}$ of \eqref{dgp:sub:eq:b01} were re-derived from the volume's coefficient recursion at orders one and two; at the exact thresholds $y=s(n)$, $(n-x_0-b_0)/\beta_0=1.198$, $1.065$, $1.006$, $0.9915$ at $n=100$, $400$, $1200$, $2000$ (against $b_1=0.949735$), and $n(n-X_2)$ stays bounded ($3.20$ to $2.40$). The non-instance classifications were confirmed. No mathematical error was found. One wording is corrected, with a dated note keeping the first wording (the front matter's ``decays like'' for an upper bound); three dated additions sharpen the note on Dousse--Mertens (the literal $m=0$ reading is false and unused), the note after Proposition~\ref{dgp:conj:prop:choosable} (it also settles the conjectures of A388711 and A387112) and the note on $s(n)/p(n)$ versus $\frac12$ (complete below $n=10$, extended to $n=2000$).
\end{writenote}

\begin{thebibliography}{99}
\raggedright
\bibitem{abek}
M. Archibald, A. Blecher, S. Elizalde and A. Knopfmacher,
\emph{Subdiagonal and superdiagonal partitions},
Afrika Matematika \textbf{36}, article 77 (2025).
\href{https://doi.org/10.1007/s13370-025-01282-0}{doi:10.1007/s13370-025-01282-0}.
\emph{Part~I:} The final article's Section~5 and Lemma~5.4 are the direct targets
and injection antecedent; Theorem~3.3 defines shifted subdiagonal
partitions and gives their generating function.
\emph{Part~II:} Section~5, especially Conjecture~5.3 and Lemma~5.4.
\emph{Part~III:} 17 pp.; published 7 April 2025.
\emph{Part~IV:} The increasing class and refined generating functions are in Section~2;
the asymptotic questions are in Section~5.
\emph{Part~V:} Theorem 2.1 and equation (6), with the empty term retained.
\wtag\ Received 12 November 2024, accepted 27 February 2025; open access under CC~BY~4.0. The conjectures are on p.~15.

\bibitem{dm}
J. Dousse and M. H. Mertens,
\emph{Asymptotic formulae for partition ranks},
Acta Arithmetica \textbf{168} (2015), 83--100.
\href{https://doi.org/10.4064/aa168-1-5}{doi:10.4064/aa168-1-5}.
\emph{Part~I:} Theorem~1.2 in the
\href{https://www.unige.ch/~doussej/rank.pdf}{author-hosted manuscript}
(arXiv:1406.6848), used here only at positive ranks.
\wtag\ Statement of Theorem~1.2 read in \href{https://arxiv.org/abs/1406.6848}{arXiv:1406.6848}, p.~2.

\bibitem{dlmfpart}
NIST Digital Library of Mathematical Functions,
\href{https://dlmf.nist.gov/27.14}{Section 27.14, Unrestricted Partitions}.
\emph{Part~I:} especially Section 27.14(iii), Asymptotic Formulas.
Classical Hardy--Ramanujan and Rademacher estimates; in particular
\href{https://dlmf.nist.gov/27.14.E9}{equation 27.14.9} supports
the refined estimate used in this report.
\emph{Part~II:} \href{https://dlmf.nist.gov/27.14#iii}{Section 27.14(iii),
Unrestricted Partitions: Asymptotic Formulas}.
Classical Hardy--Ramanujan and Rademacher theory.

\bibitem{dlmfairy}
NIST, \emph{Digital Library of Mathematical Functions},
\href{https://dlmf.nist.gov/9.9.T1}{Table 9.9.1}, first ten real zeros
of the Airy functions, checked 3 October 2026 (Part~IV).

\bibitem{oeisA238875}
OEIS Foundation Inc.,
\href{https://oeis.org/A238875}{A238875, Subdiagonal partitions}.
\emph{Part~I:} Numeric prefix and increasing-order convention checked against the
\href{https://github.com/oeis/oeisdata/blob/main/seq/A238/A238875.seq}{official data mirror},
entry revision of 14 April 2025, accessed 3 October 2026.
\emph{Part~II:} Increasing-order definition and numerical prefix checked against the
official data mirror, accessed 3 October 2026.
\wtag\ Revision \#33 of 14 April 2025 read on 6 October 2026.

\bibitem{oeisA238873}
OEIS Foundation Inc.,
\href{https://oeis.org/A238873}{A238873, Superdiagonal partitions}
(Part~III: \emph{A238873: Number of superdiagonal partitions}).
\emph{Part~I:} Numeric prefix and order convention checked against the
\href{https://github.com/oeis/oeisdata/blob/main/seq/A238/A238873.seq}{official data mirror},
entry revision of 7 October 2025, accessed 3 October 2026.
\emph{Part~III:} Definition and numerical prefix checked 3 October 2026.
\emph{Part~IV:} \emph{The On-Line Encyclopedia of Integer Sequences}, A238873, checked 3 October 2026.
\emph{Part~V:} Sequence identification; no full entry text is redistributed.
\wtag\ Revision \#36 of 7 October 2025 read on 6 October 2026.

\bibitem{oeisA387112}
\wtag\ OEIS Foundation Inc.,
\href{https://oeis.org/A387112}{A387112, Numbers with (strictly) choosable initial intervals of prime indices},
revision \#19 of 7 September 2026, read on 6 October 2026; comment of Clément Garrot, 31 August 2026.

\bibitem{lz}
Z.-G. Liu and N. H. Zhou,
\emph{Uniform asymptotic formulas for the Fourier coefficients of
the inverse of theta functions},
The Ramanujan Journal \textbf{57} (2022), 1085--1123.
\href{https://doi.org/10.1007/s11139-021-00409-8}{doi:10.1007/s11139-021-00409-8}.
\href{https://arxiv.org/abs/1903.00835v3}{arXiv:1903.00835v3},
Theorem~1.3, including rank zero (Part~II).
\wtag\ Statement of Theorem~1.3 read in arXiv:1903.00835v3, p.~5.

\bibitem{zhou}
N. H. Zhou, \emph{Eventual log-concavity of $k$-rank statistics
for integer partitions}, Journal of Number Theory
\textbf{259} (2024), 242--272.
\href{https://doi.org/10.1016/j.jnt.2024.01.003}{doi:10.1016/j.jnt.2024.01.003}.
\href{https://arxiv.org/abs/2110.11174v3}{arXiv:2110.11174v3},
Proposition~2.2 and Theorem~4.1, specialized to $k=2$, $m=0$,
and fixed $j=|r|$ (Part~II).
\wtag\ Statement and hypotheses of Theorem~4.1 read in arXiv:2110.11174v3, p.~13.

\bibitem{elizalde}
S. Elizalde, \emph{Research and publications},
\url{https://math.dartmouth.edu/~sergi/research.html}.
Primary publication list checked 3 October 2026; used only for the bounded
source assessment, not as a mathematical input (Part~III).

\bibitem{mallein}
B. Mallein, \emph{Maximal displacement of a branching random walk in
time-inhomogeneous environment}, Stochastic Processes and their
Applications \textbf{125} (2015), 3958--4019.
\href{https://doi.org/10.1016/j.spa.2015.05.011}{doi:10.1016/j.spa.2015.05.011}.
\href{https://arxiv.org/pdf/1307.4496v3}{arXiv:1307.4496v3},
Lemmas~2.6--2.7, equations~(2.15), (2.19).
\href{https://www.math.univ-toulouse.fr/~bmallein/publications/brwtie.pdf}{Author manuscript} (Part~IV).
\wtag\ Lemmas~2.6 and~2.7 and assumptions (2.12)--(2.13) read in arXiv:1307.4496v3, pp.~11 and~13.

\bibitem{ss}
F. \v{S}tampach and P. \v{S}\v{t}ov\'i\v{c}ek,
\emph{Factorization of the characteristic function of a Jacobi matrix},
Operators and Matrices \textbf{11}(1) (2017), 147--169.
\href{https://doi.org/10.7153/oam-11-11}{doi:10.7153/oam-11-11}.
Theorem 5.2, positive-sequence discussion and Example 5.2 (Part~V).
\wtag\ The write read \href{https://arxiv.org/abs/1510.01454v1}{arXiv:1510.01454v1} (6 October 2015), where the theorem is Theorem~15, the positive case follows it, and the example is Example~19 (p.~16); the journal version was not read.

\bibitem{li}
Y.-T. Li,
\emph{Airy Turning-Point Asymptotics for Ramanujan's Entire Function
$A_q(z)$}, arXiv:2607.27504v1, submitted 29 July 2026.
\href{https://arxiv.org/abs/2607.27504v1}{Versioned primary record};
\href{https://arxiv.org/html/2607.27504v1}{versioned full text}.
Theorem 1 and Corollary 2. Preprint; no journal publication verified (Part~V).
\wtag\ Unrefereed. On 6 October 2026 arXiv lists only version~1 and no journal reference. Part~V is conditional on it.

\bibitem{liwong}
Y.-T. Li and R. Wong, \emph{Global asymptotics of Stieltjes--Wigert polynomials}, arXiv:1302.5193v2, 11 June 2013.
\href{https://arxiv.org/abs/1302.5193v2}{Primary preprint record}.
The fixed-argument theorem is not used as the turning-point input (Part~V).

\bibitem{r156}
\emph{Proofs of two conjectures on subdiagonal and superdiagonal
partitions: Quantitative rank coupling and exponential separation},
Report~156 (3 October 2026).
\emph{Part~II:} preceding report in this series.
The subdiagonal result is $s(n)/p(n)=1/2+O(1/\log n)$;
no all-orders result was asserted there.
\emph{Part~III:} Earlier report in this series; it proves the
lower constant $B_*$ and explicitly leaves the exact superdiagonal rate
undetermined. The present proof does not require access to that report.
\wtag\ Printed in this report as Part~\ref{dgp:conj:part}.

\bibitem{r158}
\emph{Exact exponential growth of superdiagonal partitions},
Report~158, 3 October 2026. Antecedent for the exact rate and coarse
pointwise lower bound; unchanged by this report (Part~IV).
\wtag\ Printed in this report as Part~\ref{dgp:rate:part}; its delivered title reads ``Exact exponential growth of \emph{increasing} superdiagonal partitions''.

\bibitem{TSvol} \wtag\ \emph{Transseries and inversion}, ProveIt volume,
\file{Analysis/Transseries/docs/series-and-transseries/Transseries_And_Inversion/transseries_and_inversion.tex}:
Definition \file{p0:def:model}, Lemma \file{p0:lem:alpha-reduction} with \file{p0:eq:alpha-inverse-transport},
Theorems \file{p0:thm:lambert-core} (with \file{p0:eq:core-lambert-minus}),
\file{p0:thm:lambert-centered} (with \file{p0:eq:c-bell}), \file{p0:thm:staircase}
and \file{plt:thm:lw-template}.
\end{thebibliography}
\end{document}
